\chapter{Direct modification in Japanese} \label{direct japanese}

This chapter is dedicated to direct -- that is, attributive -- modification in Japanese. Applying valid tests for direct modification, I present empirical and theoretical evidence that this type of modification does exist in this language and discuss every group of direct modifiers in detail. In the second part of this chapter, I will focus on the lack of ordering restrictions and show how this phenomenon falls out from the theory assumed here, the existence of unspecified functional projections within the direct domain. Next, I will discuss the question of why some languages (English, German) do show strict ordering in the direct domain, some do not (Japanese), and some do only partially, namely Persian and Greek. I will offer an explanation through the \textit{Generalization on Cross-Linguistic Ordering Restrictions} that offers a first approach to predict ordering within the direct domain. Finally, I discuss some remaining questions.

Many candidates for direct-only modifiers are scattered in various works (\citealt{yamakido:2000:japanese, yamakido:2005:nature, yamakido:2007:nature, ayano:2010:revisiting, watanabe:2012:direct, morita:2010:internal, morita:2011:three, morita:2013:morph, nagano:2015:RA, nagano:2016:are}, see especially \citealt{watanabe:2017:attributive, muraki:1996:imi, muraki:2009:nihongo, muraki:2012:nihongo}), that serve as a reference for this chapter. The main insights are based on the results of my corpus study (see Section \ref{corpusanalysis}) and supplemented with interviews with consultants.\footnote{These interviews were conducted online over Zoom in the period between May and June 2021. The procedure was as follows: I showed written examples to five native speakers and immediately asked for their judgment and feedback. Then we discussed critical points. These speakers were not trained linguists.}

\section{Introduction}

As was shown in the preceding chapter, the domain for direct modification is situated inside the DP-internal domain and hierarchically lower than the domain for indirect modification. Furthermore, it was argued that the functional projections inside the direct domain are not based on fine-grained semantic categories. The direct domain is circled in the following structure of the internal DP, repeated from the previous chapter.

\begin{figure}


 \begin{tikzpicture}[baseline=0pt, remember picture]
\tikzset{level distance=30pt, sibling distance=-5pt}
 \Tree [.DP\textsubscript{internal} {} [.D′\textsubscript{internal} [. \node (indirect1) {FP\textsubscript{indirect-n}}; CP [.\node (indirect0) {F′\textsubscript{indirect-n}}; [.FP\textsubscript{Num} NumP [.F′\textsubscript{Num} [.FP\textsubscript{indirect-n+1} \node (indirect3) {CP}; [. \node (indirect2) {F′\textsubscript{indirect-n+1}}; [.dP [. \node (direct2) {FP\textsubscript{direct-n}}; XP [. \node (direct0) {F′\textsubscript{direct-n}}; [.FP\textsubscript{direct-n+1} XP [.F′\textsubscript{direct-n+1} [.$\sqrt{}$P \node (direct1) {XP}; [.$\sqrt{}$′ \edge[roof]; N ] ] F ] ] F ] ] d ] F ] ] F ] ] F ] ] D ] ]
 \node[draw,circle,fit= (direct1) (direct2) ] [color=blue] {};
\end{tikzpicture}

 \caption{Structure of the direct domain in Japanese}
 \label{fig:Structure of the Direct Domain in Japanese2}
\end{figure}



\subsection{Goals and overview}

This chapter will present the following results.

\ex. \a. Valid cross-linguistic as well as language-specific criteria for direct modification can be put forth in Japanese. The main criterion is \textit{lack of predicativity}, but this is often accompanied by other criteria such as restrictions in gradability.
\b. Japanese exhibits several groups of direct modifiers: Non-intersective modifiers, relational modifiers, direct qualitative modifiers, idiomatic modifiers, quantifiers and the language-specific group of \textit{rentaishi}.
\c. The fact that the functional projections inside the direct domain are not specified for certain semantic categories predicts the lack of ordering restrictions in linear structure.
\d. There is parametric variation as to whether a language exhibits fine-grained dedicated functional projections within the direct domain.

\subsection{Arguments against direct modification}

As mentioned in the introduction, many researchers have denied the existence of direct modification in Japanese \citep{kuno:1973:structure, hinds:1988:japanese, whitman:1981:internal, kaplan:1995:the, sproat:1991:the, baker:2003:verbal, laenzlinger:2011:elements}. And while single instances of direct modifiers have been mentioned in the literature (see \citealt{watanabe:2017:attributive}) and Japanese grammarians have long known about the category of solely attributively used lexemes, sometimes called \textit{rentaishi} (see Section \ref{rentaishi}), especially in syntactic research, Japanese is characterized as a language without direct modification. 

The main arguments are the missing mechanism for disambiguation between direct (attributive) and indirect (predicative) modification and the \is{absence of ordering restrictions}lack of ordering restrictions. Some more arguments are given in \citet{baker:2003:verbal}. Building on \citet{nishiyama:1999:adj}, who argues that Japanese adjectives always project a CP layer (see Section \ref{copulaelements}), \citet{baker:2003:verbal} explicitly claims that direct modification is lacking in Japanese and all adjectives ``form relative-clause-like structures" \citep[19]{baker:2003:verbal}. In his framework, direct modification is an operation limited to the adjective class and denotes the modification of a noun without an additional element \citep[191]{baker:2003:cat}. A visible output of direct modification is phi-features, which Japanese lacks in general. Furthermore, he puts forth a systematic lack of \is{intersectivity}non-intersectivity in Japanese and illustrates this with the equivalent of the famous \textit{beautiful dancer} example \citep{larson:1998:events}, which does not show the same ambiguity as in English between a non-intersective (subsective) and an intersective interpretation. Instead, only the intersective interpretation is available as expected under a relative clause analysis in which the intersective reading is predicatively derived (cf. \citealt{cinque:2010:adj} and Section \Ref{intersective}). This is shown in \ref{utsukushiibaker} based on \citet[3]{baker:2003:verbal}. \is{intersectivity}

\exg. utsukushi-i kashu\\
beautiful-\textsc{i} singer \\
1. $\neq$`beautiful as a singer’ [beautiful singer] $\subseteq$ [singer] (subsective)\\
2. $\approx$ `beautiful and a singer’ [beautiful singer] = [beautiful] $\cap$ [singer] (intersective) \label{utsukushiibaker}

As will be shown below, not only are ordering restrictions not a prerequisite for direct modification in Japanese, there are also good counterarguments to the alleged absence of non-intersectivity, showcasing that none of the arguments given in the literature are convincing enough to deny the absence of direct modification in this language.

\section{Criteria and tests for direct modification} \label{testsdirect}

In order to prove that direct-only modification does exist in Japanese, it is of great importance to have specific criteria as substantial guidelines to apply to potential groups of direct modifiers. When establishing such criteria, one has to take into account not only cross-linguistic, but also language-specific aspects. Typical cross-linguistic testbeds, presented throughout the preceding chapters, are \textit{predication, gradability} and \textit{tense/inflection}. These will be looked at here, sidelining the factor \textit{ordering}, and supplemented with two language-specific factors connected to different occurrences of the copula: \textit{coordination} and \textit{co-occurrence with the canonical copula}.

\subsection{Predication} \label{test predication}

Recall from Section \ref{directmodifiers} that the stand-out characteristic of direct modifiers, in Cartography and elsewhere, is \textit{lack of predicative use}. In English, non-verbal modifiers are licensed in predicative position via a copula, but this process is illicit with direct adjectives such as the \is{non-intersective modifier}non-intersective \textit{former} (\citealt{siegel:1976:capturing, larson:1998:events, cinque:2010:adj, cinque:2014:semantics} etc.).

Recall that Japanese \textit{na}-adjectives and \textit{no}-modifiers appear with the copular element \textsc{da} or the longer canonical copula \textsc{de-ar-u} in predicative position, the latter historically derived from the former \citep{urushibara:1993:syntactic, nishiyama:1999:adj, nakahara:2002:japanese, frellesvig:2010:history, miyama:2011:da}, as exemplified below for the \textit{na}-adjective \textit{nigiyaka-na} `lively’.

\ex. \ag. Ano machi-wa nigiyaka-da.\\
that city-\textsc{top} lively-\textsc{cop}\\
`That city is lively.’
\bg. Ano machi-wa nigiyaka-de-ar-u.\\
that city-\textsc{top} lively-\textsc{pred-cop-prs}\\
`That city is lively.’

On the other hand, \textit{i}-adjectives appear with the element \textsc{i} in predicative position, which is homophonous to the output of \is{MOD@\textsc{mod}}\textsc{mod} following \textit{i}-adjectives in attributive position, as argued in the preceding chapter. I will gloss \textsc{i} in predicative position as \textsc{cop}, matching the glossing for \textsc{da}.\footnote{For a precise analysis of all these elements, see Section \ref{copulaelements}. See especially Section \ref{predicativei} on why \textsc{i} in predicative position is a copula as well.}

\exg. Ano machi-wa ōki-i.\\
that city-\textsc{top} big-\textsc{cop}\\
`That city is big.’

Assuming indirect modification is the only kind available in Japanese, there should be no lexemes that appear as attributive modifiers, but are unable to appear in predicative position. However, several Japanese modifiers resist predicative use, for example the \textit{na}-adjective \textit{kanzen-na} `complete' in \ref{kanzen} \citep{yamakido:2005:nature} or the \textit{no}-modifier \textit{Nihon-no} `Japan(ese)’ in \ref{nihon direct}.

\ex. \label{kanzen} \ag. Max-ga kanzen-na baka-da.\\
Max-\textsc{nom} complete-\textsc{mod} idiot-\textsc{cop}\\
`Max is a complete idiot.' \hfill \citep[68]{yamakido:2005:nature}
\bg. *Ano baka-ga kanzen-da.\\
that idiot-\textsc{nom} complete-\textsc{cop}\\
(`That idiot is complete.’)

\ex. \label{nihon direct} \ag. Nihon-no kuruma\\
Japan-\textsc{mod} car\\
`Japanese car'
\bg. *Kono kuruma-wa Nihon-da.\\
this car-\textsc{top} Japan-\textsc{cop}\\
(`This car is Japanese.’)

Importantly, one needs to take into account other exponents of the copula such as \textsc{desu} (and the forms analyzed in Section \ref{copula}), which is a more \is{polite copula}polite version of \textsc{da}. Accordingly, the two instances of the copula behave alike. In instances where \textsc{da} is impossible, \textsc{desu} is as well.\footnote{This element, which can never appear in attributive position, will be discussed in Section \ref{desu} where it will be depicted as a combination of the \is{predicative copula}predicative copula, the \is{polite copula}polite form of the dummy copula and a tense marker for non-past. For now, I will apply the gloss \textsc{cop.pol} to highlight its function as a polite copula.}

\ex. \ag. *Ano baka-ga kanzen-desu.\\
that idiot-\textsc{nom} complete-\textsc{cop.pol}\\
(`That idiot is complete.’; polite)
\bg. *Kono kuruma-wa Nihon-desu.\\
this car-\textsc{top} Japan-\textsc{cop.pol}\\
(`This car is Japanese.’; polite)

Such examples have provided the main motivation for authors such as \citet{yamakido:2005:nature} to argue against a universal analysis of Japanese modifiers as indirect. \textit{Predication}, especially its absence, will be used as the main criterion to check for direct modification verified via the occurrence of the copula forms \textsc{da} and \textsc{i}.

It should be noted here that in Japanese, seemingly direct modifiers can, in specific instances, appear with the copula. First, predicative use can be facilitated easily in contrastive statements. An example is the following (and see \citealt{hoshi:2002:kaynean, shimada:2018:relational}).

\exg. Kono kuruma-ja-na-ku-te, ano kuruma-ga Nihon-da.\\
this car-\textsc{pred.top-neg-pred-ger} that car-\textsc{nom} Japan-\textsc{cop}\\
`Not this car, that car is Japanese.’ \label{kurumajanakute}

Here, the direct-only modifier \textit{Nihon} `Japan(ese)’ seemingly appears in predicative position. See also the following example reported in \citet{hoshi:2002:kaynean}, who argues against a direct status of \textit{omoigakena-i} `unexpected’, which \citet{yamakido:2005:nature} argues to be direct.

\exg. Kono kyaku ja-na-ku-te, ano kyaku-ga omoigakena-i no-da.\\
this guest \textsc{pred.top-neg-rky-ger} that guest-\textsc{nom} unexpected-\textsc{cop} \textsc{emp-cop}\\
`Not this visitor, that visitor is unexpected.' \\ \null \hfill \citep[11]{hoshi:2002:kaynean}

In the literature on \is{relational modifier}relational adjectives, many such instances in which otherwise non-predicative relational adjectives appear in predicative position are reported, including class shift and the availability of a \is{kind-reading}\textit{kind}-reading. This instances are referred to as \textit{adjective stranding} (\citealt{levi:1978:syntax, mcnally:2004:relational, nagano:2018:conversion}; see Section \ref{RA}). Another typical way to achieve predicative occurrence involves contrastive use, as in the Japanese examples, shown for \is{English} English below. \is{adjective stranding}

\ex. Her infection turned out to be bacterial, not viral. \null \hfill \citep[254]{levi:1978:syntax} \label{viral stranding}

In \ref{viral stranding}, the relational adjective \textit{viral} is used predicatively. In a context-neutral sentence, such a use is significantly degraded.

\ex. ??Her infection is viral.

According to \citet[255]{levi:1978:syntax}, in such instances the second noun is elided under identity at a late stage in the derivation. One piece of evidence she puts forth is that the two instances, the noun and the modifier, must be established as a set. The first instance of the noun then licenses the deletion of the second instance under phonological identity. This means \ref{viral stranding} is derived from \ref{stranded analysis}.

\ex. Her infection turned out to be (a) bacterial \st{infection}, not (a) viral \st{infection}. \\ \null \hfill \citep[254]{levi:1978:syntax}\label{stranded analysis}

Applying the data to Japanese, this means that the underlying construction of \ref{kurumajanakute} is \ref{kuruma analysis}. Since no overt modifier is present, the syntactic head \is{MOD@\textsc{mod}}\textsc{mod} is deleted alongside.

\exg. Kono kuruma-ja-na-ku-te, ano kuruma-wa Nihon-\st{no} \st{kuruma} da.\\
this car-\textsc{pred.top-neg-pred-ger} that car-\textsc{top} Japan-\textsc{mod} car \textsc{cop}\\
`Not this car, that car is \st{a} Japanese \st{car}.’ \label{kuruma analysis}

Evidence for this claim comes from the fact that this sentence is even more natural when the noun \textit{Nihon} `Japan(ese)’ is followed by a so-called \is{bound marker}\textit{bound marker}, a nominal suffix which specifies the denotation of the noun \citep{nagano:2016:are, shimada:2018:relational}. This in turn makes it easier for the two nouns to be established as a set and facilitates predicate use (see Section \ref{Ra japanese predicative} and appendix A for more on bound markers.)

\exg. Kono kuruma-wa Nihon-sei-d-a.\\
this car-\textsc{top} Japan-produce-\textsc{pred-cop.prs}\\
`This car is Japanese.’ (Lit. `is Japan-produce’)

Furthermore, the copula can oftentimes function as a placeholder for other verbs. This is the famous \textit{eel-construction} or \textit{unagi-bun} coined by \citet{okutsu:1978:boku}. In \ref{eel-sentence}, the copula \textsc{da} appears, but is understood as a lexical verb, for example \textit{tabe} `to eat’, depending on the context. Semantically, it does not express an equative relation between \textit{I} and \textit{eel}. \is{eel-construction}

\exg. Kare-wa sakana-wo tabe-ta. Boku-wa unagi-da.\\
he-\textsc{top} fish-\textsc{acc} eat-\textsc{acc} I-\textsc{top} eel-\textsc{cop}\\
`He ate a fish. I ate/had (the) eel.’ (Lit.: `I am eel.’) \label{eel-sentence}

This means that, provided the appropriate context, formerly non-predicative lexemes can become acceptable in predicative position. \textit{Eel}-constructions are very productive and usually involve a noun directly preceding the copula \citep[23–29]{okutsu:1978:boku}. Since many of the \is{relational modifier}relational modifiers discussed here are nominal in nature, it is no surprise that they can appear in such constructions. In other words, it is oftentimes not clear if we are dealing with a proper predicative phrase in the relevant instances. In the discussion at the end of this chapter, I will further discuss some indirect features of direct modifiers in Japanese.

\subsection{Gradability} \label{gradability}

A further typical characteristic often observed with direct modifiers and generally related to their less complex nature is \is{gradability}\textit{non-gradability}, the inability to express a degree. Gradability can be expressed through modification by degree heads \ref{toolarge}, degree adverbs \ref{verylarge} or the formation of the comparative or superlative \ref{thanlarge} (cf. \citealt{baker:2003:cat}).

\ex. \a. a too large building \label{toolarge}
\b. a very large building \label{verylarge}
\c. a larger building than Tokyo Tower \label{thanlarge}

It is a general property of the word class \textit{adjective} to express a degree, although not all adjectives are equally freely gradable \citep[70–74]{dixon:2010:grammatical}. A restriction in gradability is a central characteristic of both \is{relational modifier}relational \ref{relational, non-gradable} and \is{non-intersective modifier}non-intersective adjectives \ref{nonintersective, non-gradable} in most European languages (see references in the preceding chapter).

\ex. \label{relational, non-gradable} \a. *very / *too /*more nuclear energy
\b. *very / *too / *more American attack

\ex. \label{nonintersective, non-gradable} \a. *very / *more former minister
\b. *totally present editor

The semantic approach to \is{gradability}gradability makes use of scale structures \citep{kennedy:2005:scale, kennedy:2007:vagueness, morzycki:2012:adj}. Gradable adjectives are understood as \textit{relative} in contrast to \textit{absolute} or \textit{totally closed scale} adjectives. Only the former can be modified by so-called \textit{proportional modifiers}, including the common degree adverb \textit{very}, whereas the latter allow modification via \textit{endpoint-oriented modifiers} \citep{kennedy:2005:scale}, also called \is{modifiers of extension}\textit{modifiers of extension} \citep{fabregas:2014:adjectival} such as \textit{completely}, denoting a maximum scale, or \textit{slightly}, denoting a minimum scale \citep[355]{kennedy:2005:scale}. The term \textit{non-gradable modifier} will be used in this book to denote those modifiers that cannot be combined with any degree adverb and do not have access to comparative-formation.

It is important to note that lack of gradability does not automatically entail lack of predication and that gradability is not a sufficient criterion for direct modification, yet often an accompanying one. Absolute adjectives such as \textit{pregnant} or \textit{bent}, despite restrictions in gradability, can productively appear in predicative position in English. \citet{cinque:2010:adj} also provides examples of some \is{non-intersective modifier}non-intersective adjectives in \is{English} English (and \is{Italian} Italian and \is{Bulgarian} Bulgarian) that can be combined with degree adverbs.

\ex. the most probable winner \hfill \citep[47]{cinque:2010:adj}

Gradability in Japanese is expressed via the particle \textit{yori(-mo)} `than' and the optional adverb \textit{motto} `more' or by various degree adverbs (see Section \ref{intro japanese}). The most prototypical one is \textit{totemo} `very’, the equivalent of the proportional modifier \textit{very} in English. Equivalents of endpoint(-oriented) modifiers are the maximum-standard degree adverbs \textit{kanzenni} and \textit{mattaku}, both `completely’, and the minimum-standard degree adverb \textit{wazukani} `slightly’ (cf. \citealt{oshima:2019:gradability}).

\largerpage
It is often claimed that \textit{no}-modifiers are non-gradable in general \citep{kato:2003:nihongo, morita:2011:three, morita:2013:morph}. And while most are incompatible with the degree adverb for relative comparison \textit{totemo} `very’ \Ref{totemo mumei}, many are compatible with the adverb \textit{kanari} `quite’, as well as the maximum- and minimum-standard degree adverbs mentioned above \ref{kanari mumei}.\footnote{A reviewer points out that example \ref{wazukani mumei-no} is pragmatically odd, as, indeed, the standard way to express this notion would be a combination along the lines of `very famous'. Nevertheless, this example serves an illustrative illustration purpose.}

\exg. ??totemo mumei-no haiyū\\
very unknown-\textsc{mod} actor\\
`very unknown actor’ \label{totemo mumei}

\ex. \label{kanari mumei} \ag. kanari mumei-no haiyū\\
quite unknown-\textsc{mod} actor\\
`quite unknown actor'
\bg. kanzenni / mattaku mumei-no haiyū\\
completely / completely unknown-\textsc{mod} actor\\
`completely unknown actor'
\cg. wazukani mumei-no haiyū\\
slightly unknown-\textsc{mod} actor\\
`slightly unknown actor' \label{wazukani mumei-no}

Also note that many \textit{no}-modifiers, for example \textit{mumei-no} `unknown’, can productively appear inside a comparative construction with the comparative particle \textit{yori} `than’.

\exg. Tanaka-san-yori (motto) mumei-no haiyū-to at-ta.\\
Tanaka-Mr.-than more unknown-\textsc{mod} actor-\textsc{com} meet-\textsc{pst}\\
`I met with a more unknown actor than Mr. Tanaka.’

Although this shows that \is{gradability}gradability is restricted for \textit{no}-modifiers, as also corroborated in the study of \citet{oshima:2019:gradability} in which the authors judged only 22.25\% of the 200 most used \textit{no}-modifiers on the \is{BCCWJ}BCCWJ to be gradable,\footnote{They followed up on their judgment with a questionnaire survey and identified only 8.28\% of the 169 \textit{no}-modifiers they considered as gradable on a relative scale, thus with \textit{totemo} \citep[29]{oshima:2019:gradability}.} there are some \textit{no}-modifiers which are compatible with relative degree adverbs of the type \textit{totemo} `very’, although only a few. See \ref{gradable no} and \citet{oshima:2019:gradability}, as well as their data set, for more examples.

\ex. \label{gradable no} \ag. totemo totsuzen-no henka\\
very sudden-\textsc{mod} change\\
`a very sudden change’
\bg. yori ōgata-no fune\\
more large.size-\textsc{mod} ship\\
`a ship of a larger size’ \null \hfill \citep[16]{oshima:2019:gradability}

\newpage
Nevertheless, most non-predicative modifiers in Japanese resist gradability across the board, as exemplified for the \is{non-intersective modifier}non-intersective modifier \textit{shōrai-no} `future’ in \ref{shorai non-gradable} and to be discussed in the following.

\exg. *totemo shōrai-no daitōryō\\
very future-\textsc{mod} prime.minister\\
(`a very future prime minister’) \label{shorai non-gradable}

\subsection{Tense and inflection} \label{test tense}

The absence of tense information is another criterion that differentiates direct from indirect modifiers. This can be easily illustrated with \is{German} German where adjectives can be used attributively together with a copula appearing as a prenominal participle, a case of a \is{reduced relative clause}reduced relative clause (\citealt[203]{cinque:2020:relcl}; \citealt{krause:2001:rrc, douglas:2016:syntactic, harwood:2018:rrc}). This is shown in \ref{ehemalig weis} for the adjective \textit{weiß} `white’, which can be used as a predicate. The \is{non-intersective modifier}non-intersective adjective \textit{ehemalig} `former’, on the other hand, can never appear together with an inflected copula, as it can never appear with the copula in the first place \ref{ehemalig fruher}, highlighting its direct status.

\exg. das (früher einmal) weiß gewesene Hemd (German)\\
the earlier once white be.\textsc{pst.ptcp} shirt\\
`the shirt which was (once) white’ \label{ehemalig weis}

\exg. *der ehemalig gewesene Präsident (German)\\
the former be.\textsc{pst.ptcp} president\\
(`the was former president’) \label{ehemalig fruher}

As shown in Section \ref{internalstructure}, Japanese modifiers can inflect for past DP-internally, but this is impossible for non-predicative modifiers such as \textit{kanzen-na} `complete’ \Ref{kanzendatta}.\footnote{That non-intersective adjectives cannot inflect for past also applies in Korean as shown by \citet[164]{kang:2005:adj}.} \is{Korean}

\exg. *ano kanzen-d-at-ta baka\\
that complete-\textsc{pred-cop-pst} idiot\\
(`that idiot who was complete’) \label{kanzendatta} \null \hfill \citep[109]{ayano:2010:revisiting}

From a theoretical angle, these facts fall out naturally. The structure of direct modifiers is less complex and no IP layer is present, hence no overt tense can be expressed. In a nutshell, whereas the endings \textsc{de-ar-u}, \textsc{d(e)-at-ta} and \textsc{k-at-ta} always carry tense information, this is not obligatorily the case for \textsc{i}, \textsc{na} and \textsc{no}. As shown in Section \ref{indirect mod}, this only applies when they are indeed part of a relative clause (see also \citealt[793]{watanabe:2017:attributive}). Direct modifiers can therefore never co-occur with the complex copula forms instead of the monomoraic endings. 

\subsection{Co-occurrence with other copula forms} \label{copula}

I will now discuss two rather language-specific tests. These correlate with the inability of predicative use as they look at co-occurrence with two different exponents of the copula, namely the canonical copula \textsc{de-ar-u} and coordination forms.

\subsubsection{The canonical copula DE-AR-U}

The complex form \textsc{de-ar-u} is usually called the \textit{standard} or \textit{canonical} copula \citep{urushibara:1993:syntactic, nishiyama:1999:adj, nakahara:2002:japanese}. Recall (for instance from Section \ref{intro japanese}) that this canonical copula co-occurs with \textit{na}-adjectives and \textit{no}-modifiers and that it is not only mutually exclusive with the copula \textsc{da}, but also with the attributive exponents of \is{MOD@\textsc{mod}}\textsc{mod}, namely \textsc{na} and \textsc{no}. This is repeated for convenience for the \textit{na}-adjective \textit{nigiyaka-na} `lively’ \ref{nigiyaka convenience} and the \textit{no}-modifier \textit{shujinkō} `main protagonist’ \ref{taro convenience}.\footnote{For \textit{i}-adjectives, in order to paraphrase \textit{i} with the corresponding \textit{ku-ar-u}, the right context is needed. See Section \Ref{copulaelements}.}

\ex. \label{nigiyaka convenience} \ag. nigiyaka-na machi\\
lively-\textsc{mod} city\\
\bg. nigiyaka-de-ar-u machi\\
lively-\textsc{pred-cop-prs} city\\
`a city which is lively’

\ex. \label{taro convenience} \ag. Tarō-ga shujinkō-no monogatari\\
Taro-\textsc{nom} main.protagonist-\textsc{mod} story\\
\bg. Tarō-ga shujinkō-de-ar-u monogatari\\
Taro-\textsc{nom} main.protagonist-\textsc{pred-cop-prs} story\\
`story in which Taro is the main protagonist' \\ \null \hfill \citep[66]{watanabe:2010:notes}

That \textsc{no} is not (always) an instance of the copula can be shown in the following minimal pair in which \textsc{no} can yield both an equative, as well as a possessive meaning, whereas \textsc{de-ar-u}, on the other hand, only causes the nominal modifier to have an equative reading, highlighting its predicative nature (see also \citealt{kato:2003:nihongo} among others).

\exg. sensei-no tomodachi\\
teacher-\textsc{mod} friend\\
1. $\approx$ `(my) friend who is a teacher' (equative)\\
2. $\approx$ `friend of the teacher' (possessive) \label{senseinotomodachi}

\exg. sensei-de-ar-u tomodachi\\
teacher-\textsc{pred-cop-prs} friend\\
1. $\approx$ `(my) friend who is a teacher' (equative)\\
2. $\neq$`friend of the teacher' (possessive)

Cross-linguistically, there is agreement that the copula is a licensor of non-verbal elements in predicative position, thus linking the predicate to its subject, and that it is in essence meaningless (\citealt[1–6]{arche:2019:main}, see \citealt{moro:1997:raising, moro:2000:dynamic, bowers:2002:pred} among others for historical overview). The appearance of the copula itself is thus strongly linked to indirect modification. On the other hand, ``non-paraphrasability with a copular construction" is a sign of direct modification \citep[291–292]{alex:2007:nounphrase}.

This means that when direct-only \textit{na}-adjectives and \textit{no}-modifiers are used attributively, \textsc{de-ar-u}, or \textsc{d(e)-at-ta}, should never occur instead of the different outputs of \is{MOD@\textsc{mod}}\textsc{mod} (see also \citealt{ayano:2010:revisiting}). This is shown below for a direct \textit{na}-adjective \ref{kanzen no cop} and a direct \textit{no}-modifier \ref{nihon no cop}. Note that this fact relates again to the lack of tense of direct modifiers and their smaller structure which does not contain projections for these complex elements.

\exg. *ano kanzen-de-ar-u baka\\
that complete-\textsc{pred-cop-prs} idiot\\
(`that idiot who is complete’) \null \hfill \citep[108]{ayano:2010:revisiting} \label{kanzen no cop}

\exg. *Nihon-de-ar-u kuruma\\
Japan-\textsc{pred-cop-prs} car\\
(`car which is Japanese’) \label{nihon no cop}

\subsubsection {The coordinative form DE} \label{coordination copula}

Besides the canonical and the polite form of the copula, direct-only modifiers also cannot co-occur with the coordinative form of the copula \textsc{de}. They are even illicit in coordinative constructions in which the first conjunct bares the copula form. \is{coordination} \is{coordinative copula}

\largerpage[-2]
In Japanese, different word classes use different methods of coordination, meaning there is no universal coordinator at the word level similar to English `\textit{and}'. Nouns are connected via the comitative/conjunction postposition -\textit{to} (or the postposition -\textit{ya} for inexhaustive coordination).\footnote{The comitative postposition -\textit{to} further differs from the other elements to be reviewed in that it can also appear after the last noun in a list format and even together with (case) particles. See examples and discussion in \citet{kuno:1973:structure} and \citet{hiraiwa:2014:constraining} against \citet{tatsumi:2018:splitting}.}

\exg. hon-to shinbun\\
book-and newspaper\\
`books and newspapers'     \hfill nominal coordination, exhaustive\par

\largerpage
Verbs (and generally inflecting word classes), on the other hand, are \is{coordination}coordinated via the suffix -\textsc{te}, which implies a consecution of actions and possibly additional connections between two clauses, for instance causality.\footnote{-\textsc{te} also functions as an aspectual marker and is often glossed \textsc{prog} or \textsc{ger(und)}. I will adopt the label \textsc{ger}, but note that this element has nothing to do with nominalization \textit{per se}. See \citet{hasegawa:1996:study, mihara:2013:tekei, nakatani:2013:predicate, nakatani:2016:complex}.} -\textsc{te} attaches to the nominalization form (\textit{renyōkei}) and can undergo various phonological assimilation processes. The element \textit{koron-de} in \ref{koronde} consists of the verb \textit{korob} `to fall’, inflected to the nominalization form \textit{korobi}, to which -\textsc{te} is added and phonological assimilation takes place.\footnote{The combination /bi te/ was present still in Classical Japanese, but has become /nde/ as a result of historical sound change (\is{onbin@\textit{onbin}}\textit{onbin}). This form can now be interpreted as analytical. In Modern Japanese, the voiced variant /de/ of the coordinative form appears after all consonant-stem verbs ending in a nasal or a voiced obstruent. A similar phonological assimilation process happens before the past marker -\textsc{ta} \citep{ito:2015:word}.} The nominalization form of vowel-stem verbs, such as \textit{tabe} `to eat’, is syncretic with the infinitival form and the coordinative marker is always just /te/ as shown in \ref{tabete}. \is{Classical Japanese}

\ex. \ag. Ken-wa koron-de nai-ta.\\
Ken-\textsc{top} fall-\textsc{ger} cry-\textsc{pst}\\
`Ken fell down and cried.’ \\ \null \hfill verbal coordination, consonant verb \citep[357]{hiraiwa:2012:mechanism} \label{koronde}
\bg. Bangohan-wo tabe-te o-furo-ni hait-ta.\\
dinner-\textsc{acc} eat-\textsc{ger} \textsc{hon}-bath-\textsc{ill} enter-\textsc{pst}\\
`(I) ate dinner and (afterwards) took a bath.' \\ \null \hfill verbal coordination, vowel verb \label{tabete}

\textit{i}-adjectives are coordinated via the same suffix, which attaches to the inflection form /ku/, in Section \ref{copulaelements} argued to be the \is{predicative copula}predicative copula. In more elaborate contexts, bare /ku/ (the bare predicative copula) can be used for coordination, although this is highly restricted and does rarely appear in spoken Japanese. See below. \is{coordinative copula}

\ex. \ag. taka-ku-te mazu-i okashi\\
expensive-\textsc{pred-ger} unpleasant-\textsc{mod} sweet\\
`expensive and terrible sweets' \\ \null \hfill \textit{i}-adjective, coordinated with -\textsc{te}
\bg. taka-ku mazu-i okashi\\
expensive-\textsc{pred} unpleasant-\textsc{mod} sweet\\
`expensive and terrible sweets' \\ \null \hfill \textit{i}-adjective, coordinated without -\textsc{te}

When \textit{na}-adjectives and \textit{no}-modifiers appear as the first \is{coordination}conjunct, they exhibit the element \textsc{de}, exemplified below for a \textit{na}-adjective.\footnote{The second conjunct, the modifier directly preceding the noun can never appear with -\textsc{te/de}.

\par

\ex. \ag. *mazu-i taka-ku-te okashi\\
unpleasant-\textsc{mod} expensive-\textsc{pred-ger} sweet\\
\bg. *benri-na jōbu-de kuruma\\
practical-\textsc{mod} robust-\textsc{pred} car\\

\par


A reviewer asks why a conjecture in which a coordinated modifier carries \is{MOD@\textsc{mod}}\textsc{mod} in addition to the coordinative form of the copula, so for instance *\textit{benri-de-na}, is ruled out. I hypothesize that this is for the same reason why combinations of \textsc{mod} with any form of the canonical copula (\textsc{de-ar-u}) are ruled out, although perhaps more opaque. It seems that the overt realization of the head of PredP, which is the predicative copula (cf. Section \ref{chapter pred}), blocks the overt occurrence of \textsc{mod} as well, the same as the head of IP.}

\exg. benri-de jōbu-na kuruma\\
practical-\textsc{pred} robust-\textsc{mod} car\\
`a practical and robust car'

This element is homophonous to the voiced allophone /de/ of the coordination suffix \textsc{te}, which occurs after verbs ending in a voiced obstruent or nasal. They might look like the same element, however, they should be kept apart and \textsc{de} co-occurring with \textit{na}-adjectives and \textit{no}-modifiers is, in fact, an instance of the copula. Several pieces of evidence for this are given in \citet{hiraiwa:2012:mechanism}, and see also \citet{bloch:1970:bloch}. The first piece of evidence comes from distributional patterns. As the \is{predicative copula}predicative copula \textsc{da} is mutually exclusive with the canonical copula \textsc{de-ar-u} reviewed above, the coordinative copula \textsc{de} corresponds also to a longer form: The coordinative form of the canonical copula \textsc{de-at-te}. Note the morphological setup, shown in \ref{set up de-at-te}, in which, again, the suffix -\textsc{te} is added to the nominalization form.\footnote{The historical nominalization form of verbs of the \textit{aru}-paradigm (which \textit{ar}- is the prime member of) was /ri/. Phonological assimilation has led to deletion of /r/ and gemination of /t/, so to a sound change from /ri te/ to /tte/. Again, the bare nominalization form \textsc{de-ar-i} can be used for coordination and behaves essentially like \textsc{de-at-te}.} \is{coordinative copula}

\exg. benri-de-at-te jōbu-na kuruma\\
practical-\textsc{pred-cop-ger} robust-\textsc{mod} car\\
`a car that is practical and robust' \label{set up de-at-te}

This means that both \textsc{de} and \textsc{de-at-te}, as well as -\textit{to}, appear in contexts where nouns are coordinated. However, the former two elements are used for \textit{clausal coordination} of a \textit{no}-modifier with another modifier, in the sense that both exhibit predicativity. The postposition -\textit{to} can never be used in clausal instances. This is illustrated below (and cf. examples in \citealt[363]{hiraiwa:2012:mechanism}). \is{coordinative copula} \is{coordination}

\ex. \ag. Hanako-wa gakusei-de (-at-te) Nihon-jin da.\\
Hanako-\textsc{top} student-\textsc{pred} (-\textsc{cop}-\textsc{ger}) Japan-person \textsc{cop}\\
`Hanako is a student and Japanese.’
\bg. Hanako-wa gakusei-de (-at-te) hon-wo yo-mu koto-ga suki-da\\
Hanako-\textsc{top} student-\textsc{pred} (-\textsc{cop}-\textsc{ger}) book-\textsc{acc} read-\textsc{prs} \textsc{comp-nom} like-\textsc{cop}\\
`Hanako is a student and likes reading books.’
\cg. *Hanako-wa gakusei-to hon-wo yo-mu koto-ga suki-da\\
Hanako-\textsc{top} student-and book-\textsc{acc} read-\textsc{prs} \textsc{comp-nom} like-\textsc{cop}\\
(`Hanako is a student and likes reading books.’)

\largerpage
More evidence comes from the connective element \textit{katsu}, which \citet[364]{hiraiwa:2012:mechanism} calls a \textit{sentential conjunction marker}. \textit{Katsu} necessarily adds an additional state of affairs, is restricted to written speech and can only follow nouns, \textit{na}-adjectives and verbs in their nominalization form. That it must express a predication relation is illustrated in the following sentence, where a pure coordination reading, in which the person under description is the teacher of both Ken and Naomi, cannot be achieved. The only possible interpretation is `this person is Ken and this person is also Naomi's teacher’.\footnote{This sets \textit{katsu} clearly apart from \textit{shi} which has to appear with verbs inflected for tense. See also \citet{fenger:2020:words}.}\is{coordinative copula}

\newpage
\exg. *Kono hito-wa Ken katsu Naomi-no sensei-da.\\
this person-\textsc{top} Ken katsu Naomi-\textsc{mod} teacher-\textsc{cop}\\
Under interpretation: `This person is Ken's and Naomi's teacher.’

The same holds true of -\textsc{de}. This element also shows a restriction to clausal, that is predicative, contexts, giving further evidence that it is connected to indirect modification.

\exg. Kono hito-wa Ken-de Naomi-no sensei-da.\\
this person-\textsc{top} Ken-\textsc{pred} Naomi-\textsc{mod} teacher-\textsc{cop}\\
1. $\neq$ `This person is Ken's and Naomi's teacher.’\\
2. only: $\approx$ `This person is Ken, and this person is Naomi’s teacher.’

The predicative, nature of \textsc{de} can be further highlighted via its interaction with \textit{katsu}. If two entities are predicated on one noun, one of these elements has to appear, or both. The meaning is unchanged. \is{coordinative copula} \is{coordination}

\ex. \ag. Kono hito-wa Ken-no sensei *(katsu) Naomi-no sensei-da.\\
this person-\textsc{top} Ken-\textsc{mod} teacher \textsc{katsu} Naomi-\textsc{mod} teacher-\textsc{cop}\\
\bg. Kono hito-wa Ken-no sensei-*(de) Naomi-no sensei-da.\\
this person-\textsc{top} Ken-\textsc{mod} teacher-\textsc{pred} Naomi-\textsc{mod} teacher-\textsc{cop}\\
\cg. Kono hito-wa Ken-no sensei-de (katsu) Naomi-no sensei-da.\\
this person-\textsc{top} Ken-\textsc{mod} teacher-\textsc{pred} \textsc{katsu} Naomi-\textsc{mod} teacher-\textsc{cop}\\
`This person is Ken's teacher and (also) Naomi's teacher.’ \\ \null \hfill \citep[364–365]{hiraiwa:2012:mechanism}

This shows that, clearly, \textsc{de} is an instance of the copula. This in turn predicts that non-predicative modifiers are banned from coordination with \textsc{de} (and also \textsc{de-at-te}), which is borne out, as illustrated below for the non-predicative \textit{na}-adjective \textit{kanzen-na} `complete’ and the \is{relational modifier}\textit{no}-modifier \textit{Nihon-no} `Japan(ese)’.

\ex.\ag. *Max-ga kanzen-de baka-da.\\
Max-\textsc{nom} complete-\textsc{pred} idiot-\textsc{cop}\\
\bg. *Max-ga kanzen-de-at-te baka-da.\\
Max-\textsc{nom} complete-\textsc{pred-cop-ger} idiot-\textsc{cop}\\
(`Max is complete and Max is an idiot.’)

\ex. \ag. *Nihon-de jōbu-na kuruma\\
Japan-\textsc{pred} robust-na car\\
\bg. *Nihon-de-at-te jōbu-na kuruma\\
Japan-\textsc{pred-cop-ger} robust-na car\\
(`a Japanese and robust car’)

An even more striking observation is that non-predicative modifiers cannot even appear in \is{coordination}coordinative structures inside the DP as the second conjunct, regardless of whether the modifier to their left can appear in predicative position or not. As shown in \ref{kirei second conjunct}, although \textit{kirei-na} `beautiful’ can appear predicatively and can be coordinated, it cannot appear in coordination structures with a non-predicative modifier such as \textit{Chūgoku-no} `Chinese’, causing the DP to be ungrammatical.

\exg. *kirei-de Chūgoku-no kabin\\
beautiful-\textsc{pred} China-\textsc{mod} vase\\
(`a beautiful and Chinese vase’) \label{kirei second conjunct}

\section{Direct modifiers in Japanese} \label{section direct modifiers}

The following diagnostics for direct modifiers in Japanese were established in the previous section: \textit{Lack of predication, non-gradability, lack of tense/inflection, no co-occurrence with canonical copula} and \textit{lack of coordination}. These tests will help to identify the relevant groups of direct modifiers. I will start with a discussion of those groups of direct modifiers that are cross-linguistically attested, and finally come to the language-specific group of \textit{rentaishi}.

Recall that with the notion \textit{direct-only}, I refer to those modifiers that truly can never appear in predicative position. For many of the modifiers discussed in the following pages, predicative use, and this extends to the other characteristics, can be achieved. For example, the \is{relational modifier}relational modifier \textit{Chūgoku} is direct when denoting an origin, `Chinese, from China’. Its principal meaning, denoting something related to the country `China’, however, is possible in predicative position, which is why modifiers of \textsc{nationality} are often attested in the corpus in combination with the copula. As mentioned, group 5 of the data set contains some modifiers for which predicative use is marginal, but nevertheless attested, meaning those are not direct-only. Reflecting the obvious skew, which needed to be at least 5$>$1, towards attributive use, I classified such modifiers as \textit{prototypically attributive} and also sorted them to this group (cf. Section \ref{corpusanalysis}).

\subsection{Non-intersective modifiers} \label{non-intersective Japanese}

A cross-linguistically prominent group of direct-only modifiers, outlined in Section \ref{intersective}, is \is{non-intersective modifier}non-intersective modifiers of the type \textit{former}. Those are characterized by the missing set intersection between their denotation and that of the noun they modify and can be sub-classified into \textit{truly non-intersective modifiers}, which can never be used predicatively, \textit{subsective} modifiers, which can be used predicatively but then exhibit an intersective meaning, \textit{privative} modifiers and \textit{adnominal degree modifiers} \citep{vendler:1957:verbs, bolinger:1967:adj, larson:1998:events, morzycki:2016:mod}.

\subsubsection{Non-intersective adjectives}

Several \is{non-intersective modifier}non-intersective adjectives in Japanese are presented in the PhD thesis of \citet{yamakido:2005:nature} (cf. also \citealt{yamakido:2000:japanese, yamakido:2007:nature}), who argues explicitly against a unified relative clause analysis of Japanese adjectives, and in \citet{ayano:2010:revisiting}. The authors show that \is{intersectivity}ambiguous examples of the type \textit{old friend} in English are disambiguated via two separate lexemes in Japanese. The \textit{i}-adjective \textit{furu-i} can be translated as `old’, but cannot be used to denote the age of a human entity. In other words, it only gives a subsective reading when modifying human nouns, where its intersective interpretation is blocked \citep[66–67]{yamakido:2005:nature}. In order to denote a person's age, other lexemes have to be used, such as \textit{kōrei-na} `old’. This is illustrated below with the relevant notation.\footnote{As \citet[511, Footnote 7]{watanabe:2012:direct} notes, the phrase \textit{furu-kara-no} old-from-\textsc{mod} `from a long time ago' might be a better and more natural alternative to \textit{furu-i}. See also \citet[155–156]{tatsuki:2009:nihongo}.}

\exg. furu-i tomodachi\\
old-\textsc{mod} friend \null \hfill \citep[81]{yamakido:2005:nature}\\
1. $\approx$ `long-time friend’ [furu-i tomodachi] $\subseteq$ [tomodachi] (subsective)\\
2. $\neq$ `aged friend’ [furu-i tomodachi] = [furu-i] $\cap$ [tomodachi] (intersective)

\exg. kōrei-na tomodachi\\
old-\textsc{mod} friend \null \hfill \citep[67]{yamakido:2005:nature}\\
1. $\neq$ `long-time friend’ [kōrei-na tomodachi] $\subseteq$ [tomodachi] (subsective)\\
2.$\approx$ `aged friend’ [kōrei-na tomodachi] = [kōrei-na] $\cap$ [tomodachi] (intersective)

Lack of intersective interpretation suggests lack of predicative use, which in turn suggests direct modification status \citep{larson:1998:events, cinque:2010:adj}. As expected, this adjective cannot be used predicatively, at least not with human entities (cf. \citealt[67]{yamakido:2005:nature}).

\exg. *Watashi-no tomodachi-wa furu-i.\\
I-\textsc{mod} friend-\textsc{top} old-\textsc{cop}\\
(`My friend is old (aged).’)

Another case of \is{intersectivity}ambiguity resolution presents itself in the \textit{beautiful dancer} example. \citet[3]{baker:2003:verbal} is right when he says that the direct equivalent \textit{utsukushi-i / kirei-na kashu} `beautiful singer’ only has an intersective reading, denoting the beauty of a person irrespective of their singing skills. However, other lexemes allow only a \is{non-intersective modifier}non-intersective, that is a subsective, interpretation of beauty pertaining to the dancing or singing skills of a person. Such lexemes are \textit{takumi-na} `skillful’ \citep[101, Footnote 2]{ayano:2010:revisiting}, \textit{migoto-na} \citep[786]{watanabe:2017:attributive}, which can be translated as `beautiful’, but can only be used with reference to events, or \textit{suteki-na} `wonderful'. This contrast is shown below.

\exg. utsukushi-i kashu\\
beautiful-\textsc{mod} singer \null \hfill \citep[3]{baker:2003:verbal}\\
1. $\neq$ `beautiful as a singer’ [beautiful singer] $\subseteq$ [singer] (subsective)\\
2. $\approx$ `beautiful and a singer’ [beautiful singer] = [beautiful] $\cap$ [singer] (intersective)

\exg. Hanako-wa takumi-na / migoto-na / suteki-na kashu-da.\\
Hanako-\textsc{top} skilfull-\textsc{mod} / beautiful-\textsc{mod} / wonderful-\textsc{mod} singer-\textsc{cop}\\
1. $\approx$ `beautiful as a singer’ [takumi-na kashu] $\subseteq$ [kashu] (subsective)\\
2. $\neq$ `beautiful and a singer’ [takumi-na kashu] = [takumi-na] $\cap$ [kashu] (intersective)

This ambiguity resolution via two different lexemes points towards the fact that one is a direct, and the other an indirect modifier. Similar ambiguity resolutions have been observed in Korean \citep[161]{kang:2005:adj}, and German \citep[69]{maienborn:2020:revisiting}. \is{German} \is{Korean}

Japanese also has truly non-intersective adjectives, among those the \textit{i}-adjective \textit{omoigakena-i} `unexpected' and the \textit{na}-adjective \textit{kanzen-na} `complete' \citep[67–68]{yamakido:2005:nature}, as well as the \textit{i}-adjective \textit{tayumina-i} `non-ceasing' and the \textit{na}-adjective \textit{sakan-na} `vigorous' given by \citet{ayano:2010:revisiting}. \textit{Omoigakena-i} `unexpected', for example, can only express an ``external adverbial reading" \citep[64]{yamakido:2005:nature}. Below I illustrate the ban on predicative use, the corresponding adverbial reading and the missing intersectivity.

\ex. \ag. Omoigakena-i kyaku-ga ki-ta.\\
unexpected-\textsc{mod} guest-\textsc{nom} come-\textsc{pst}\\
`An unexpected guest came.' \null \hfill \citep[67]{yamakido:2005:nature}
\bg. *Ano kyaku-ga omoigakena-i.\\
that guest-\textsc{nom} unexpected-\textsc{cop}\\
(`That guest is unexpected.’)
\cg. Omoigakena-ku kyaku-ga ki-ta.\\
unexpected-\textsc{adv} guest-\textsc{nom} come-\textsc{pst}\\
`Unexpectedly, a guest came.'
\d. [omoigakena-i kyaku] $\neq$ [omoigakena] $\cap$ [kyaku]

As mentioned in the previous section, a ban on predicative use goes alongside a ban on the canonical copula \textsc{de-ar-u} \ref{kanzendearu} \citep[108--109]{ayano:2010:revisiting} and the \is{coordinative copula}coordinative form of the copula \ref{kanzendeari}, shown for direct-only \textit{kanzen-na} `complete’.\footnote{Some corpus results contain the attributive construction \textit{kanzen-naru} + NP. In this case, the attributive form (\textit{rentaikei}) of the premodern copula \textit{nari} occurs, thus this is still a case of attributive, not predicative use. See the discussion in Section \ref{complexdirect}.}

\ex. \ag. Max-ga kanzen-na baka-da.\\
Max-\textsc{nom} complete-\textsc{mod} idiot-\textsc{cop}\\
`Max is a complete idiot.' \null \hfill \citep[68]{yamakido:2005:nature}
\bg. *Ano baka-ga kanzen-da.\\
that idiot-\textsc{nom} complete-\textsc{cop}\\
(`That idiot is complete.’)
\cg. *ano kanzen-de-ar-u baka\\
that complete-\textsc{pred-cop-prs} idiot \label{kanzendearu}\\
(`that idiot who is complete’)
\dg. *ano kanzen-de / de-ar-i / de-at-te kinpatsu-na baka\\
that complete-\textsc{pred} / \textsc{pred-cop-ryk} / \textsc{pred-cop-ger} blond-\textsc{mod} idiot \label{kanzendeari}\\
(`that complete and blond idiot’)

Finally, the absence of tense is also attested, illustrated with \is{non-intersective modifier}\textit{furu-katta} `was old' (cf. also example in \citealt[109]{ayano:2010:revisiting}).

\exg. *Watashi-no tomodachi-wa furu-k-at-a.\\
I-\textsc{mod} friend-\textsc{top} old-\textsc{pred-cop-pst}\\
(`My friend was old (aged).’)

However, one needs to be careful not to draw early conclusions from these results as each of those adjectival lexemes exhibits certain particularities. First, as \citet{ayano:2010:revisiting} admits, the lexemes \textit{omoigakena-i} `unexpected', \textit{tayumina-i} `non-ceasing’ and \textit{takumi-na} `skillful' can in fact be inflected into past and then even be used predicatively. I report his judgments below.

\ex. \ag. Ano kyaku-ga omoigakena-k-at-ta.\\
that guest-\textsc{nom} unexpected-\textsc{pred-cop-pst}\\
`That guest was unexpected.'
\bg. ?Shinpo-ga tayumina-k-at-ta.\\
progress-\textsc{nom} non.ceasing-\textsc{pred-cop-pst}\\
`Progress was non-ceasing.'
\cg. ?Ano utaite-ga takumi-d-at-ta.\\
that singer-\textsc{nom} skillful-\textsc{pred-cop-pst}\\
`That singer was skillful.' \null \hfill \citep[109]{ayano:2010:revisiting}

\citet{watanabe:2017:attributive} alludes to the very high CP status of the adverbial counterpart of \textit{omoigakena-i} and shows that this adjective can appear above numerals. As outlined in Section \ref{japanesedp}, the functional projection hosting numerals is inside the indirect domain, therefore direct modifiers should not be able to cross this border.\footnote{\citet{hoshi:2002:kaynean} adds that predicative use can be ameliorated when modifiers are contrasted, although recall the examples at the beginning of this chapter for why this might be an orthogonal issue.}

\exg. Omoigakena-i futari-no kyaku-ga soko-ni i-ta.\\
unexpected-\textsc{mod} two-\textsc{mod} guest-\textsc{nom} there-\textsc{loc} be-\textsc{pst}\\
`Two unexpected guests were there.’ \null \hfill \citep[790]{watanabe:2017:attributive}

A further potential caveat is that many of the \is{non-intersective modifier}lexemes presented above have a rather peripheral lexematic status. \textit{Furu-i} is banned from a combination with human entities in referring to age, possibly for pragmatic reasons.\footnote{The concept of honorificity might forbid this combination. As \citet[123]{nagano:2015:RA} point out, \textit{furu-i} does not correspond to \textit{waka-i} `young' but rather to \textit{atarashi-i} `new'.} \textit{Migoto-na} `beautiful' also cannot be used in combination with human entities, but only refers to events \citep[786]{watanabe:2017:attributive}. \textit{Kanzen-na}, on the other hand, might not even be a lexical adjective but what Morzycki calls an \is{adnominal degree modifier}`adnominal degree modifier' (also: \is{adnominal degree modifier}`adnominal degree morpheme'; \citealt{morzycki:2009:degree, morzycki:2012:adj, morzycki:2012:faces}).\footnote{In fact, \citet{morzycki:2014:where} uses Japanese \textit{kanzen-na} as one example for this category of modifiers. Also recall that the corresponding adverbial form \textit{kanzenni} `completely' is an endpoint-oriented modifier \citep{oshima:2019:gradability}.} \textit{Tayumina-i} `non-ceasing' and \textit{omoigakena-i} `unexpected' end in the suffix -\textit{na} for negation, which inflects like an \textit{i}-adjective.\footnote{\textit{Tayumina-i} is derived from \textit{tayumi}, which is a noun but has a marginal status in Modern Japanese. \textit{Omoigakena-i} has lexicalized from the negative form of the verb \textit{omoigakeru} `to expect’ \citep[803]{watanabe:2017:attributive}.}

These modifiers therefore seem to constitute a special case, and it becomes obvious that there are only a few non-intersective modifiers among Japanese \textit{i}- and \textit{na}-adjectives.

\subsubsection{Non-intersective \textit{no}-modifiers}

In fact, most non-intersective lexemes -- the equivalents of well-known examples such as \textit{former, alleged} or \textit{future} -- are found among Japanese \textit{no}-modifiers. \citet[122]{nagano:2015:RA} give a few examples, but many more can be found in the data set. Below are given examples of truly \is{non-intersective modifier}non-intersective modifiers, and \is{adnominal degree modifier}adnominal degree modifiers \citep{morzycki:2014:where}, for all of which predicative use was limited or non existent.\footnote{Again, speakers might have the feeling that predicative use can be achieved to different extents via establishing contrastive contexts. Not only were occurrences low to non-existent in the corpus, predicative use is always tested for context-neutral utterances, as far as that is possible.}

\ex. \emph{Non-intersective modifiers in Japanese with element \textsc{no}}
\a. \textit{yui'itsu-no} `sole, unique'
\b. \textit{jūzen-no} `previous, former’
\c. \textit{tenpu-no} `innate’
\d. \textit{shōbi-no} `urgent’
\e. \textit{shōrai-no} `future'

\ex. \emph{Adnominal degree modifiers in Japanese with element \textsc{no}}
\a. \textit{mattaku-no} `complete(ly)'
\b. \textit{hontō-no} `true(ly), real(ly)'

These modifiers cannot be used predicatively, hence fulfill the main criterion of direct modification. This is illustrated below for \textit{shōrai-no} `future’ and \textit{tenpu-no} `innate’.

\largerpage
\ex. \ag. shōrai-no daitōryō\\
future-\textsc{mod} prime.minister\\
`a future prime minister’
\bg. *Kono daitōryō-wa shōrai-da.\\
this prime.minister-\textsc{top} future-\textsc{cop}\\
(`This prime minister is future.’)

\ex. \ag. tenpu-no sainō\\
innate-\textsc{mod} talent\\
`innate talent’
\bg. *Kare-no sainō-wa tenpu-da.\\
he-\textsc{mod} talent-\textsc{top} innate-\textsc{cop}\\
(`His talent is innate.’)

These modifiers are also restricted in \is{gradability}gradability, first illustrated via combinations with the canonical degree adverb and comparative constructions, either with \textit{motto} `more’ or the particle \textit{yori} `than’.

\ex. \ag. *totemo shōrai-no daitōryō\\
very former-\textsc{mod} prime.minister\\
(`a very future prime minister’)
\bg. *motto shōrai-no daitōryō\\
more future-\textsc{mod} prime.minister\\
(`a more future prime minister’)
\cg. *Kono daitōryō-wa ano daitōryo-yori shōrai-da.\\
this prime.minister-\textsc{top} that prime.minister-than future-\textsc{cop}\\
(`This prime minister is more future than that prime minister.’)

Above, it was shown that the degree adverb \textit{kanari} `quite, rather' is more permissive, however the combination of \textit{kanari} with most \is{non-intersective modifier}non-intersective modifiers seems to be out as well, exemplified below for \textit{shōrai-no} `future’, \textit{yui’itsu-no} `unique’ and the \is{adnominal degree modifier}nominal degree modifier \textit{hontō-no} `true’.

\ex. \ag. *kanari shōrai-no daitōryō\\
quite future-\textsc{mod} prime.minister\\
(`a quite future prime minister’)
\bg. *kanari yui'itsu-no hōhō\\
quite sole-\textsc{mod} method\\
(`a quite sole method’)
\cg. *kanari hontō-no tomodachi\\
quite true-\textsc{mod} friend\\
(`a quite true friend’)

Encoding a tense interpretation is also ungrammatical.

\exg. *shōrai-d-at-ta daitōryō\\
future-\textsc{pred-cop-pst} prime.minister\\
(`a prime minister who was future’)

Finally, there is one \is{privative modifier}privative modifier, namely \textit{nise-no} `fake’, for which predicative use is degraded but not entirely ungrammatical, as is the case also with the corresponding \textit{fake} in English \citep{cinque:2014:semantics}.

\exg. ?Kono pisutoru-wa nise-da.\\
this gun-\textsc{top} fake-\textsc{cop}\\
`This gun is fake.'

This shows that \is{non-intersective modifier}non-intersective modifiers do exist in Japanese and most of them appear with the element \textsc{no}. They exhibit many characteristics of their cross-linguistic counterparts: they cannot be used predicatively, are not gradable and are not complex, all features that make them candidates for direct-only modifiers.\footnote{Other instances of non-intersective modifiers can be found among the group of \is{rentaishi@\textit{rentaishi}}\textit{rentaishi} discussed in Section \ref{rentaishi}.}

\subsubsection{Hierarchies and ordering} \label{ordering intersective}

Recall from Section \ref{svenonius} that \citet{svenonius:2008:position} argues that it is only modifiers of different \is{gradability}gradability and intersectivity that exhibit \is{adjective ordering restrictions}ordering preferences and assigns two different projections, SortP for subsective, but gradable modifiers, and nP for non-gradable intersective modifiers (see also \citealt{dekany:2011:profile, truswell:2009:attributive, kim:2019:syntax}). He predicts that the latter type are closer to the noun (in linear order, shown through the symbol `$\succ$'). The examples below illustrate this point.

\ex. \a. tiny red hats (gradable/subsective $\succ$ non-gradable/intersective)
\b. *?red tiny hats (non-gradable/intersective $\succ$ gradable/subsective) \\ \null \hfill\citep[36]{svenonius:2008:position}

But what about truly non-intersective modifiers? In \is{German} German, \textit{früher} `former’, \textit{heutig} `today's' or \textit{gestrig} `yesterday’s’ among others, systematically have to precede intersective and subsective adjectives (see corpus study in \citealt{trost:2006:das}, as well as \citealt{eichinger:1991:ganz}).\footnote{In German linguistics, such modifiers are also called `\textit{referentiell}', meaning `referential’ \citep{trost:2006:das}, or `Situativa' \citep{eichinger:1979:syntaktische, eichinger:1991:ganz}, because they situate a referent in space and time.}

\ex. \ag. der gestrige schöne Tag (German)\\
the yesterday beautiful day\\
`yesterday's beautiful day'
\bg. *der schöne gestrige Tag\\
the beautiful yesterday day\\

Structural precedence of non-intersective adjectives has been tested for \is{Greek} Greek and \is{English} English by \citet{panayidou:2013:in} based on scope effects. Assuming both adjectives are direct, she detects differences in scope between the ordinary \is{intersectivity}intersective adjective \textit{bright} and the truly non-intersective adjective \textit{former} in English depending on their order. When both adjectives modify the noun \textit{student} in the order \textit{former} $\succ$ \textit{bright}, two readings are possible: A reading where \textit{former} takes \is{scope}scope over the combination of \textit{bright} and the noun \textit{student} at the same time, denoting a student who is no longer bright, and a reading where both modifiers take individual scope over the noun and \textit{former} skips the intervening adjective, denoting a person who was once a student and is or was bright, irrespective of their affiliation \citep[106]{panayidou:2013:in}.

\ex. the former bright student (ambiguous)\\
1. $\approx$ `a student who is no longer bright’\\
2. $\approx$ `a person who is no longer a student and is/was bright’

If the order is reversed to \textit{bright} $\succ$ \textit{former}, \textit{bright} cannot modify solely the noun without skipping over the intervening adjective and only the first reading is \is{scope}available \citep[106]{panayidou:2013:in}.

\ex. the bright former student (only reading 1)\\
1. $\approx$ `a bright person who is no longer a student’\\
2. $\neq$ `a person who was bright as a student and is no longer one’

This shows that we are dealing with a marked \is{adjective ordering restrictions}order. Since a common reason for marked orders -- besides different kinds of modification (which is refuted in \citealt[§ 3.5.2]{panayidou:2013:in}) -- is movement out of canonical positions, the author concludes that \textit{former} $>$ \textit{bright}, and therefore \textit{non-intersective} $>$ \textit{intersective} is the underlying structural hierarchy \citep[107–111]{panayidou:2013:in}. This yields the following preliminary structural hierarchy based on intersectivity.

\ex. Hierarchy based on Intersectivity \citep{trost:2006:das, svenonius:2008:position, panayidou:2013:in}\\
Non-intersective $>$ Subsective $>$ Intersective $>$ N

On the contrary, \citet{ayano:2010:revisiting} argues that \is{non-intersective modifier}non-intersective adjectives in Japanese have to appear lower than other adjectives \ref{ordering1} as well as lower than \is{numeral}numerals \ref{ordering2} as in the following examples (\textit{furu} will be glossed as `long.time' now to avoid confusion).

\ex. \label{ordering1} \ag. monoshizuka-na furu-i yūjin\\
quiet-\textsc{mod} long.time-\textsc{mod} friend\\
`quiet long-time friend' \hfill \citep[104]{ayano:2010:revisiting}
\bg. *furu-i monoshizuka-na yūjin\\
long-time-\textsc{mod} quiet-\textsc{mod} friend\\
\cg. monomidaka-i kanzen-na baka\\
burningly.curious-\textsc{mod} complete-\textsc{mod} idiot\\
`burningly curious complete idiot' \hfill \citep[105]{ayano:2010:revisiting}
\dg. *kanzen-na monomidaka-i baka\\
complete-\textsc{mod} burningly.curious-\textsc{mod} idiot\\

\ex. \label{ordering2} \ag. San-nin-no furu-i yūjin-ga ki-ta.\\
three-\textsc{cl}-\textsc{mod} long.time-\textsc{mod} friend-\textsc{nom} come-\textsc{pst}\\
`Three long-time friends came.' \hfill \citep[106]{ayano:2010:revisiting}
\bg. *Furu-i san-nin-no yūjin-ga ki-ta.\\
long.time-\textsc{mod} three-\textsc{cl}-\textsc{mod} friend-\textsc{nom} come-\textsc{pst}\\
\cg. futari-no kanzen-na baka\\
two-\textsc{mod} complete-\textsc{mod} idiot\\
`two complete idiots' \hfill \citep[107]{ayano:2010:revisiting}
\dg. *kanzen-na futari-no baka \\
complete-\textsc{mod} two-\textsc{mod} idiot\\

From this, he concludes that non-intersective modifiers in Japanese are systematically in a lower domain than all other direct modifiers and presents his own hierarchy. \is{intersectivity}

\ex. Hierarchy for Japanese based on Intersectivity \citep{ayano:2010:revisiting}\\
Intersective $>$ Non-Intersective $>$ N

\citet{watanabe:2017:attributive} adds the following \is{adjective ordering restrictions}piece of data regarding the order of \textit{utsukushi-i} (intersective) and \textit{migoto-na} (subsective), those modifiers that divide the meaning of `beautiful' as described above. He judges the \is{non-intersective modifier}subsective modifier best closest to the noun but mentions that this contrast ``may not be as sharp as we want" \citep[790]{watanabe:2017:attributive}.

\ex. \ag. utsukushi-i migoto-na utaite\\
beautiful-\textsc{mod} wonderful-\textsc{mod} singer\\
`beautiful person who sings beautifully/wonderfully' \label{utsukushi-i migoto-na utaite}
\bg. ??migoto-na utsukushi-i utaite\\
wonderful-\textsc{mod} beautiful-\textsc{mod} singer \null \hfill \citep[790]{watanabe:2017:attributive} \\

The observations in \citet{ayano:2010:revisiting} and \citet{watanabe:2017:attributive} can, however, be explained by adhering to the direct/indirect dichotomy, rather than a structural hierarchy inside the direct domain. First, that non-intersective modifiers appear lower than \is{numeral}numerals is expected since numerals are merged in the indirect domain, explaining \ref{ordering2}. The fact that the \is{non-intersective modifier}non-intersective modifier \textit{kanzen-na} `complete’ is best closer to the noun in \ref{ordering1} follows naturally if we assume that \textit{monoshizuka-na} `quiet’ and \textit{monomidaka-i} `burningly curious’ are indirect.\footnote{There might even be a completely different reasons at stake here, namely phonological length or prosody. \textit{Monoshizuka-na} `quiet' is twice as long as \textit{furu-i} `long time' (6 vs. 3 morae), and it is not clear why the author did not use the more common \textit{shizuka-na} `quiet' with only 4 morae. However, in my survey, I could not find significant results for the effect of mora size on ordering as reported in Section \ref{optionality}.} Concerning the example in \citet{watanabe:2017:attributive}, the order \textit{utsukushi-i} $\succ$ \textit{migoto-na} is also expected under the assumption that these modifiers are indirect and direct respectively, in line with \citet[22–24]{cinque:2010:adj}. Similar to the \textit{beautiful beautiful dancer} or \textit{visible visible stars} examples (Section \ref{cartography Cinque}), the modifier closer to the noun is direct. Their respective positions in the DP are illustrated in \figref{Structure of utsukushi-i migoto-na utaite} (omitting ModP for better readability).

\begin{figure}
\centering
\begin{tikzpicture}[baseline=0pt, remember picture]
\tikzset{level distance=25pt, sibling distance=5pt}
\Tree [.DP\textsubscript{internal} {} [.D′\textsubscript{internal} [.FP\textsubscript{indirect} [.AP \edge[roof]; utsukushi-i ] [.F′\textsubscript{indirect} [.dP [.FP\textsubscript{direct} [.AP \edge[roof]; migoto-na ] [.F′\textsubscript{direct} [.NP \edge[roof]; utaite ] F ] ] d ] F ] ] D ] ]
\end{tikzpicture}
\caption{Structure of \ref{utsukushi-i migoto-na utaite}}
\label{Structure of utsukushi-i migoto-na utaite}
\end{figure}

The same structural positions are at stake in the following minimal pair. The non-intersective modifier \textit{shōrai-no} `future' must be closer to the noun than the \textit{na}-adjective \textit{majime-na} `hardworking'. Again, an indirect, or at least ambiguous, modifier must precede the direct-only modifier showing its higher position within the DP.

\ex. \ag. majime-na shōrai-no purosakkāsenshu\\
hardworking-\textsc{mod} future-\textsc{mod} football.player\\
`a hardworking person who will be a professional football player in the future’
\bg. ??shōrai-no majime-na purosakkāsenshu\\
future-\textsc{mod} hardworking-\textsc{mod} pro.football.player\\

Crucially, a hierarchy based on \is{intersectivity}intersectivity can hardly be established in Japanese and all these examples seem rather to involve a direct/indirect dichotomy. Recall from Chapter \ref{japanesedp} that truly non-intersective modifiers also do not engage in \is{absence of ordering restrictions}ordering hierarchies with other types of direct modifiers, for example the equivalents of \is{relational modifier}relational adjectives, to be reviewed below. This is exemplified in \ref{tada} with the \is{non-intersective modifier}non-intersective modifier \textit{tada-no} `mere’. \ref{tadada} shows that this modifier cannot occur predicatively.\footnote{Only with the meaning `mere’ is it never possible in predicative position. \textit{Tada-no} can also have the meaning `for free’, in which case it can be used predicatively. Therefore it is sorted in \is{scope}group 1 in the data set.}

\ex. \label{tada} non-intersective vs. relational modifier
\ag. tada-no Doitsu-no gakusei\\
mere-\textsc{mod} German-\textsc{mod} student \label{tada1}\\
\bg. Doitsu-no tada-no gakusei\\
German-\textsc{mod} mere-\textsc{mod} student\\
`a mere German student' \label{tada2}

\exg. *Kono gakusei-wa tada-da.\\
this student-\textsc{top} mere-\textsc{cop}\\
(`This student is mere.’) \label{tadada}

This discussion has shown that non-intersective modifiers do exist in Japanese and that they predominantly can be found in the group of \textit{no}-modifiers. Furthermore, it was shown that the well-known ambiguous cases in \is{English} English are resolved in Japanese, which is a substantial problem for \citet{baker:2003:verbal} who correlates the apparent lack of non-intersectivity with a lack of direct modification. Finally, since non-intersective modifiers are direct, they are merged lower than indirect modifiers and appear closer to the noun, explaining the mismatches between the assumptions of \citet{svenonius:2008:position} and \citet{ayano:2010:revisiting}.

\subsection{Relational modifiers}
\label{RAJapanese}

The second cross-linguistic group of direct-only modifiers are \is{relational modifier}\textit{relational adjectives}, although in Japanese it might be better to speak of \textit{relational modifiers}. Those modifiers were discussed in Section \ref{RA} with regard to European languages (cf. \citealt{fabregas:2007:internal, rainer:2013:can, bortolotto:2016:RA}).

Japanese, in addition to \is{Arabic} Arabic, was one of the two Non-European languages looked at to verify this category cross-linguistically in \citet{bisetto:2010:RA}, which explains the surprising amount of literature on this category with regard to Japanese \citep{nagano:2015:RA, nagano:2016:are, nagano:2018:conversion, shimada:2018:relational, morita:2010:internal, morita:2011:three, morita:2013:morph}. \citet{bisetto:2010:RA} focuses on derivational factors and analyzes the suffix morpheme -\textit{teki} \textjpn{的} `-ic' as the derivational suffix dedicated to form relational adjectives in Japanese.\footnote{Examples include \textit{keizai-teki} `economical' derived from \textit{keizai} `economy' or \textit{shukan-teki} `subjective' derived from \textit{shukan} `subjectivity’, cf. also \citet[71–73]{bisetto:2010:RA}.} However, as shown in \citet[48–50]{nagano:2016:are}, these derivates are not true equivalents of relational adjectives, since they in fact can be used predicatively and are gradable. Rather, the equivalents of relational adjectives appear in Japanese unequivocally with the element \textsc{no} and are not subject to morphological processes. Those equivalents seem to be nominal not only on the deep, but also on the surface level.

In the following, I will first show that the equivalents of \textit{thematic} and \textit{classificatory} relational adjectives exist in Japanese and then give an overview of the categories named in \citet{nagano:2016:are}, discussing some of them in sufficient detail in order to show which are candidates for direct modifiers, and where deviations might occur.

\subsubsection{Thematic and classificatory relational modifiers}

In Section \ref{RA}, it was outlined that relational adjectives in European languages come in two types: \textit{classificatory} and \textit{thematic}. Only the equivalents of classificatory relational adjectives are discussed in the literature with regard to Japanese. An example is given below. Observe that, contrary to the English equivalent \textit{environmental crisis}, no derivational processes are observable on the morphological level.

\exg. kankyō-no kiki \\
environment-\textsc{mod} crisis\\
`environmental crisis’

As observable above, the equivalents of the \is{relational modifier}classificatory adjective \textit{environmental} is a \textit{no}-modifier which expresses a certain relationship to the noun it modifies. Crucially however, \textit{kankyō} is a noun itself. Besides the absence of morphological derivation, this can also be exemplified on the syntactic level via productive use with case particles.

\exg. Kankyō-ga daiji-da.\\
environment-\textsc{nom} important-\textsc{cop}\\
`The environment is important.’

As far as competing constructions go, \textit{of}-phrases (\textit{problems of the economy}) or PPs are not available, yet Japanese makes frequent use of lexical compounds, often Sino-Japanese four character lexical compounds, with the same meaning.

\exg. kankyō-kiki (\textjpn{環境危機})\\
environment-crisis\\
`environmental crisis’

Thematic relational adjectives have not been identified as a separate type of relational modifiers with regard to Japanese, probably because such modifiers are nominals in the first place. Recall that thematic relational adjectives are, in languages such as English, restricted to the modification of event or action nouns \citep{grimshaw:1991:argument}. A direct equivalent of such examples would be \ref{machi hatten} in which \textit{machi-no} `urban’ could be interpreted as a \textit{thematic relational modifier} since it expresses a subject role.

\exg. machi-no hatten\\
city-\textsc{mod} development\\
`urban development / development of the city’ \label{machi hatten}

Event nouns in Japanese usually belong to the category of so-called \textit{Verbal Nouns} \citep{teramura:1982:nihongo, oshima:2019:gradability}. Such nouns can either appear as true nouns, as in \ref{machi hatten}, or as predicates of a clause together with the light verb \textit{su} `to do’. \ref{machi-ga hatten} is a paraphrase of \ref{machi hatten}, in which \textit{hatten} serves as the predicate of the matrix clause and \textit{machi} now appears as a noun with the nominative case particle.

\exg. Machi-ga hatten-shi-ta.\\
city-\textsc{nom} development-do-\textsc{pst}\\
`The city developed.’ \label{machi-ga hatten}

Since such modifiers appear very nominal, no special discussion on their character seems necessary at first. Nevertheless, in Japanese, neither the equivalents of classificatory \ref{kiki}, nor of thematic \is{relational modifier}relational adjectives \ref{machi} can serve as the predicate of a sentence together with the copula.

\ex. \ag. *Kiki-ga kankyō-da.\\
crisis-\textsc{nom} environment-\textsc{cop}\\
(`The crisis is environmental.’) \label{kiki}
\bg. *Hatten-ga machi-da.\\
development-\textsc{nom} city-\textsc{cop}\\
(`The development is urban.’) \label{machi}

\subsubsection{Different categories}
\largerpage[1.5]
A further obstacle for the precise analysis of relational modifiers in Japanese is the abundance of possible categories scattered in the literature, for which it is not clear if all really fulfill the relevant syntactic criteria. A very detailed list of categories is presented in \citet[53–54]{nagano:2016:are}, who understands relational adjectives in Japanese as non-predicative modifiers with a nominal character. In the table below, I give her categories and compare them to the categories presented in hierarchies of relational adjectives written in the cartographic research field \citep{rae:2010:ordering, bortolotto:2016:RA}, to older approaches to relational adjectives in \is{English} English (\citealt{downing:1977:creation, levi:1978:syntax, warren:1984:classifying}), to standard hierarchies of qualitative adjectives, and to the PP hierarchy of \citet{takamine:2010:hierarchy}.\footnote{This
    table is most of all meant to serve as a reference. See also \citet[93–94]{bortolotto:2016:RA} as well as \citet[16]{rainer:2013:can} for similar comparisons.}

\begin{table}[h]
\caption{Categories of Japanese relational modifiers compared to other studies}
\begin{tabularx}{\textwidth}{lQ@{}p{3cm}}
\lsptoprule
\textbf{Nagano (2016)} & \textbf{In other hierarchies} & \textbf{Example} \\
\midrule

Material & \citet{scott:2002:stacked, laenzlinger:2005:french, laenzlinger:2011:elements}: \textit{Material} (QA)\newline
\citet{rae:2010:ordering, bortolotto:2016:RA}: \textit{Material} (RA) & \textit{komugi-no pan} `wheat(en) bread' \\
& \citet{takamine:2010:hierarchy}: \textit{Material} (PPs) & \\
\tablevspace
Nationality & \citet{sproat:1991:the}: \textit{Provenance} (QA)\newline
 \citet{scott:2002:stacked, cinque:2010:adj, laenzlinger:2005:french, laenzlinger:2011:elements}: \textit{Nationality/Origin} (QA) \newline
\citet{alex:2011:ea, arseni:2014:EA}: \textit{ethnic adjectives} & \textit{Chūgoku-no kabin} `Chinese vase'  \\
\midrule
\end{tabularx}
\end{table}
\clearpage

\begin{table}
\begin{tabularx}{\textwidth}{lQp{2.8cm}}
\midrule
\textbf{Nagano (2016)} & \textbf{In other hierarchies} & \textbf{Example} \\
\midrule

Shape & \citet{sproat:1991:the}: \textit{Shape} (QA)\newline
 \citet{scott:2002:stacked, cinque:2010:adj, laenzlinger:2005:french, laenzlinger:2011:elements}: \textit{Shape} (QA)& \textit{sankaku-no heya}\newline `triangular room'\\
& & \\
\tablevspace
Style/Pattern & no direct equivalent & \textit{shima-no uwagi} \newline `striped jacket' \\
\tablevspace
Type & potentially \citet{downing:1977:creation}: \textit{comparison} (RA)
 \newline potentially \citet{warren:1984:classifying}: \textit{comparant-compared} (RA)& \textit{Yōroppa-no keizai}  `European economy' (= European type of economy) \\
\tablevspace
Tendency towards & no direct equivalent & \textit{ame-no kisetsu}\newline  `rainy season' \\
\tablevspace
Line/Family/Group & close to \textit{group adjectives} \citep{grimshaw:1991:argument, cetnarowska:2015:categorial} & \textit{Surabu-no gengo}\newline  `Slavic languages' \\
\tablevspace
Location & \citet{downing:1977:creation}: \textit{place} (RA)\newline
    \citet{levi:1978:syntax}: \textit{in} (RA)  & \textit{umi-no seibutsu}\newline  `marine life'\\
& \citet{warren:1984:classifying}: \textit{place-object} (RA) & \\
& \citet{rae:2010:ordering, bortolotto:2016:RA}: \textit{Location} (RA) & \\
& \citet{takamine:2010:hierarchy}: \textit{Location} (PPs) & \\
\tablevspace
Time & \citet{downing:1977:creation}: \textit{time} (RA)\newline  \citet{levi:1978:syntax}: \textit{in} (RA) & \textit{yoru-no hōmon} \newline `nocturnal visit'\\
& \citet{warren:1984:classifying}: \textit{time-object} (RA) & \\
& \citet{rae:2010:ordering, bortolotto:2016:RA}: \textit{Time} (RA) & \\
& \citet{takamine:2010:hierarchy}: \textit{Time} (PPs) & \\

\tablevspace
Topic & \citet{levi:1978:syntax}: \textit{about} (RA)\newline  \citet{rae:2010:ordering, bortolotto:2016:RA}: \textit{Matter} (RA)  & \textit{keitai-no mondai} `morphological problem(s)' \\
\midrule
\end{tabularx}
\end{table}
\clearpage
\begin{table}[t]
\begin{tabularx}{\textwidth}{lQp{4cm}}
\midrule
\textbf{Nagano (2016)} & \textbf{In other hierarchies} & \textbf{Example} \\
\midrule
Status & \citet{downing:1977:creation}: \textit{occupation} (RA)\newline \citet{warren:1984:classifying}: \textit{affected object-actor} (RA)  & \textit{fuku-daitōryō-no ninki}\newline `vice-presidential term'\\
& included in \textit{group adjectives} (\citealt{grimshaw:1991:argument, cetnarowska:2015:categorial}) & \\
\tablevspace
Possession & potentially \textit{part-whole} in \citet{warren:1984:classifying}\newline \citet{nikolaeva:2012:possession, nikolaeva:2020:mixed, spencer:2017:denominal}: non-canonical modification carried out by nouns or \is{relational modifier}relational adjectives & \textit{yotsu-ashi-no dōbutsu}\newline `four-legged animal'\\
\lspbottomrule
\end{tabularx}
\end{table}

It would take us too far afield to discuss each of these \is{relational modifier}categories above. Instead, I will concentrate on a sample of relevant categories, most importantly \textsc{nationality}, \textsc{material} and \textsc{shape}. The reason for this is that these categories also appear in well-known hierarchies for \textit{qualitative} adjectives. It will be shown that modifiers of \textsc{nationality} and \textsc{material} have a strong direct character in Japanese. Modifiers of \textsc{shape}, on the other hand, do not. In the second part of this section, I will put forth data concerning combination with the different forms of the copula and discuss ordering and possible hierarchies.

\subsubsection{NATIONALITY and ethnic modifiers}

Modifiers of \textsc{nationality} are quite intriguing in Japanese. This category is quite low in ordering hierarchies of qualitative adjectives in the literature (cf. \citealt{scott:2002:stacked}). Recall, on the other hand, the group of \is{ethnic modifier}\textit{ethnic adjectives}, which has been put forth for various languages as toponymic relational adjectives denoting an origin (\citealt{bartning:1984:aspects, bosque:1996:postnominal, gunkel:2008:constraints, alex:2011:ea} a. o.). It is unclear whether the term \textit{ethnic adjectives} is applicable only with regard to thematic, or also to classificatory adjectives denoting an origin. Recall that the latter type can be used predicatively in \is{English} English \ref{bedsheet}, while \textsc{nationality} adjectives with a thematic reading, those modifying event nouns, cannot \ref{attack} (cf. \citealt{alex:2011:ea}).

\ex. \a. These bed sheets are Italian. \label{bedsheet}
\b. *This attack was Italian. \label{attack}

In the literature on Japanese, only equivalents of \textsc{nationality} modifiers with object nouns are given \citep{morita:2011:three, nagano:2016:are}, for example \ref{nihon-no kuruma}.\footnote{Note that in this example, the modifier \textit{Nihon-no} does not necessarily denote an origin, or nationality, but might denote several other relations to the head noun, including \textsc{location} on a reading `in Japan’.}

\exg. Nihon-no kuruma\\
Japan-\textsc{mod} car\\
`Japanese car' \label{nihon-no kuruma}

Similar to thematic \is{relational modifier}relational modifiers, the equivalents of thematic \is{ethnic modifier}ethnic adjectives have probably not been discussed due to the fact that they just seem to highlight a case of modification by a noun. Arguably, an example could be the following, given in passing by \citet[48]{nagano:2016:are}. 

\exg. Nihon-no kōgeki\\
Japan-\textsc{mod} attack\\
`Japanese attack'

\largerpage
It is a crucial observation that \textsc{nationality} modifiers in \is{ethnic modifier}Japanese, also those with a classificatory meaning, are only used attributively, never predicatively \ref{kurumanihonda} and are non-gradable \ref{totemonihon}, \ref{yorinihon}. Importantly, gradability is not even possible with endpoint modifiers as shown in \ref{kanzenni}. This crucially differentiates them from \is{English} English classificatory adjectives.

\ex. \ag. Nihon-no kuruma\\
Japan-\textsc{mod} car\\
`Japanese car'
\bg. *Kono kuruma-wa Nihon-da.\\
this car-\textsc{top} Japan-\textsc{cop}\\
(`This car is Japanese.’) \label{kurumanihonda}
\cg. *totemo Nihon-no kuruma\\
very Japan-\textsc{mod} car\\
(`very Japanese car’) \label{totemonihon}
\dg. *Watashi-no kuruma-wa sensei-no kuruma-yori Nihon (no kuruma) da.\\
I-\textsc{mod} car-\textsc{top} teacher-\textsc{mod} car-than Japan \textsc{mod} car \textsc{cop}\\
(`My car is more Japanese than my teacher's (car).’) \label{yorinihon}
\eg. *kanzenni Nihon-no kuruma\\
completely Japan-\textsc{mod} car\\
(`completely Japanese car’) \label{kanzenni}

What holds for \is{relational modifier}classificatory, also holds for the equivalents of thematic \is{ethnic modifier}ethnic modifiers as shown below.

\ex. \ag. Nihon-no kōgeki\\
Japan-\textsc{mod} attack\\
`Japanese attack'
\bg. *Kono kōgeki-wa Nihon-da.\\
this attack-\textsc{top} Japan-\textsc{cop}\\
(`This attack is Japanese.’)
\cg. *totemo / kanzenni Nihon-no kōgeki\\
very / completely Japan-\textsc{mod} attack\\
(`very / completely Japanese attack’)

While in \is{English} English, \textsc{theme} readings of ethnic adjectives have been argued to be impossible \citep{kayne:1981:ecp, giorgi:1991:syntax},\footnote{Nevertheless such readings have been marginally reported in English \citep{arseni:2014:EA} and seem productive in Polish \citep{cetnarowska:2015:categorial}.

\par

\ex. the French defeat \null\hfill \citep[23]{arseni:2014:EA}

\exg. chiń-s-ka pora\v{z}ka (Polish)\\
China-\textsc{ra-nom.sg.f} destruction.\textsc{nom.sg.f}\\
`destruction of China' \null\hfill \citep[136]{cetnarowska:2015:categorial}

\par


\citet[20]{rainer:2013:can} mentions that online he found ``dozens of examples of the type \textit{the Russian invasion by Napoleon, the Afghan invasion by Russia, the Polish invasion by Germany, the Chinese invasion by Japan}, and the like, many of which apparently occur in authoritative documents".} no such restrictions are applicable to the Japanese equivalents, presumably for the fact that they are simply nouns. \is{Polish}

\largerpage[2]
Furthermore, object and subject roles can both be expressed inside the same sentence as in \ref{sub-ob} (cf. \citealt{murasugi:1991:noun}). In such cases, the first modifier is usually interpreted as the subject and the second, the one adjacent to the noun, as the object. As expected, Japanese exhibits a rather free behavior of word order here as well and the reverse interpretations are not \is{relational modifier}impossible.\footnote{Many other thematic roles can be expressed within the DP, but for most, an additional postposition needs to appear before \textsc{mod}.}

\exg. Amerika-no Itaria-no hakai \label{sub-ob}\\
America-\textsc{mod} Italy-\textsc{mod} destruction\\
`destruction of Italy by America’ \\ \null \hfill \textit{Amerika}=subject, \textit{Itaria}=object

Finally, modifiers of \textsc{nationality} can modify human nouns as the following example given in \citet[91]{morita:2011:three} shows. Again, predicative use is impossible, and so is gradability.

\ex. \ag. Nihon-no haiyū\\
Japan-\textsc{mod} actor\\
`Japanese actor'
\bg. *Kono haiyū-wa Nihon-da.\\
this actor-\textsc{top} Japan-\textsc{cop}\\
(`This actor is Japanese.’)
\cg. *totemo Nihon-no haiyū\\
very Japan-\textsc{mod} actor\\
(`very Japanese actor’)
\dg. *Kono haiyū-wa sono haiyū-yori Nihon-no haiyū da.\\
this actor-\textsc{top} that actor-than Japan-\textsc{mod} actor \textsc{cop}\\
(`This actor is more Japanese than that actor.’)
\eg. *kanzenni Nihon-no haiyū\\
completely Japan-\textsc{mod} actor\\
(`completely Japanese actor’)

This might be different for other languages. \citet{arseni:2014:EA} have argued that all \is{ethnic modifier}ethnic modifiers have an underlying origin-reading which arguably permits predicative use, as in the following example in \is{English} English.

\ex. Guillem is French. \hfill \citep[23]{arseni:2014:EA}

In \is{German} German, as in Japanese, such examples are clearly marked as also mentioned in \citet{kotowski:2016:adj}. %\is{German}

\exg. ??Guillem ist französisch. (German)\\
Guillem is French\\
`Guillem is French’

In order to express the \is{ethnic modifier}nationality of a human entity in predicative position in Japanese, the suffix -\textit{jin} \textjpn{人}, Lit. `human' needs to be added. Again, the same holds true in German \ref{franzose}. \is{relational modifier}

\exg. Kono gekai-wa Nihon-*(jin)-da. (Japanese)\\
this surgeon-\textsc{top} Japan-person-\textsc{cop}\\
`This surgeon is Japanese.’ (= a Japanese person)

\exg. Guillem ist Franzose. (German)\\
Guillem is Frenchman\\
`Guillem is French’ (= a French person) \label{franzose}

A remaining question is how the thematic relation, the theta roles within DPs such as \ref{sub-ob} are being established. In the cartographic research literature, dedicated projections are given for modifiers that express a \textsc{theme} or \textsc{agent} reading \citep{cinque:1993:evidence, laenzlinger:2005:french, laenzlinger:2011:elements, bortolotto:2016:RA}. As for their derivation, \citet{laenzlinger:2011:elements} predicts that modifiers in the specifier positions of these projections have moved there from lower positions. Taking the parallel to the verbal domain seriously, nouns with a \textsc{theme} reading start out as complements of N, those with an \textsc{agent} reading as specifiers of an additional nP projection (see also \citealt{adger:2003:core}). One derivation could look as in \figref{Structure of Amerika-no Itaria-no hakai}. Recall that the first modifier is typically interpreted as the agent, but that the reading can also be reversed, which underlines the unspecific nature of the functional projections within the direct domain.

 \begin{figure}
\centering
\begin{tikzpicture}[baseline=0pt, remember picture]
\tikzset{level distance=25pt, sibling distance=5pt}
\Tree [.DP\textsubscript{internal} {} [.D′ [.FP [.NP\textsubscript{1} \edge[roof]; {Amerika} ] [.F′ [.FP [.NP\textsubscript{2} \edge[roof]; {Itaria} ] [.F′ [.nP $<$NP\textsubscript{1}$>$ [.n′ [.NP {} [.N′ $<$NP\textsubscript{2}$>$ [.N hakai ] ] ] n ] ] F ] ] F ] ] D ] ]
\end{tikzpicture}
\caption{Structure of \textit{Amerika-no Itaria-no hakai} (= \ref{sub-ob})}
\label{Structure of Amerika-no Itaria-no hakai}
\end{figure}

To summarize, modifiers of \textsc{nationality} are a productive source of direct modification in Japanese and the relevant lexemes clearly have a nominal status.

\subsubsection{Modifiers of MATERIAL}

\textsc{material} is in Japanese also exclusively expressed via \textit{no}-modifiers. These cannot be used predicatively \ref{komugida} and are not gradable \ref{totemokomugi}, as judged unequivocally so by all speakers I consulted in relevant constructions. One example is \textit{komugi-no} `wheat(en)’, and recall also the examples with \textit{chīku-no} `teak’ from Section \ref{japanesedp}. \is{relational modifier}

\ex. \ag. komugi-no pan\\
wheat-\textsc{mod} bread\\
`wheat(en) bread'
\bg. *Kono pan-wa komugi-da.\\
this bread-\textsc{top} wheat(en)-\textsc{cop}\\
(`This bread is wheat(en).’) \label{komugida}
\cg. *totemo komugi-no pan\\
very wheat-\textsc{mod} bread\\
(`very wheat(en) bread’) \label{totemokomugi}

In contrast to modifiers of \textsc{nationality}, modifiers of \textsc{material} can be modified via endpoint modifiers \citep{kennedy:2005:scale, fabregas:2014:adjectival, oshima:2019:gradability}. Cf. Sections \ref{RA} and \ref{gradability}. \is{relational modifier}

\exg. kanzenni komugi-no pan\\
completely wheat(en)-\textsc{mod} bread\\
`completely wheaten bread’ (= bread consisting completely of wheat)

Finally, note that also modifiers of \textsc{material} are nominals and occur productively with case particles as shown below.

\exg. Ano komugi-wa oishi-i.\\
that wheat-\textsc{top} tasty-\textsc{cop}\\
`That wheat is tasty.’

Restrictions in predication and gradability are the reason why modifiers of \textsc{material} are closer to relational than to qualitative modifiers \citep{nagano:2016:are, morita:2011:three}. Their syntactic restrictions qualify them as direct modifiers.

\subsubsection{Modifiers of SHAPE are not direct-only}

The category \textsc{shape} is not a part of any study on \is{relational modifier}relational adjectives in European languages, but is rather a central category of qualitative adjectives in \citet{sproat:1991:the, cinque:1993:evidence, cinque:2010:adj}, \citet{scott:2002:stacked} and \citet{laenzlinger:2005:french, laenzlinger:2011:elements}. Why is it treated then as a relational modifier in \citet{nagano:2016:are} and also \citet{morita:2010:internal, morita:2011:three}? First, note that the category \textsc{shape} is expressed in Japanese via \textit{i}- and \textit{na}-adjectives, but also via \textit{no}-modifiers and that there are lexemes, such as \textit{maru} `round' and \textit{shikaku} `square’, which can appear with either \textsc{i} or \textsc{no} in attributive position.\footnote{The \is{BCCWJ}BCCWJ yields 405 results with -\textsc{i} and 74 with -\textsc{no} for \textit{shikaku}, and even 33 with -\textsc{na}. It is a matter of debate whether we should predict to different spell-out forms of \is{MOD@\textsc{mod}}\textsc{mod} or simply two different lexemes which differ in word class status and thus lead to a different form of the syntactic head. The latter option seems more attractive in the system presented here, but, at the moment, nothing hinges on this.}

\ex. \ag. maru-i / maru-no tēburu\\
round-\textsc{mod} / round-\textsc{mod} table\\
`round table'
\bg. shikaku-i / shikaku-no heya\\
square-\textsc{mod} / square-\textsc{mod} room\\
`square room'

The authors predict that both \is{gradability}gradability and predicativity are marked for the \textit{no}-variant of these lexemes, making them essentially equivalents of true relational adjectives. However, not only are there examples of predicative use in the corpus,\footnote{12 occurrences for \textit{shikaku-da}, which is still very little making it a prototypical direct modifier, but 88 for \textit{maru-da}.} predicative use is also accepted by many speakers, including the ones I asked, as shown below.

\ex. \ag. Kono heya-wa shikaku-da.\\
this room-\textsc{top} square-\textsc{cop}\\
`This room is square.’
\bg. Kono tēburu-wa maru-da.\\
this table-\textsc{top} round-\textsc{cop}\\
`This table is round.’

Regarding \is{gradability}gradability, the \textit{no}-modifier-variant is marked in combination with the degree adverb \textit{totemo} `very’, compared to the \textit{i}-variant, showing that this is not for semantic reasons. None of my consultants accepted the combination of the degree adverb with \textit{maru-no} \ref{totemomaru}, and only two out of five judged \textit{totemo shikaku-no} `very square' to be marginally accep	abref{totemoshikaku}.

\ex. \ag. *totemo maru-no tēburu\\
very round-\textsc{mod} table\\
`very round table' \label{totemomaru}
\bg. ??totemo shikaku-no heya\\
very square-\textsc{mod} room\\
`very square room' \label{totemoshikaku}

However, the following combinations involving gradability are perfectly fine.

\ex. \ag. motto shikaku-no heya\\
more square-\textsc{mod} room\\
`more square room’
\bg. omot-ta-yori maru-no tēburu\\
think-\textsc{pst}-than round-\textsc{mod} table\\
`a table (which is) rounder than X thought’

\largerpage
Indeed, there are some prototypical direct \is{relational modifier}\textit{no}-modifiers in the data set that arguably express shape, such as \textit{hishigata-no} `diamond-shaped, rhombic' (37 attr., 4 pred.; cf. also \citealt[98]{morita:2011:three}) and \textit{yūkei-no} `tangible' (29 attr., 0 pred.), suggesting that direct \textsc{shape} modifiers do exist. However, importantly this is not the case for all modifiers of this category, unlike \textsc{nationality} or \textsc{material}.

This discussion extends to modifiers of the category \textsc{color}: The four basic color terms \textit{aka} `red', \textit{ao} `blue', \textit{shiro} `white' and \textit{kuro} `black' can also appear with -\textsc{i} and -\textsc{no}, that is as adjectives or \textit{no}-modifiers.\footnote{A complementary distribution pointed out in the literature \citep{sawada:1992:meishi, kishita:2005:shikisai} is that the \textit{no}-variant, the color noun \citep{okami:2012:two, okami:2016:on}, serves the purpose of identification and must receive a ``contrastive interpretation" \citep{koopman:2016:further} thus is infelicitous in contexts where the referent does not need to be identified.} Again, both versions are possible \is{gradability}predicatively, contrary to what is argued in \citet{morita:2011:three}, but the \textit{no}-variant is only marginally gradable \ref{totemoshiro}.\footnote{Note that direct-only modifiers of \textsc{color} do appear in the data set to be discussed in Section \ref{QM}.}

\ex. \ag. shiro-i / shiro-no kōto\\
white-\textsc{mod} / white-\textsc{mod} coat\\
`white coat' \null \hfill \citep{koopman:2016:further}
\bg. Kono kōto-wa shiro-da.\\
this coat-\textsc{top} white-\textsc{cop}\\
`This coat is white.’
\cg. *totemo shiro-no kōto\\
very white-\textsc{mod} coat\\
(`very white coat’) \label{totemoshiro}

\subsubsection{Other copula forms}

As outlined above, the absence of predicative use correlates with the impossible co-occurrence with other copula forms, including \textsc{de} used for coordination. This seems to be borne out for non-predicative \is{relational modifier}relational modifiers in Japanese as exemplified below. Crucially, as noted in Section \ref{coordination copula}, these modifiers can neither appear as the second conjunct in coordinative structures \ref{kirei second}.\footnote{One apparent counterexample was the material modifier \textit{komugi} `wheat(en)' where /de could be understood as a postposition, denoting the material of origin so \textit{komugi-de kata-i pan} could have the interpretation `hard bread made from wheat'.}  \is{coordinative copula}

\exg. *Chūgoku-de kirei-na kabin\\
China-\textsc{pred} beautiful-\textsc{mod} vase\\
(`a Chinese and beautiful vase’)

\exg. *kirei-de Chūgoku-no kabin\\
beautiful-\textsc{pred} China-\textsc{mod} vase\\
(`a beautiful and Chinese vase’) \label{kirei second}

These data points highlight that \textsc{de} has to appear in phrasal contexts and always involves a predication structure (\citealt{bowers:1993:predication}, see Section \ref{copulaelements}). Even if only one modifier involved in the coordination structure is banned from predication, coordination becomes ungrammatical.

Furthermore, as exemplified in \ref{komugi dearu}, in no instance was the canonical copula \textsc{de-ar-u} allowed in combination with non-predicative relational modifiers. This again speaks against analyses that treat \textsc{no} as some predicative element.

\exg. *komugi-de-ar-u pan\\
wheat-\textsc{pred-cop-prs} bread\\
(`bread which is wheat(en)’) \label{komugi dearu}

\subsubsection{Hierarchy and ordering of relational modifiers}

\citet{bortolotto:2016:RA} establishes an \is{adjective ordering restrictions}ordering hierarchy for \is{relational modifier}relational modifiers based on Romance, given here in terms of structural selection.

\ex. \emph{Hierarchy of relational adjectives (based on \citealt[180]{bortolotto:2016:RA}})\\
\textsc{time} $>$ \textsc{location} $>$ \textsc{agent} $>$ \textsc{instrument} $>$ \textsc{matter} ($>$) \textsc{theme} $>$ N

Also recall that it is a characteristic of relational adjectives to be strictly adjacent to their head noun and not to allow permutation with qualitative adjectives, as shown below for \is{English} English.

\ex. \a. big wooden table
\b. *wooden big table \null \hfill \citep[66]{bisetto:2010:RA}

In the DP structure presented in this book, no specific functional projections dedicated to the different categories of relational modifiers are included inside the direct domain. This predicts (correctly) that relational modifiers in Japanese are not ordered according to a specific hierarchy. However, \citet{watanabe:2012:direct} argues in favor of such a hierarchy. First, he mentions that modifiers of \textsc{nationality} and \textsc{material} have to be closer to their head noun than modifiers of other categories.\footnote{I give his judgments below, but it is not clear whether this is an issue of grammaticality or preference. Note also that he uses as contrast the non-predicative modifier \textit{chiisa-na} `small' (see Section \ref{rentaishi}).}

\ex. \textsc{dimension} $\succ$ \textsc{material}
\label{chiisa1} \ag. chiisa-na ki-no hashi\\
small-\textsc{mod} wood-\textsc{mod} bridge\\
\bg. ??ki-no chiisa-na hashi\\
wood-\textsc{mod} small-\textsc{mod} bridge\\
`small wooden bridge' \hfill \citep[507]{watanabe:2012:direct}

\ex. \textsc{dimension} $\succ$ \textsc{nationality}
\label{chiisa2} \ag. chiisa-na Chūgoku-no kabin\\
small-\textsc{mod} China-\textsc{mod} vase\\
\bg. ??Chūgoku-no chiisa-na kabin\\
China-\textsc{mod} small-\textsc{mod} vase\\
`small Chinese vase' \hfill \citep[507]{watanabe:2012:direct}

Next, the author argues that the categories \textsc{nationality} and \textsc{material} reflect the structural \is{adjective ordering restrictions}hierarchy \textsc{nationality} $>$ \textsc{material} in linear order when occurring together.

\largerpage
\ex. \textsc{nationality} $\succ$ \textsc{material}
\label{hokuo} \ag. hokuō-no ki-no isu \\
North.Europe-\textsc{mod} wood-\textsc{mod} chair\\
\bg. *ki-no hokuō-no isu\\
wood-\textsc{mod} North.Europe-\textsc{mod} chair\\
`North European wooden chair' \hfill \citep[508]{watanabe:2012:direct}

\ex. \textsc{nationality} $\succ$ \textsc{material}
\label{Chiri} \ag. Chiri-no kin-no kubikazari\\
Chile-\textsc{mod} gold-\textsc{mod} necklace\\
\bg. *kin-no Chiri-no kubikazari\\
gold-\textsc{mod} Chile-\textsc{mod} necklace\\
`Chilean gold necklace' \hfill \citep[508]{watanabe:2012:direct}

However, first, this assessment is not shared by all \is{absence of ordering restrictions}authors. \citet[Footnote 14]{nagano:2016:are} emphasizes that she shares the judgment for \ref{chiisa1} and also for \ref{hokuo}, but not for \ref{chiisa2} and \ref{Chiri}, meaning she does not have a preferred \is{relational modifier}order. She suggests that the ordering restrictions in the other cases might be phonologically conditioned \citep[Footnote 14]{nagano:2016:are}. More precisely, speakers including Nagano might prefer a shorter modifier to the right of a \is{mora length}longer one. This is the case for \textit{chiisa} `small' (3 morae) and \textit{hokuō} `North European' (4 morae) preceding \textit{ki} `wood(en)' (1 mora). It is not however for \textit{chiisa} `small' vs. \textit{Chūgoku} `China, Chinese' (3 vs. 4 morae) and \textit{Chiri} `Chile, Chilean' vs. \textit{kin} `gold(en)' (both 2 morae).

Furthermore, \textit{ki} `wood(en)’ is the only monomoraic modifier in the data set and might be the only one that triggers such ordering restrictions (cf. Section \ref{optionality}). Concerning \ref{Chiri}, Ken Hiraiwa (p.c.) mentioned to me that some speakers might think of \textit{kin-no kubikazari} `gold(en) necklace' as a set phrase which is why separating the two entities is judged as not good.

I have tested for several ordering preferences among relational modifiers and in combinations with qualitative modifiers and could not detect any ordering restrictions. Importantly, in order to exclude the factors \textit{phonology} (the number of morae) and \textit{word class status} (\textit{i}-, \textit{na}-adjective, \textit{no}-modifier), I tried to verify ordering restrictions for different combinations and scenarios (see Section \ref{optionality}), however, in no case did I trigger an increase or decrease in preference for any speaker in any pair. In contrast, note the awkward English translation for almost all b) examples in the examples given below.

\ex. \textsc{nationality} vs. \textsc{material}
\ag. Denmāku-no chīku-no isu\\
Denmark-\textsc{mod} teak-\textsc{mod} chair\\
`Danish teak chair'
\bg. chīku-no Denmāku-no isu\\
teak-\textsc{mod} Denmark-\textsc{mod} chair\\
`teak Danish chair'

\newpage
\ex. \textsc{quality} vs. \textsc{nationality}
\ag. kirei-na Chūgoku-no kabin\\
beautiful-\textsc{mod} China-\textsc{mod} vase\\
`beautiful Chinese vase'
\bg. Chūgoku-no kirei-na kabin\\
China-\textsc{mod} beautiful-\textsc{mod} vase\\
`Chinese beautiful vase'

\ex. \textsc{material} vs. \textsc{quality} (\textsc{shape})
\ag. komugi-no kata-i pan\\
wheat-\textsc{mod} hard-\textsc{mod} bread\\
`hard wheat(en) bread'
\bg. kata-i komugi-no pan\\
hard-\textsc{mod} wheat-\textsc{mod} bread\\
`wheat(en) hard bread'

\ex. \textsc{color} vs. \textsc{nationality}
\ag. ao-i Chūgoku-no kabin\\
blue-\textsc{mod} China-\textsc{mod} vase\\
`blue Chinese vase'
\bg. Chūgoku-no ao-i kabin\\
China-\textsc{mod} blue-\textsc{mod} vase\\
`Chinese blue vase'

Note also the following \is{relational modifier}contrast given in \citet{nagano:2016:are}, which is identical to \is{absence of ordering restrictions}Watanabe’s example \ref{chiisa2} with the difference that the \textit{i}-adjective \textit{chiisa-i} `small’ is used instead of the \is{rentaishi@\textit{rentaishi}}\textit{rentaishi} \textit{chiisa-na} `small’.

\ex. \textsc{dimension} vs. \textsc{nationality}
\ag. chiisa-i Chūgoku-no kabin\\
small-\textsc{mod} China-\textsc{mod} vase\\
`small Chinese vase'
\bg. Chūgoku-no chiisa-i kabin\\
China-\textsc{mod} small-\textsc{mod} vase\\
`Chinese small vase' \null \hfill \citep[57]{nagano:2016:are}

These facts emphasize that neither \is{absence of ordering restrictions}\textit{ordering restrictions} nor \textit{adjacency to the head noun} are valid criteria for the equivalents of relational modifiers in Japanese, which is to be expected considering the broad direct domain of Japanese.

\subsubsection{Further issues: Nominal status and predicative use} \label{Ra japanese predicative}

There are two remaining issues concerning \is{relational modifier}relational modifiers. First, since relational modifiers are rather nominal, they make up a prototype of direct nominals in Japanese as also mentioned in \citet{watanabe:2012:direct} for modifiers of \textsc{nationality} and \textsc{material}. This is shown below via productive combination with case particles \ref{chugoku kawatta} and the inability to be \is{nominalization}nominalized \ref{chugoku sa}.

\exg. Chūgoku-ga kawat-ta.\\
China-\textsc{nom} change-\textsc{pst}\\
`China has changed.’ \label{chugoku kawatta}

\exg. *Chūgoku-sa\\
China-ness \label{chugoku sa}\\

A further piece of evidence for nominal status is the fact that these modifiers can be modified themselves leading to \is{bracketing ambiguity}bracketing ambiguities. As shown in \ref{katai komugi}, the \textit{i}-adjective \textit{kata-i} `hard’ can either modify \textit{pan} `bread’, skipping over \textit{komugi} `wheat(en)’ which modifies \textit{pan} `bread’. On an alternative reading, \textit{kata-i} modifies \textit{komugi} `wheat(en)’ yielding the DP \textit{kata-i komugi} `hard wheat’ which in turn modifies \textit{pan}.\footnote{Such bracketing ambiguities can be compared to cases of what \citet{nikolaeva:2020:mixed} call \textit{syntagmatic mixing}. In languages such as Selkup and other Samoyedic languages, relational adjectives, albeit being subject to derivational operations, can still be modified. In the Japanese case however, this seems closer to what \citet{spencer:2005:towards} calls \textit{indeterminacy}.} \is{nominalization}

\exg. kata-i komugi-no pan\\
hard-\textsc{mod} wheat-\textsc{mod} bread\\
1. $\approx$ `hard wheat(en) bread’\\
2. $\approx$ `bread out of hard wheat' \label{katai komugi}

Another cross-linguistic phenomenon concerns predicative use. Recall that \is{relational modifier}relational adjectives in various languages can be used predicatively, either as a result of a class shift to a qualitative adjective, or when they have a \is{kind-reading}\textit{kind-reading} \citep{mcnally:2004:relational}. Another mechanism was noun ellipsis, facilitated in contrastive statements. In Japanese, as briefly mentioned in Section \ref{test predication}, relational modifiers also become possible in predicative position when they are connected to a certain kind of suffix, referred to as \is{bound marker}\textit{bound marker} or \textit{relational noun} in \citet{nagano:2016:are, shimada:2018:relational}. When these markers, which roughly correspond in their meaning to relational modifiers, are added to a relational modifier, a new entity called \textit{expanded modifier} in \citet[54]{nagano:2016:are} is formed. These expanded modifiers are not only more specified in their meaning, they are also productive in predicative position. This is illustrated below where the pair \textit{chīzu} `cheese’ and \textit{iri} `put inside’ form \textit{chīzu-iri} `(with) cheese put inside’. As \citet{levi:1978:syntax} has argued for similar cases in \is{English} English, establishing modifier and modified noun as a set facilitates predicative use. Interested readers are referred to Appendix A for an overview. \is{Selkup}

\ex. \ag. chīzu-iri-no pan\\
cheese-put.in-\textsc{mod} bread\\
`bread with cheese put inside'
\bg. kono pan-wa chīzu-*(iri)-da\\
this bread-\textsc{top} cheese-put.in-\textsc{cop} \\
`This bread is with cheese inside.’

Nouns to which the adjective-forming suffix -\textit{teki} is added also become possible in predicative position and gradable. As mentioned at the beginning of this chapter, this suffix was argued originally by \citet{bisetto:2010:RA} to form relational adjectives in Japanese, but the examples below, illustrated with \textsc{nationality} adjectives with a \textsc{theme} reading, clearly show that a class shift to a qualitative adjective occurred (see also \citealt[48]{nagano:2016:are}).

\ex. \label{like} \ag. Ano kōgeki-wa America-*(teki)-d-at-ta.\\
this attack-\textsc{top} America-\textsc{adj}-\textsc{cop-pred-pst}\\
`This attack was (typically) American.'
\bg. totemo America-*(teki)-na kōgeki\\
very America-\textsc{adj}-na attack\\
`very American (= typical for Americans) attack'

\subsubsection{Summary: Characteristics of relational modifiers in Japanese}

Taking all these facts into account, Japanese does exhibit a category that is equivalent to the cross-linguistic category of \is{relational modifier}relational adjectives. Due to their unclear lexical status, they are referred to as \textit{relational modifiers}. They show the following characteristics:

\ex. \a. They generally cannot be used predicatively.
\b. They are not gradable.
\c. They cannot be coordinated.
\d. They can be used as real nouns.
\e. However, they do not show the typical ordering restrictions and strict adjacency to the head noun associated with relational adjectives in other languages.

\subsection{Direct qualitative modifiers} \label{QM}

Those authors who denied the existence of direct modification in Japanese \citep{sproat:1991:the, laenzlinger:2011:elements} only investigated modifiers of categories such as \textsc{dimension}, \textsc{shape} and \textsc{quality}, those categories which are related to \textit{qualitative} modification in European languages, where \textit{qualitative} is used contra \textit{relational}. In a sense, those categories form the most basic adjectival categories. The fact that no direct modifiers seem to be available for these categories in Japanese was thus directly related to the alleged \is{absence of ordering restrictions}absence of direct modification. Importantly, as also mentioned in the introduction, most modifiers of these categories are not direct-only in other languages either, and the status of a certain qualitative adjective in English only become obvious in contexts of ordering \citep{cinque:2010:adj, laenzlinger:2011:elements}.

However, Japanese exhibits several lexemes that one could classify as direct-only qualitative modifiers and these cases are worth discussing.\footnote{I have annotated (some) modifiers for their semantic category, but, of course, this is sometimes quite subjective. In any case, since the relevant functional projections are not active in Japanese, the semantic category of a modifier is not crucial.} Some potentially direct modifiers denoting \textsc{shape} were discussed above (\textit{hishigata-no} `diamond-shaped, rhombic’; \textit{yūkei-no} `tangible’). Examples of direct modifiers of the category \textsc{dimension} are the lexemes \textit{ōki} `big' and \textit{chiisa} `small’. Those can appear either with \textsc{i} or with \textsc{na}, without a substantial change in meaning. However, they never appear with the copula \textsc{da} \ref{ieokida}, only the \textit{i}-variant can be used predicatively.

\ex. \ag. ōki-i / chiisa-i ie\\
big-\textsc{mod} / small-\textsc{mod} house\\
`big / small house'
\bg. Kono ie-wa ōki-i / chiisa-i.\\
this house-\textsc{top} big-\textsc{cop} / small-\textsc{cop}\\
`This house is big / small.’
\cg. ōki-na / chiisa-na ie\\
big-\textsc{mod} / small-\textsc{mod} house\\
`big / small house'
\dg. *Kono ie-wa ōki-da / chiisa-da.\\
this house-\textsc{top} big-\textsc{cop} / small-\textsc{cop}\\
(`This house is big / small.’) \label{ieokida}

There is also no use as adverbials or inchoatives attested for the \textit{na}-version of these lexemes, contrary to the corresponding \textit{i}-adjective, where \textsc{i} inflects to \textsc{ku}. 

\exg. Ōki-\textbf{ku} / *ōki-\textbf{ni} nat-ta.\\
big-\textsc{adv} / big-\textsc{adv} become-\textsc{pst}\\
`Became big.’

Therefore, the \textit{na}-variants of these lexemes are also classified as \is{rentaishi@\textit{rentaishi}}\textit{rentaishi}, the topic of Section \ref{rentaishi}. An additional factor to discuss is that they are in fact gradable as illustrated below. Remember, however, that \is{gradability}non-gradability is an accompanying factor of non-predicative modifiers, not a prerequisite \citep[47]{cinque:2010:adj}. 

\exg. Ichi-ban ōki-na no-wo eran-da.\\
number-one big-\textsc{mod} \textsc{ln}-\textsc{acc} choose-\textsc{pst}\\
`X chose the biggest one.’ \hfill \citep[1: 28]{sekine:1937:fukutaishi}

Potential other direct modifiers of \textsc{dimension} are documented in the data set, for instance \textit{ōgata-no} `large-scale' (709 attr., 73 pred.) and \textit{atsude-no} `thick’ (189 attr., 2 pred.). However, both are, in principle, possible in predicative position, albeit only marginally attested. This holds especially for \textit{atsude-no} whose possibility in predicative position was corroborated by some speakers despite its very few occurences.

Clearer cases of direct modifiers are found among the category \textsc{quality}, which is \textit{per se} one of the more, perhaps the most, versatile category. For instance, \textit{gokujō-no} `of highest quality' (142 attr. 5 pred.) or \textit{seirai-no} `innate’ (92 attr., 0 pred.) yield few to none predicative occurences.

Direct modifiers of the category \textsc{color} are \textit{shikkoku-no} `pitch black' (165 attr., 2 pred.) and \textit{konpeki-no} `deep blue’ (46 attr., 0 pred.).

\largerpage
It is interesting to note that direct modifiers of some less canonical categories given in the very fine-grained hierarchy of \citet[106, 114]{scott:2002:stacked}, can be found as well. Examples include \textit{gokkan-no} `freezing cold’ (44 attr., 5 pred.) representing the category \textsc{temparature} or \textit{kyūsoku-na} `fast, rapid’ (599 attr., 4 pred.) for the category \textsc{speed}, illustrated below for \textit{gokkan-no}.

\ex. \ag. gokkan-no kisetsu\\
freezing.cold-\textsc{mod} season\\
`freezing cold season'
\bg. ??Ima-no kisetsu-wa gokkan-da\\
now-\textsc{mod} season-\textsc{top} freezing.cold-\textsc{cop}\\
`The season at the moment is freezing cold.’

Finally, there are some direct-only modifiers that do not fit any semantic category present in established hierarchies, such as \textit{kidai-no} `rare’ (71 attr., 0 pred.) and \textit{kaishin-no} `spot-on; critical’ (58 attr., 0 pred.).

\ex. \ag. kaishin-no ronbun\\
spot.on-\textsc{mod} essay\\
`a spot-on paper; a paper written to one's full satisfaction'
\bg. *Kono ronbun-wa kaishin-da.\\
this paper-\textsc{top} spot.on-\textsc{cop}\\
(`This paper is spot-on.’)

\ex. \ag. kidai-no sainō\\
rare-\textsc{mod} talent\\
`rare talent’ \label{kidai}
\bg. *Kare-no sainō-wa kidai-da.\\
he-\textsc{mod} talent-\textsc{top} rare-\textsc{cop}\\
(`His talent is rare.’)

Since the functional projections inside the direct domain are not governed by semantic categories, the problem of where to sort to sort these modifiers, i.e, into which functional projections, disappears. They are merged in the direct domain, but in an unspecific functional projection, as depicted in \figref{Structure of kidai-no saino}.

\begin{figure}
\centering
\begin{tikzpicture}[baseline=0pt, remember picture]
\tikzset{level distance=22pt}
 \Tree [.DP\textsubscript{internal} {} [.D′\textsubscript{internal} [.FP\textsubscript{direct} [.ModP {} [.Mod′ [.XP \edge[roof]; kidai ] [.Mod no ] ]] [.F′\textsubscript{direct} [.NP \edge[roof]; sainō ] F ] ] D ] ]
\end{tikzpicture}
\caption{Structure of \ref{kidai}}
\label{Structure of kidai-no saino}
\end{figure}

\subsection{Idiomatic expressions} \label{idiomatic}

Another source for direct modification are \is{idiomatic modifier}modifiers involved in idiomatic expressions, although it might be better to speak of \textit{modifiers in non-free phrases}.\footnote{See \citet{melcuk:2012:phraseology} and \citet{kohlich:2024:nonfree} for the difference between idiomatic expressions and collocations.} Recall that in English, only prenominal adjectives can elicit an idiomatic reading with their corresponding head noun, as for example \textit{white} in \textit{white lie} \citep[87–88]{cinque:2010:adj}. The only idiomatic expression with a single modifier given in the literature is as follows, given in \citet[122]{nagano:2015:RA}. As expected, predicative use is impossible with the idiomatic meaning.

\ex. \ag. aka-no tan'in\\
red-\textsc{mod} stranger\\
`a total stranger' (Lit. `a red stranger')
\bg. *Kono tan'in-wa aka-da\\
this stranger-\textsc{top} red-\textsc{cop}\\
only meaning: `This stranger is red.’

Several more idiomatic modifiers in Japanese can be identified, where ``idiomatic" refers to modifiers that only co-occur with a certain noun, or a certain group of nouns with similar meaning. Thus, it is no prerequisite that a new meaning is formed, as in the example \textit{white lie}. One example of an idiomatic modifier, already mentioned in Section \ref{japanesedp}, is \textit{kōki-no} `curious’. This lexeme yields 69 results in attributive position, against zero with the copula. 61 out of those 69 occurrences are with nouns with the meaning `gaze, look’, among them \textit{me} (also `eye'), \textit{manazashi} `gaze' and \textit{shisen} `gaze, look'. The restriction in predicativity, and also gradability, is shown below.\footnote{Compared to \textit{me}, \textit{manazashi} and \textit{shisen} can be described as nouns with the same meaning, but from a higher register. Ken Hiraiwa has told me that to his native speaker ears the modification of \textit{manazashi} `gaze' via \textit{kōki-no} `curious’ seems to be truly idiomatic.} \is{idiomatic modifier}

\ex. \ag. Shikashi kare-ga kōki-no me-de mi-rare, ōku-no hito-ga kare-no hanashi-wo kiki-ta-gat-ta koto[...]\\
but he-\textsc{nom} curious-\textsc{mod} eye-\textsc{ins} see-\textsc{pass} many-\textsc{mod} people-\textsc{nom} he-\textsc{mod} story-\textsc{acc} hear-\textsc{vol}-seem-\textsc{pst} fact\\
`But the fact that he was viewed with curious eyes (a curious gaze) and many people wanted to hear his story [...]’\\ Ogawa, Ryō (2002): Doreishōnin Soniē - 18-seiki Furansu no doreikōki to Afurika-shakai. Tokyo: Yamakawa Shuppansha, via BCCWJ, 2023/11/09.
\bg. *Kare-ga mi-rare-ta me-wa kōki-dat-ta.\\
he-\textsc{nom} see-\textsc{pass-pst} eye-\textsc{top} curious-\textsc{cop-pst}\\
(`The eye/gaze he was viewed with was curious.’)
\cg. *totemo kōki-no me\\
very curious-\textsc{mod} eye\\
(`very curious gaze’)

Other examples include the modifiers \textit{higō-no} `unexpected’ (47 attr., 0 pred.), which occurs only with \textit{shi} `death’ and \textit{saigo} `end’ (also in the sense of `death’) \ref{higo-no} and \textit{anmoku-no} `tacit’ (350 attr., 0 pred.), mostly with \textit{ryōkai} `agreement’.\footnote{Two things are noteworthy here: First, the combination of \textit{higō} `unexpected’ and the two nouns are in 32 out of 47 cases part of the larger expression \textit{higō-no shi/saigo-wo toge-ru} `to meet an unexpected death’. Second, the modifier \textit{anmoku-no} `tacit’ does not yield a single instance of predicative use in the \is{BCCWJ}BCCWJ corpus, but some instances can be found on Google. Nevertheless, it is also given in the literature \citep{muraki:2012:nihongo, lehmann:2015:japanisch} as a non-predicative modifier.}

\ex. \label{higo-no} \ag. Nanninka-no Kirisuto-kyōtō-wa higō-no shi-wo toge-ta.\\
several.people-\textsc{mod} Christianity-believer-\textsc{top} unexpected-\textsc{mod} death-\textsc{acc} meet-\textsc{pst}\\
`Several Christians met an unexpected death.’\\ Ōno, Kazumichi (2001): Translation of Michelet, Jules: \textit{Bible de l’humanité}. Tokyo: Fujiwara Shoten, original: 1864, via BCCWJ, 2023/11/09.\footnote{Note that potentially, sticking to the religious dimension of the modifier, a translation along the line of `undignified’ might be more appropriate.}
\bg. *Nanninka-no Kirisuto-kyōtō-ga toge-ta shi-wa higō-d-at-ta.\\
several.people-\textsc{mod} Christianity-believer-\textsc{nom} meet-\textsc{pst} death-\textsc{top} unexpected-\textsc{pred-cop-pst}\\
(`The death several Christians met was unexpected.’)

These data points suggest three things. First, \is{idiomatic modifier}idiomatic expressions do exist in Japanese. Second, they seem to offer another source for direct modification. Third, they are mostly found in the domain of \textit{no}-modifiers.

As shown in Section \ref{japanesedp}, \citet{svenonius:2008:position} assumes that idiomatic modifiers are merged as specifiers of category-less root phrases $\sqrt{}$P \citep{marantz:1997:no}. These root phrases are embedded below nP and receive their category from NP. Descriptively speaking, this means that idiomatic modifiers are at the very core (root) of the noun \citep[36–37]{svenonius:2008:position}. Hence, their position is the lowest inside the direct domain, and therefore within the whole DP.

Concrete evidence for this claim comes from the fact that nothing can intervene between the idiomatic modifier and the head noun, i.e. the idiomatic modifier cannot be higher in the structure. This is shown for \is{English} English in \citet[36–37]{svenonius:2008:position}, and also holds in Japanese \ref{koki low}.

\ex. \label{koki low} \ag. minna-no kōki-no manazashi\\
everyone-\textsc{mod} curious-\textsc{mod} gaze\\
`everyone's curious gaze'
\bg. *kōki-no minna-no manazashi\\
curious-\textsc{mod} everyone-\textsc{mod} gaze\\

Crucially, this behavior is not due to \textit{minna-no} `everyone’s’. As shown below, this modifier can intervene between non-idiomatic modifiers and the head noun.

\ex. \ag. minna-no aka-i doresu\\
everyone-\textsc{mod} red-\textsc{mod} dress\\
\bg. aka-i minna-no doresu\\
red-\textsc{mod} everyone-\textsc{mod} dress\\
`everyone's red dresses'

The adjacency facts apply to other idiomatic modifiers as well, for example the combination \textit{aka-no tan'in} `red stranger', tested with an intervening relative clause.

%

\ex. \ag. [dare-mo shira-na-i] aka-no tan'in\\
everyone-\textsc{emp} know-\textsc{neg}-\textsc{mod} red-\textsc{mod} stranger\\
`the total stranger that no one knows'
\bg. *aka-no [dare-mo shira-na-i] tan'in\\
red-\textsc{mod} everyone-\textsc{emp} know-\textsc{neg-mod} stranger\\

From these facts, I derive the integration of \is{idiomatic modifier}idiomatic modifiers as depicted in \figref{Structure idiomatic DP}, repeated from Section \ref{japanesedp}, but now including the internal structure of the idiomatic modifier under ModP. Recall that the root phrase equals the label `NP' in the system of \citet{svenonius:2008:position}. I display root phrases as dedicated functional projections throughout (especially in the larger trees) to highlight their status and facilitate readibility.

 \begin{figure}
\centering
\begin{tikzpicture}[baseline=0pt, remember picture]
\tikzset{level distance=25pt, sibling distance=5pt}
\Tree [.DP\textsubscript{internal} {} [.D′\textsubscript{internal} [.FP\textsubscript{indirect} {} [.F′\textsubscript{indirect} [.dP [.FP\textsubscript{direct} {} [.F′\textsubscript{direct} [.nP [.$\sqrt{}$P [.ModP {} [.Mod′ [.XP\textsubscript{idiomatic} \edge[roof]; kōki ] [.Mod no ] ] ] [.$\sqrt{}$′ \edge[roof]; manazashi ] ] n ] F ] ] d ] F ] ] D ] ]
 \end{tikzpicture}
\caption{Structure of DP with idiomatic modifier}
\label{Structure idiomatic DP}
\end{figure}

Two more things are noteworthy. First, this root phrase \textit{only} hosts \is{idiomatic modifier}idiomatic modifiers in Japanese. \citet[130]{kim:2019:syntax}, who adopts the category $\sqrt{}$P from \citet{svenonius:2008:position}, argues that this phrase hosts idiomatic modifiers, but also modifiers which receive an internal theta role from the noun, thus compound-like constructions, in Korean. As evidence, she shows that no other modifier can intervene between such modifiers and the head noun in Korean either, parallel to idiomatic constructions. In the following example, \textit{kyengcey} `economy' is the internal argument of \textit{kayhyek} `renovation' and must be immediately adjacent to the noun \citep[130]{kim:2019:syntax}, essentially reflecting the category of \is{relational modifier}relational modifiers discussed above. \is{Korean}

\ex. \ag. sin kyengcey kayhyek (Korean)\\
new economy renovation\\
`a new renovation of the economy'
\bg. *kyengcey sin kayhyek\\
economy new renovation \null \hfill \citep[130]{kim:2019:syntax}\\

However, this is not the case in Japanese, as the (approximate) translation of the examples shows. Here, permutation is easily possible.\footnote{This also extends to some examples given by \citet{laenzlinger:2011:elements}, such as \ref{machi-no saikin-no}. Here, a thematic modifier with a \textsc{theme} role (\textit{machi-no} `city’) and another modifier (\textit{saikin-no} `recent’), according to \citet{laenzlinger:2011:elements} have to appear in this order, thus showing the opposite pattern to Kim’s.

\par

\exg. machi-no saikin-no hakai\\
city-\textsc{mod} recent-\textsc{mod} destruction\\
`the recent destruction of the city’ \hfill \citep[227–228]{laenzlinger:2011:elements} \label{machi-no saikin-no}

\par


However, this DP is probably not very natural to begin with. Furthermore, without context, permutation is easily possible as well.}

\ex. \ag. arata-na keizai-no kaikaku\\
new-\textsc{mod} economy-\textsc{mod} renovation\\
\bg. keizai-no arata-na kaikaku\\
economy-\textsc{mod} new-\textsc{mod} renovation\\
`new renovation of the economy'

Second, \citet{okami:2012:two} has argued that those \textsc{color} modifiers which can appear either with \textsc{i} or with \textsc{no} can only receive an \is{idiomatic modifier}idiomatic reading in the \textit{i}-variant. She gives the following examples.

\ex. \ag. Tarō-to Hanako-wa aka-i / *aka-no ito-de musub-are-te-i-ru.\\
Taro-and Hanako-\textsc{top} red-\textsc{mod} / *red-\textsc{mod} string-\textsc{ins} bind-\textsc{pass-ger}-be-\textsc{prs}\\
`Taro and Hanako are destined to be together.’ (Lit. `bound by a red string’)
\bg. Kare-ni-wa kuro-i / *kuro-no uwasa-ga ar-u.\\
he-\textsc{dat-top} black-\textsc{mod} / *black-\textsc{mod} rumor-\textsc{nom} be-\textsc{prs}\\
`There are dark rumors about him.’ \hfill \citep[96]{okami:2012:two}

Interestingly, such color items do allow permutation when occurring with another modifier, as opposed to the examples above.

\ex. \ag. Kare-ni-wa hen-na kuro-i uwasa-ga ar-u.\\
he-\textsc{dat-top} strange-\textsc{mod} black-\textsc{mod} rumor-\textsc{nom} \textsc{cop-prs}\\
\bg. Kare-ni-wa kuro-i hen-na uwasa-ga ar-u.\\
he-\textsc{dat-top} black-\textsc{mod} strange-\textsc{mod} rumor-\textsc{nom} \textsc{cop-prs}\\
`There are strange dark rumors about him.’

This might suggest that these are in fact not \is{idiomatic modifier}idiomatic modifiers, or at least they are not merged in the same position (the root phrase).

\subsection{Quantificational modifiers}

As already shown in Section \ref{quantificational modifiers}, another source for direct modifiers are quantifiers or more precisely \is{quantificational modifier}quantificational modifiers. Many quantificational modifiers are direct-only, while some are only prototypically attributive. Examples of quantificational modifiers which cannot be used predicatively include \textit{kazukazu-no} `many’, \textit{amata-no} `a lot, many’ and \textit{shuju-no} `various’.

\ex. \ag. Kare-wa kazukazu-no shō-wo tot-ta.\\
he-\textsc{top} many-\textsc{mod} prize-\textsc{acc} take-\textsc{pst}\\
`He received many prizes.’
\bg. *Kare-ga tot-ta shō-wa kazukazu-desu.\\
he-\textsc{nom} take-\textsc{pst} prize-\textsc{top} many-\textsc{cop.pol}\\
(`The prizes he received were many.’)

\ex. \ag. Kare-wa shuju-no / amata-no shikaku-wo mot-te-i-ru.\\
he-\textsc{top} various-\textsc{mod} / many-\textsc{mod} qualification-\textsc{acc} have-\textsc{ger}-be-\textsc{prs}\\
`He has various / many qualifications.’
\bg. *Kare-ga mot-te-i-ru shikaku-wa shuju-desu / amata-desu.\\
he-\textsc{nom} have-\textsc{ger}-be-\textsc{prs} qualification-\textsc{top} various-\textsc{cop.pol} / many-\textsc{cop.pol}\\
(`The qualifications he has are various / many.’)

In Section \ref{quantificational modifiers}, it was argued that \is{quantificational modifier}quantificational modifiers are merged first inside the direct domain, since they exhibit direct behavior, but can subsequently move to the \is{left periphery}left periphery, as they can precede not only demonstratives \ref{kazukazuano} but also other direct-only modifiers \ref{takusankomugi4}. Probably they do this to check a [quantificational] feature.

\ex. \ag. kazukazu-no ano mondai\\
many-\textsc{mod} that problem\\
`those many problems’ \label{kazukazuano}
\bg. takusan-no komugi-no pan\\
a.lot-\textsc{mod} wheat-\textsc{mod} bread\\
`many wheat(en) (pieces of) bread’ \label{takusankomugi4}

\subsection{Rentaishi} \label{rentaishi}

Finally, I will discuss the controversial word class of \is{rentaishi@\textit{rentaishi}}\textit{rentaishi} \textjpn{連体詞}, also called \is{fukutaishi@\textit{fukutaishi}}\textit{fukutaishi} \textjpn{副体詞}, arguably a language-specific source for direct-only modifiers. This word class is comprised of several lexemes with different morphology, derived from other word classes and now fossilized so that they are only used as attributive modifiers. This is the reason why \textit{rentaishi} are typically named by many Japanese linguists working in more traditional frameworks as non-predicative lexemes whose only function is to modify nouns making them the only group of direct modifiers acknowledged in the literature for a long time (see the newer works mentioned throughout this chapter).

The word class \textit{rentaishi} makes up a controversial case, first, because it was created more upon historical and ideological than on syntactic grounds at the beginning of the 20th century, and second, since it is not accepted as an individual word class by all linguists, both in- and outside of Japan. This is also the reason why many linguists would probably not classify the lexemes presented in the following pages as direct-only modifiers, since they strongly resemble fixed expressions. Nevertheless, I will show that they are quite close to idiomatic, and sometimes also non-intersective, modifiers; and since the relevant lexemes only serve as attributive modifiers, they deserve a discussion here. I will give a brief introduction, and then discuss different subgroups. For more information the reader is referred to \citet{kohlich:2020:rentaishi, kohlich:2023:uninflecting}. \is{word classes}

\subsubsection{Introduction: Terminology, diachrony, criteria, subgroups}

The term \is{rentaishi@\textit{rentaishi}}\textit{rentaishi} literally translates to \textit{adnominal-noun-word} and is in English usually translated as `adnoun’ or `prenoun’. German equivalents are \textit{Adnomen}, sometimes \textit{Adnominalie} \citep{rickmeyer:1995:japanische}. \textit{Rentaishi} were incorporated into the Japanese word class system at the beginning of the Meiji-Period (1868--1912). The important linguists of the time, especially MATSUSHITA Daizaburō (in \citealt{matsushita:1976:kaisen}), alongside many others \citep{yamada:1908:nihon, mitsuya:1932:bunpo, sekine:1937:fukutaishi, kieda:1938:koto, tokieda:1941:kokugogaku, yuzawa:1979:kokugo}, realized that there are several lexemes in Japanese that are restricted to attributive use, do not inflect and whose sole function is to modify nouns. They also realized that such lexemes bear striking similarities to the \textit{defective} adjectives in Western (mostly Dutch and English) languages, which only inflect dependent on the noun they modify \citep{sugimoto:1967:kindai, sugimoto:1991:kokugo, kim:2006:kindai}. The new word class \textit{rentaishi} was created to incorporate all those defective lexemes, becoming an equivalent of the Western adjective. \is{Meiji-Period}

\largerpage
Another important criterion of this \is{rentaishi@\textit{rentaishi}}word class is the morphological diversity of its members. Being fossilized derivates from other word classes and frozen in their morphological form and usage domain explains not only the restriction to attributive position, it also predicts the absence of other inflectional forms. This is a significant point of deviation from all major Japanese word classes which since the language reform of the Meiji-Period have been predominantly categorized on the basis of their shared morphological paradigms \citep{nitta:2000:bun}.

Starting with an example, there are several lexemes which look verbal on a morphological level -- because they are indeed derived from verbs -- but only appear as prenominal modifiers. In fact, all those verbal \is{rentaishi@\textit{rentaishi}}\textit{rentaishi} are derived from verbs fossilized in their attributive inflectional form (\textit{rentaikei}). They include verbal derivates ending in the simple non-past form such as \textit{aru} `(a) certain’ \ref{aru},\footnote{This \textit{aru} is homophonous to the existential/copular verb \textit{ar}-, which is why some authors \citep{lehmann:2015:japanisch} treat these two elements as the same. However, the copula verb \textit{ar}- is generally incompatible with animate nouns, therefore it cannot be the same element as \textit{aru} in \ref{aru}.} derivates ending in the morpheme -\textit{iu}, which signaled passive, potential or spontaneity, and itself inflected for attributive form, such as \textit{iwayuru} `so-called’ \ref{iwayuru}, and those ending in the past morpheme -\textsc{ta}, such as \textit{taishita} `considerable’ \ref{taishita}. None of these modifiers can be used predicatively or is gradable, highlighting the role of \textit{rentaishi} for the present study.

\ex. \label{aru} \ag. aru gakusei\\
certain student\\
`a certain student'
\bg. *Gakusei-wa aru.\\
student-\textsc{top} certain\\
(`This student is a certain.’)
\cg. *motto / totemo aru gakusei\\
more / very certain student\\
(`a more / very certain student’)

\ex. \label{iwayuru} \ag. iwayuru gakusei\\
so-called student\\
`a so-called student'
\bg. *Kono gakusei-wa iwayuru.\\
this student-\textsc{top} so-called\\
(`This student is so-called.’)
\cg. *motto / totemo iwayuru gakusei\\
more / very so-called student\\
(`a more / very so-called student’)

\ex. \label{taishita} \ag. taishita sa\\
considerable difference\\
`a considerable difference’
\bg. *Sa-wa taishita.\\
difference-\textsc{top} considerable\\
(`The difference is considerable.’)
\cg. *motto / totemo taishita sa\\
more / very considerable difference\\
(`more / very considerable difference’)

Other potential members of \is{rentaishi@\textit{rentaishi}}\textit{rentaishi} are derived from adjectives and appear in the attributive form of the premodern copula \textit{nari} or \textit{tari}.\footnote{Those two can be viewed as allomorphs where \textit{tari} was restricted to Sino-Japanese lexemes.} Again, the relevant lexemes in Modern Japanese are restricted to attributive use and are non-gradable.

\ex. \ag. tannaru machigai\\
mere misunderstanding\\
`a mere misunderstanding'
\bg. *Machigai-wa tannaru\\
misunderstanding-\textsc{top} mere\\
(`The misunderstanding is mere.’)
\cg. * motto / totemo tannaru machigai\\
more / very mere misunderstanding\\
(`more / very mere misunderstanding’)

\ex. \ag. shutaru gen'in\\
main reason\\
`the main reason'
\bg. *Gen'in-wa shutaru.\\
reason-\textsc{top} main\\
(`The reason is main.’)
\cg. *motto / totemo shutaru gen'in\\
more / very main reason\\
(`more / very main reason’)

\newpage
Recall the lexemes \textit{ōki} `big' and \textit{chiisa} `small’, which can either co-occur with \textsc{i} or with \textsc{na}, but never with the copula \textsc{da}. The \textit{na}-variant of these lexemes has traditionally also been called \is{rentaishi@\textit{rentaishi}}\textit{rentaishi}.\footnote{Another such case is \textit{okashi} `strange’. The question remains whether there is one stem with two possible attributive morphemes, or if all three of these lexemes are derived from an \textit{i}-adjective-counterpart, have become fossilized and are now lexicalized as uninflectable \textit{rentaishi}. Both analyses are available (see also \citealt[Footnote 15]{oshima:2021:nihongo}), but I will assume that \textit{ōki-i} and its counterparts are \textit{i}-adjectives, whereas \textit{ōki-na} belong to a different word class. This conditions the spell-out of \textsc{mod}. See also Section \ref{complexdirect}.}

\exg. *Kono ie-wa ōki-da.\\
this house-\textsc{top} big-\textsc{cop}\\
(`This house is big.’) \null \hfill (= \ref{ieokida})

Finally, there are some lexemes which end in /ga/, but are counted as \textit{rentaishi}, namely \textit{waga} `my/our' and \textit{onoga} `my’, and some that end in /no/, such as \textit{honno} `only, mere’.\footnote{The nominative case particle, -\textit{ga} is the diachronic genitive marker, accounting for the archaic flavor of DPs in which the relevant \textit{rentaishi} ending in /ga/ appear.}

In this sense, it is debatable whether one should analyze \textit{honno} as a fossilized lexeme, with /no/ being part of the stem, or whether one should maintain the analysis presented here of /no/ as a syntactic head. An argument for the former is that the first result for \textit{honno} on the \is{CHJ} Corpus of Historical Japanese (CHJ) dates back to the year 900. Maintaining an analysis of \textit{rentaishi} as a closed class of historical lexemes, I tentatively opt for this analysis, although nothing immediately hinges on this. Note, however, that this also means that not every direct-only \textit{no}-modifier (sorted in group 5) in the data set is a \textit{rentaishi}, and \textit{rentaishi} is here not treated as a feature, but as a word class of historical derivates, despite recent efforts of several authors to argue for exactly that \citep{muraki:1996:imi, muraki:2009:nihongo, muraki:2012:nihongo, muraki:2015:hinshi, oshima:2021:nihongo, oshima:2021:redefining, abe:2022:no}. See also \citet{kohlich:2023:uninflecting} for discussion.\is{CHJ}

\tabref{tab:listrentaishi} gives an overview of the morphological subgroups of \textit{rentaishi}.\footnote{This
    list is compiled from \citet{mitsuya:1932:bunpo, kieda:1938:koto, suzuki:1972:nihon, komatsu:1973:rentaishi, tokieda:1941:kokugogaku, mio:1958:hanashi, kai:1980:rentaishi, kim:2006:kindai, matsubara:2009:nihongo}, see \citet[181–182]{kohlich:2020:rentaishi}. Traditionally, also the demonstratives \textit{kono} `this' and \textit{kon-na} `such a’ have been analyzed as \textit{rentaishi} in the relevant studies. I will not include them in the present study where they are treated as demonstratives.}

\begin{table}[t]
\caption{List of morphological groups among the group of \textit{rentaishi}}
\label{tab:listrentaishi}
\begin{tabularx}{\textwidth}{lQQ}
\lsptoprule
\textbf{Group} & \textbf{Examples} & \textbf{Translation} \\
\midrule
verbal \textit{rentaishi} & verbs inflected for attributive form, non-past & \textit{aru} `(a) certain'; \textit{saru} `previous'; \textit{akuru} `following'; \textit{kitaru} `ahead' \\
\tablevspace
& verbs combined with verbal element -\textit{iu}, which is inflected for attributive form & \textit{iwayuru} `so-called'; \textit{arayuru} `all kinds of'\\
\tablevspace
& verbs ending in past morpheme & \textit{taishita} `considerable'; \textit{tonda} `remarkable'\\
\tablevspace
adjectival \textit{rentaishi} & three lexemes in combination with element -\textsc{na} & \textit{ōkina} `big'; \textit{chiisana} `small'; \textit{okashina} `strange'\\
\tablevspace
& lexemes ending in attributive forms of premodern copula -\textit{naru}/\textit{taru} & \textit{tannaru} `mere'; \textit{shutaru} `main'\\
\tablevspace
\textit{rentaishi} with \textit{ga} and \textit{no} & lexemes ending in -\textit{ga}, often in certain archaic sayings & \textit{waga} `my/our'; \textit{onoga} `my'\\
\tablevspace
& lexemes ending in -\textit{no} & \textit{rei-no} `ordinary, aforementioned'; \textit{hon-no} `only, mere'\\
\lspbottomrule
\end{tabularx}
\end{table}
\newpage
\subsubsection{Summary: \textit{rentaishi}}

The morphologically diverse group of \is{rentaishi@\textit{rentaishi}}\textit{rentaishi} ``created'' in the beginning of the past century constitutes another group of direct-only modifiers in Japanese, and a language-specific one as such. From a syntactic point of view, it is striking that these relics from Classical Japanese serve an important role in modern syntactic theory. They also fit an important characteristics argued by \citet{cinque:2010:adj} for direct-only modifiers: they are \textit{functional}, meaning they make up a closed group. This is to be expected given the fact that they are fossilized. Furthermore, from a perspective of word class status, these lexemes resist a unifying analysis based on morphological criteria and rather the absence thereof lends itself as a criterion.

Also, note that there are several \is{rentaishi@\textit{rentaishi}}\textit{rentaishi} with similar meanings to \is{non-intersective modifier}non-intersective modifiers \citep{kamp:1995:prototype, panayidou:2013:in} discussed at the beginning of this section. These include \textit{aru} `a certain’, \textit{tannaru} `mere’ and \textit{shutaru} `main’. It is remarkable that an awareness of the solely attributive character of these lexemes existed long before the concept of non-intersectivity had been established.

In a nutshell, this group of historically fossilized attributive modifiers is highly relevant for the discussion on direct modification and brings in a language-specific dimension, adding to the equivalents of cross-linguistic groups of direct modifiers looked at so far.

\subsection{Interim summary}

The preceding discussion has shown that Japanese does exhibit direct modifiers, contrary to what has often been claimed in the literature. When taking stock and applying a large scale empirical validation, the number of direct modifiers is much larger than it is suggested according to the few studies concerned with them. Furthermore, it was again highlighted that the vast majority of direct modifiers can be found among the group of \textit{no}-modifiers.

The following groups of direct modifiers could be established.

\ex. \a. Non-intersective modifiers
\b. Relational modifiers
\c. Qualitative modifiers of standard categories (\textsc{quality} etc.), potentially of different categories
\d. Idiomatic modifiers
\e. Quantificational modifiers
\f. \textit{Rentaishi}

\section{Implications and discussion} \label{discussiondirect}

Having established the groups of direct modifiers present in Japanese, I will now discuss the structure of the direct domain more closely, with special focus on the consequences and implications Japanese offers for Cartography.

First, the following characteristics of the direct domain were identified:

\ex. \a. The direct domain does not consist of functional projections dedicated to specific semantic categories, but of unspecific functional projections.
\b. This has the consequence that no hierarchical order according to these specific functional projections exists, hence no ordering restrictions of modifiers are expected in linear order either.
\c. Idiomatic modifiers occupy a specific position, directly adjacent to the noun.
\d. Quantificational modifiers are direct, but typically move out of the direct domain to the left periphery.

These points inevitably raise critical questions on a cross-linguistic dimension and for any cartographic account. An immediate question is, why is Japanese so lenient when it comes to ordering, but other languages are not, i.e. why do, for instance, English or German dictate strict \is{absence of ordering restrictions}ordering on their modifiers implicating that structural hierarchies are at stake in such languages? But then what is the reason that such hierarchies are active in \is{English} English, but not in Japanese? Since it is a core principle of Cartography that structures are uniform, such deviations should not exist, and thus pose a serious challenge. This also extends to the concept of Universal Grammar (UG) in general. From a UG perspective, the question arises whether this distribution is simply random and, furthermore, whether we only find these two extremes cross-linguistically. Or could it be possible that there exist languages that exhibit a mixed behavior, that are situated between English and Japanese? If so, how would children learn to determine whether they need to acquire ordering restrictions of direct modifiers based on specific categories or not?

These issues will be addressed in this section. I will argue that although the general spine, the structure of the direct domain, is uniform, its inventory is subject to parametric differences. Furthermore, I will show that there are in-between cases and demonstrate this with two languages: \is{Persian} Persian and \is{Modern Greek} Modern Greek. These observations will then lend themselves to cross-linguistic generalizations. In the next subsection, I will discuss ordering in Japanese in more detail.

\subsection{More on ordering in Japanese}

As presented in Section \ref{sectioncarto}, a core assumption of Cartography is that direct qualitative modifiers (generally adjectives) show \is{adjective ordering restrictions}ordering restrictions when two or more occur within the same DP (see \citealt{scott:2002:stacked, cinque:2010:adj, laenzlinger:2011:elements}). This means that whether modifiers adhere to the expected ordering principles or not directly allows to draw generalizations regarding their modification type. The non-canonical order of a \textsc{color} and a \textsc{dimension} adjective in languages such as \is{English} English, for instance, can straightforwardly be analyzed as the result of \textit{red} being merged in the indirect domain, thus higher. Otherwise, the surface order \textit{big $\succ$ red} should be impossible given an out-of-the-blue sentence with no specific context. This is shown below.

\ex. [red\textsubscript{indirect} [big\textsubscript{direct} [book]]]

On the other hand, it was sufficiently demonstrated that ordering and modification type are orthogonal issues in Japanese and that (most) modifiers, whether ambiguous or direct-only, are not expected to show ordering restrictions. Thus, the situation given above does not hold automatically.

The general freedom in \is{absence of ordering restrictions}ordering between such qualitative modifiers was quantitatively measured in an earlier study \citep{kohlich:2018:die}. In this study, which was carried out on the \is{BCCWJ}BCCWJ, I tracked combinations of modifiers belonging to the categories \textsc{quality}, \textsc{dimension}, \textsc{shape} and \textsc{color} inside and across the two canonical adjective groups within the same DP in order to establish whether their semantic category, as well as their word class, has an effect on ordering.

The results are displayed in \tabref{adjective ordering BCCWJ}. The leftmost column depicts the relevant categories of the first adjective in a given DP, the second adjective corresponds to the labels written horizontally in the top row.

\begin{table}
\centering
\begin{tabular}{lrrrr}
\lsptoprule
\textbf{Order} & \textsc{quality} & \textsc{dimension} & \textsc{shape} & \textsc{color} \\
\midrule
\textsc{quality}    & 290 & 40 & 16 & 33 \\
\textsc{dimension}  & 45  & 57 & 16 & 18 \\
\textsc{shape}      & 8   & 4  & 5  & 6 \\
\textsc{color}      & 15  & 19 & 10 & 5 \\
\lspbottomrule
\end{tabular}
\caption{Ordering among Japanese adjectives on the BCCWJ}
\label{adjective ordering BCCWJ}
\end{table}

 Although admittedly there are not many co-occurrences in the first place, which has to do with the design and size of the corpus (see \citealt{kohlich:2023:suppl}), the results confirm a \is{absence of ordering restrictions}lack of ordering restrictions and confirm the overall theme of the present study.\footnote{Seemingly relevant is the skew towards \textsc{quality} $\succ$ \textsc{color} over the reverse order, but then again, both categories do not show restrictions with regard to \textsc{dimension}. On the other hand, \textsc{quality} seems to precede \textsc{shape}. In general, the numbers presented here are too little to draw any substantial conclusions from.} \is{BCCWJ}

The present study adds to this observation that ordering restrictions are also absent amongst and across direct-only modifiers. This holds for pairs of \is{relational modifier}relational modifiers, combinations of relational and qualitative modifiers, and relational with \is{non-intersective modifier}non-intersective modifiers. The different positions modifiers exhibit in linear order on the surface are simply the results of the different positions they occupy inside the DP and the order of merge. For instance, the two orders in example \ref{tada} above fall out from two different merge positions. The non-intersective modifier \textit{tada-no} `mere’ precedes the relational modifier \textit{Doitsu-no} `German’ in linear order if it occupies a higher position, i.e. if it is merged later. On the other hand, the reverse order \textit{Doitsu-no tada-no gakusei}, Lit. `German mere student', ungrammatical in English, is a result of \textit{tada-no} occupying a lower position, being merged earlier. These options are illustrated in \figref{tada vs. doitsu}.

% \begin{figure}
% \begin{subfigure}{\textwidth}
% \centering
% \resizebox{.8\textwidth}{!}{
%  \begin{tikzpicture}[baseline=0pt, remember picture]
% \tikzset{level distance=25pt, sibling distance=5pt}
% \Tree [.DP\textsubscript{internal} {} [.D′\textsubscript{internal} [.FP\textsubscript{indirect} {} [.F′\textsubscript{indirect} [.dP [.FP\textsubscript{direct} [.ModP {} [.Mod′ [.XP \edge[roof]; tada ] [.Mod no ] ]] [.F′\textsubscript{direct} [.FP\textsubscript{direct} [.ModP {} [.Mod′ [.XP \edge[roof]; Doitsu ] [.Mod no ] ] ] [.F′\textsubscript{direct} [.NP \edge[roof]; gakusei ] F ] ] F ] ] d ] F ] ] D ] ]
% \end{tikzpicture}
% }
% \caption{Structure of \textit{tada-no Doitsu-no gakusei} `mere German student’ (= \ref{tada1})}
% \label{structure tada-no Doitsu-no gakusei}
% \end{subfigure}
%
% \begin{subfigure}{\textwidth}
% \centering
% \resizebox{.7\textwidth}{!}{
%  \begin{tikzpicture}[baseline=0pt, remember picture]
% \tikzset{level distance=25pt, sibling distance=5pt}
% \Tree [.DP\textsubscript{internal} {} [.D′\textsubscript{internal} [.FP\textsubscript{indirect} {} [.F′\textsubscript{indirect} [.dP [.FP\textsubscript{direct} [.ModP {} [.Mod′ [.XP \edge[roof]; Doitsu ] [.Mod no ] ]] [.F′\textsubscript{direct} [.FP\textsubscript{direct} [.ModP {} [.Mod′ [.XP \edge[roof]; tada ] [.Mod no ] ]] [.F′\textsubscript{direct} [.NP \edge[roof]; gakusei ] F ] ] F ] ] d ] F ] ] D ] ]
% \end{tikzpicture}
% }
% \caption{Structure of \textit{Doitsu-no tada-no gakusei} `German mere student’ (= \ref{tada2})}
% \label{structure Doitsu-no tada-no gakusei}
% \end{subfigure}
%
% \caption{Two possible orders as result of different structures for \ref{tada}}
% \label{tada vs. doitsu}
%
% \end{figure}

\begin{figure}
\small
\begin{subfigure}{.49\textwidth}
\centering
\begin{forest}
for tree={
  l sep=25pt,
  s sep=5pt,
}
[DP\textsubscript{internal}
  [{}]
  [D′\textsubscript{internal}
    [FP\textsubscript{indirect}
      [{}]
      [F′\textsubscript{indirect}
        [dP
          [FP\textsubscript{direct}
            [ModP
              [{}]
              [Mod′
                [XP
                  [tada, roof]
                ]
                [Mod
                  [no]
                ]
              ]
            ]
            [F′\textsubscript{direct}
              [FP\textsubscript{direct}
                [ModP
                  [{}]
                  [Mod′
                    [XP
                      [Doitsu, roof]
                    ]
                    [Mod
                      [no]
                    ]
                  ]
                ]
                [F′\textsubscript{direct}
                  [NP
                    [gakusei, roof]
                  ]
                  [F]
                ]
              ]
              [F]
            ]
          ]
          [d]
        ]
        [F]
      ]
    ]
    [D]
  ]
]
\end{forest}
\caption{Structure of \textit{tada-no Doitsu-no gakusei} `mere German student’ (= \ref{tada1})}
\label{structure tada-no Doitsu-no gakusei}
\end{subfigure}
\begin{subfigure}{.49\textwidth}
\centering
\begin{forest}
for tree={
  l sep=25pt,
  s sep=5pt,
}
[DP\textsubscript{internal}
  [{}]
  [D′\textsubscript{internal}
    [FP\textsubscript{indirect}
      [{}]
      [F′\textsubscript{indirect}
        [dP
          [FP\textsubscript{direct}
            [ModP
              [{}]
              [Mod′
                [XP
                  [Doitsu, roof]
                ]
                [Mod
                  [no]
                ]
              ]
            ]
            [F′\textsubscript{direct}
              [FP\textsubscript{direct}
                [ModP
                  [{}]
                  [Mod′
                    [XP
                      [tada, roof]
                    ]
                    [Mod
                      [no]
                    ]
                  ]
                ]
                [F′\textsubscript{direct}
                  [NP
                    [gakusei, roof]
                  ]
                  [F]
                ]
              ]
              [F]
            ]
          ]
          [d]
        ]
        [F]
      ]
    ]
    [D]
  ]
]
\end{forest}
\caption{Structure of \textit{Doitsu-no tada-no gakusei} `German mere student’ (= \ref{tada2})}
\label{structure Doitsu-no tada-no gakusei}
\end{subfigure}

\caption{Two possible orders as result of different structures for \ref{tada}}
\label{tada vs. doitsu}

\end{figure}

Although both modifiers are situated in the direct domain, which has been verified viz. their status as direct modifiers, their exact position is not predetermined and this scenario is at stake for all direct modifiers with interchangeable ordering.

\subsection{Language variation in ordering} \label{language variation}

Assuming a broad direct domain explains the absence of ordering restrictions in Japanese. But it yet fails to explain stricter ordering restrictions in languages such as \is{English} English or \is{German} German, which over the years have justified the inventory of functional projections based on specific semantic categories \citep{cinque:2010:adj, cinque:2025:superiority, laenzlinger:2011:elements}. The question is: what makes Japanese so different from English, i.e. what determines whether a language has a broad domain with unspecific, or a more articulated domain with specific functional projections?

\subsubsection{Universal Spine Hypothesis and Units of Language} \label{units of language} 

To answer these questions, I draw ideas from Wiltschko's \textit{Universal Spine Hypothesis} \is{Universal Spine Hypothesis}(USH, \citealt[24–26]{wiltschko:2014:universal}), which offers a radically different perspective from \is{Units of Language}Cartography. I argue that this hypothesis enriches this research program in various ways and offers a solution to the questions presented above.\footnote{For the specific spine and its projections proposed by Wiltschko as an alternative to the ones proposed here, see \citet[§ 2]{wiltschko:2014:universal}.} I adopt the following main ingredients (cf. \citealt[§ 1, 2]{wiltschko:2014:universal}). \is{Universal Spine Hypothesis} \is{Units of Language}

\newpage
\ex. \a. There is a universal syntactic spine which holds across all languages and is part of UG.
\b. Across this spine, we find a universal inventory of hierarchically ordered categories.
\c. Whether these categories surface in a specific language, however, is subject to language-specific factors.
\d. Each language has its own language-specific entities, called \textit{Units of Language} in \citet[1]{wiltschko:2014:universal}, which derive language-specific categories from universal categories.

The basic idea is that of a universal DP-spine with three different domains and functional heads within these domains licensing modifiers and giving landing sites. This general DP-spine and its three domains are uniform cross-linguistically: the left periphery, the indirect domain, and the direct domain. The nature of the functional heads are shared with standard cartographic assumptions as well. Functional heads within the direct and indirect domain license direct and indirect modifiers respectively and the left periphery serves mainly as a landing site for modifiers which move there for various reasons. \is{Universal Spine Hypothesis} \is{Units of Language}

Furthermore, while the general spine is fixed, and thus variation limited across domains cross-linguistically, variation in the architecture, and thus in the surface order, comes into play within the same domain and this variation is language-specific (for a similar idea in the clause see \citealt{ramchand:2014:deriving}). Specifically, depending on its language-specific inventory, or its Units of Language (UoLs; \citealt{wiltschko:2014:universal}), each language realizes the categories in each domain in a different way, albeit those differences are systematic.

Languages such as \is{English} English possess the Units of Language to realize specific categories of direct modifiers such as \textsc{quality} and \textsc{color}. They then combine those with the cross-linguistic spine and inventory of functional projections for direct modifiers to form the language-specific categories F\textsubscript{\textsc{quality}} and so on. These functional heads license the respective modifiers in the structure and are ordered hierarchically. In linear order, this \is{Universal Spine Hypothesis} \is{Units of Language} predicts the fixed order of these direct modifiers.

Japanese, on the other hand, lacks the specific Units of Language for \textsc{quality} and similar categories, therefore does not realize specific functional projections which license modifiers of \textsc{quality} and which are ordered hierarchically. However, the universal DP is shared and this includes the direct domain. The difference between the functional heads and projections this domain contains and those of the \is{English} English type is that the former are unspecific, and not specified for modifiers of specific semantic meanings, word classes, or levels of gradability, and other factors. This means, the functional heads within the Japanese direct domain are also not ordered in a specific hierarchy, and this correctly predicts the absence of ordering restrictions on the surface. \is{Universal Spine Hypothesis} \is{Units of Language}

From a cross-linguistic perspective, and a perspective of UG, this means that children must acquire the relevant Units of Language required by the language they are learning in order to see which categories need activating and are embedded into a hierarchy.

Nevertheless, it is evident that the general hierarchy of functional projections in a given domain, if a given language realizes them, is uniform as well, considering the fact that the cross-linguistically attested ordering restrictions exemplified by modifiers of different categories are very similar across languages and variation can be explained by the application of different movement operations \citep{cinque:2023:linearization, cinque:2025:superiority}. Therefore, the process of acquiring the relevant hierarchy is not random, but conditioned on acquiring the relevant functional categories. \is{Universal Spine Hypothesis} \is{Units of Language}

As argued in \citet{wiltschko:2014:universal}, the Universal Spine Hypothesis fares better than its alternatives in a lot of scenarios. The Universal Base Hypothesis (UBH) of early generativism \citep{chomsky:1965:aspects} assumes uniformity, universality and identical structural make-up across all languages, the very core idea of present Cartography \citep[68]{cinriz:2010:carto}. However, it faces two concrete problems: The first are overly specific categories in a certain language that should be present in all languages, yet are not attested. Equally problematic are cases of missing universal categories in a specific language (cf. \citealt[10–19]{wiltschko:2014:universal} and § 2), the exact problem observable in Japanese, when it comes to ordering, and essentially the problem of plenitude \citep{larson:2021:rethinking}. \is{problem of plenitude}

The No Base Hypothesis (NBH), on the other hand, attributes a special status to all languages in the sense that every category is language-specific. There are several counter arguments to this hypothesis, but only considering the data points here, it would not explain why the category of direct modifiers is attested in almost every language and why we observe such similar patterns in ordering restrictions across languages. Such a striking commonality does not come about by chance, but instead by cross-linguistic grammatical factors. Without a grammatical baseline, \citet[19–23]{wiltschko:2014:universal} asks, how come learners of so many different languages all acquire the same patterns?

One might object that the Universal Base Hypothesis is, nevertheless, active in Japanese, but that this language makes more use of movement than other languages. There is a logical problem with this objection. If there was an underlining order of modifiers, then we should be able to identify one order as the base order, or, alternatively, identify factors that ameliorate or deteriorate one order over the other. It should not be the case that both orders are equally perfectly acceptable. That is, if we ask speakers, and especially if we look at corpus data, we should see that one order is still prevailing. That is because movement is always an additional option that the speaker has to perform and the listener has to parse. Therefore, I argue, the option of two different underlying orders is much more attractive than assuming one underlying order plus movement. \is{Universal Spine Hypothesis} \is{Units of Language}

Having established that Japanese lacks the language-specific inventories to realize fine-grained functional projections for direct-only modifiers, while languages such as English do realize said projections, leading to a strict marking of ordering in the latter language, the question is how far language variation is allowed. Can languages do whatever they want? Or are there some principles they have to adhere to? Taking the strict English-type and the loose Japanese-type as two extremes, one would expect that there is some ``middle-type'' or even a continuum across languages. This is essentially what is predicted by the claims made above. There should be languages that have the Units of Language for some categories but not for others.

If we find such a language, the next question is which intermediate levels of ordering emerge. That means, if some categories are ordered strictly, but some are not, which are ones that show strict ordering? Are they always the same or can they be different? Such an in-between language would offer the necessary data for a generalization. In the following, I will show that direct-only, \is{non-intersective modifier}non-intersective and \is{relational modifier}relational modifiers, cross-linguistically are more likely to demand strict ordering, than qualitative modifiers, which are ambiguous in the first place in a lot of languages. In order to argue for this, I will focus on two languages in which qualitative modifiers are ordered quite freely, but once direct-only modifiers are involved, ordering restrictions apply. These languages will be Persian and Modern Greek. The reason for these choice are the existence of previous accounts on ordering in the DP in both languages and the comparatively easy testability.

\subsubsection{Persian} \label{persiananalysis}

I start the discussion with Persian.\footnote{Specifically Fārsi, the Persian variant spoken in Iran. I wish to thank Nasimeh Bahmanian very much for providing many of the examples and long discussions. For previous discussion and similar examples see also \citet{kahne:2014:revisiting}. However, some of his examples were judged as not so clear, which is why I elicited additional examples.} The relevant predictions are easy to test because modifiers with a few exceptions (demonstratives, superlatives, etc.) occur only postnominally. They are in this position always connected to the noun (or any preceding element) via the so called \is{E{\.z}āfe@\textit{E{\.z}āfe}}\textit{E{\.z}āfe} morpheme (which in fact could be an instance of \is{MOD@\textsc{mod}}\textsc{mod} as it shows a similar distribution to \textsc{no}; see \citealt{toosarvandani:2014:the, gn:2024:morpho}). \is{Universal Spine Hypothesis} \is{Persian}

Like Japanese, but different from \is{English} English, Persian does not show \is{adjective ordering restrictions}ordering restrictions among qualitative modifiers, as illustrated with the familiar example of a \textsc{dimension} and a \textsc{color} modifier. Modifiers of these categories, which most likely have an adjectival status, are acceptable in either order in Persian and no order seems to be more grammatical or preferred over the \is{absence of ordering restrictions}other.

\ex. \textsc{dimension} vs. \textsc{color}
\ag. miz-e qermez-e bozorg\\
table-\textsc{ez} red-\textsc{ez} big\\
\bg. miz-e bozorg-e qermez\\
table-\textsc{ez} big-\textsc{ez} red\\
`big red table'

Similar facts hold for modifiers of \textsc{material} and \textsc{dimension}, as well as \textsc{nationality} and \textsc{quality}.

\ex. \textsc{dimension} vs. \textsc{material}
\ag. sandali-ye felezzi-ye bozorg\\
chair-\textsc{ez} metal-\textsc{ez} big\\
\bg. sandali-ye bozorg-e felezzi\\
chair-\textsc{ez} big-\textsc{ez} metal\\
`big metal chair'

\ex. \textsc{quality} vs. \textsc{nationality}
\ag. goldān-e qašang-e \v{z}āponi\\
vase-\textsc{ez} beautiful-\textsc{ez} Japanese\\
\bg. goldān-e \v{z}āponi-ye qašang\\
vase-\textsc{ez} Japanese-\textsc{ez} beautiful\\
`beautiful Japanese vase'

Note that modifiers of \textsc{material} and \textsc{nationality} can be used predicatively in Persian, in contrast to Japanese. Therefore, I analyze them as qualitative adjectives, analogous to the hierarchy of \citet{scott:2002:stacked}, albeit not as direct-only modifiers. \is{Universal Spine Hypothesis} %\is{Persian}

\exg. Sandali felezzi ast / \v{z}āponi ast.\\
chair metal \textsc{cop} / Japanese \textsc{cop}\\
`The chair is metal(ic) / Japanese.’

However, direct-only modifiers do exist and they demand \is{adjective ordering restrictions}ordering restrictions. This will be shown first for \is{relational modifier}relational adjectives, labeled as such since they are in fact adjectival derivates from nouns, formed by the suffix -\textsc{i}. One example is \textit{haste’i} `nuclear’ derived from the noun \textit{haste} `nucleus, core’. This modifier, just like relational adjectives in other languages, cannot be used predicatively. \is{Universal Spine Hypothesis} \is{Persian}

\exg. *Fizikdān haste-i ast.\\
physicist nuclear-\textsc{ra} \textsc{cop}\\
(`The physicist is nuclear.’)

In \ref{fizikdan}, it is shown, with an example from \citet{kahne:2014:revisiting}, who draws the same conclusion, that the \is{relational modifier}relational adjective has to occur closer to the noun than the qualitative adjective \textit{javān} `young’, which is not a direct-only adjective. The reversed order is illicit. Assuming the postnominal position is a result of NP-movement in a \is{pied-piping}pied-piping fashion \citep{ross:1967:constraints, cinque:2010:adj}, the order of the postnominal modifiers reflects the mirror image of the underlying hierarchy. \is{Persian}

\ex. \label{fizikdan} \ag. fizikdān-e haste-i-ye javān\\
physicist-\textsc{ez} nuclear-\textsc{ra-ez} young\\
`young nuclear physicist'
\bg. *fizikdān-e javān-e haste-i\\
physicist-\textsc{ez} young-\textsc{ez} nuclear-\textsc{ra} \hfill \citep[18]{kahne:2014:revisiting}\\

In the following, I illustrate two more ordering pairs of relational and qualitative adjectives. In each case, the relational adjective must be adjacent to the head noun. Additionally, I show the illicit predicative use of the relational adjective.

\ex. \ag. doktor-e omun-i-ye javān\\
doctor-\textsc{ez} general-\textsc{ra-ez} young\\
`young general practicioner'
\bg. *doktor-e javān-e omun-i\\
doctor-\textsc{ez} young-\textsc{ez} general-\textsc{ra}\\
\cg. *Doktor omun-i ast.\\
doctor general-\textsc{ra} \textsc{cop}\\
(`The doctor is general.’)

\ex. \ag. kāx-e pādšāh-i-ye qašāng\\
palace-\textsc{ez} royal-\textsc{ra-ez} beautiful\\
`beautiful royal palace'
\bg. ??kāx-e qašāng-e pādšāh-i\\
palace-\textsc{ez} beautiful-\textsc{ez} royal-\textsc{ra}\\
\cg. *Kāx pādšāh-i ast.\\
palace royal-\textsc{ra} \textsc{cop}\\
(`The palace is royal.’)

Finally, when \is{relational modifier}a relational adjective co-occurs with a \is{non-intersective modifier}non-intersective modifier such as \textit{qabli} `former’, also direct-only in Persian as shown in \ref{qabli}, the surface \is{adjective ordering restrictions}order is \textit{relational} $\succ$ \textit{non-intersective}. \is{Persian}

\ex. \label{nuclearformer} \ag. dānešmand-e haste-i-ye qabli\\
scientist-\textsc{ez} nuclear-\textsc{ra-ez} former\\
`former nuclear scientist'
\bg. ??dānešmand-e qabli-ye haste-i\\
scientist-\textsc{ez} former-\textsc{ez} nuclear-\textsc{ra}\\

\exg. *Dānešmand qabli ast.\\
scientist former \textsc{cop}\\
(`The scientist is former.’) \label{qabli}

These facts are in line with the literature and findings in other languages. Recall that, truly non-intersective modifiers precede intersective direct modifiers in \is{German} German and \is{English} English (\citealt{trost:2006:das, panayidou:2013:in}; see Section \ref{ordering intersective}), suggesting a higher position for non-intersective modifiers. On the other hand, relational modifiers are hierarchically very low inside the direct domain (\citealt{bortolotto:2016:RA}; see Section \ref{ordering relational}). \is{Persian}

In general, these findings suggest that Persian has the language-specific inventory to form specific projections for relational, and potentially non-intersective, modifiers, but that it lacks the inventory to realize specific functional projections reserved for qualitative modifiers of \textsc{dimension}, \textsc{quality} and so on. In other words, it possesses only the \is{Units of Language}Units of Language responsible for direct-only modifiers (cf. also \citealt[18]{kahne:2014:revisiting}).

 \begin{figure}[t]
\centering
 \begin{tikzpicture}[baseline=0pt, remember picture]
\tikzset{level distance=25pt, sibling distance=5pt}
 \Tree [.DP\textsubscript{internal} {} [.D′\textsubscript{internal} [.FP\textsubscript{indirect} {} [.F′\textsubscript{indirect} [.... [.dP [.FP\textsubscript{non-intersective} XP [.F′\textsubscript{non-intersective} [.FP\textsubscript{direct-n} XP [.F′\textsubscript{direct-n} [.FP\textsubscript{direct-n+1} XP [.F′\textsubscript{direct-n+1} [.FP\textsubscript{relational} XP [.F′\textsubscript{relational} NP F ] ] F ] ] F ] ] F ] ] d ] {} ] F ] ] D ] ]
\end{tikzpicture}
\caption{DP-Structure of Persian}
\label{DP structure Persian}
\end{figure}

I therefore suggest that the direct domain in Persian looks as in \figref{DP structure Persian}. The general DP structure is the same as in Japanese. Furthermore, the direct domain also contains several functional projections which are not based on specific categories, another characteristic shared with Japanese. The important point of deviation is that \is{Persian} Persian does have certain functional projections which are based on specific categories, namely the projections that host relational adjectives. These are placed in a fixed position closer to the head noun. Tentatively, I also assume projections for non-intersective modifiers, which must be higher than those for relational modifiers, judging from \ref{nuclearformer}. Note that the surface order is the result of \is{pied-piping}roll-up movement of the \textit{whose-pictures} type after which all modifiers appear in the mirror-image order \citep{cinque:2017:micro, cinque:2023:linearization, kahne:2014:revisiting}. I do not display the \is{E{\.z}āfe@\textit{E{\.z}āfe}}\textit{E{\.z}āfe} morpheme as a separate element, although it is also a very good candidate for a syntactic head selecting modifiers as also argued in \citet{philip:2012:linkers, toosarvandani:2014:the, franco:2015:linkers} and \citet{gn:2024:morpho}.\footnote{Intriguingly, some combinations of relational with qualitative adjectives were judged better in the unexpected order, potentially for unrelated reasons such as phonology and/or prosody. According to Nasimeh Bahmanian, the following example can be uttered in both orders, but interestingly, the qualitative adjective \textit{šadid} `severe’ seems to be a bit better closer to the head noun than the relational adjective \textit{zistmohiti} `environmental’.

\par

\ex. \ag. bohrān-e šadid-e zistmohit-i\\
crisis-\textsc{ez} severe-\textsc{ez} environmental-\textsc{ra}\\
\bg. bohrān-e zistmohit-i-ye šadid\\
crisis-\textsc{ez} environmental-\textsc{ra-ez} severe\\
`severe environmental crisis'

\par

}



\subsubsection{Greek}

Modern Greek \is{adjective ordering restrictions}behaves similarly to Persian, although predictions are a bit more difficult to test.\footnote{I am very grateful to Jim Kotopoulis for providing the examples in this section and taking the time to discuss them with me. Unless indicated otherwise, the examples are his.} This is so because modifiers can, in principle, appear on both sides of the noun, and because they can appear with an additional determiner, a construction dubbed \textit{determiner spreading} \citep{androutsopoulou:1995:ds}. This option becomes necessary with postnominal modifiers \citep{alex:1998:adj, alex:2003:adj, campos:2004:ds, panagiotidis:2011:ds, lekakou:2012:poly, panayidou:2013:in}. An illustration is the following \is{absence of ordering restrictions}example where the two postnominal modifiers must occur with their own determiner. This is optional when they appear in prenominal position (adapted from \citealt[303]{alex:1998:adj}). \is{Universal Spine Hypothesis} \is{Modern Greek} \is{determiner spreading}

\ex. \ag. to vivlio *(to) kokino *(to) meghalo\\
\textsc{det} book \textsc{det} red \textsc{det} big\\
\bg. to meghalo (to) kokino (to) vivlio\\
\textsc{det} big \textsc{det} red \textsc{det} book\\
`the big red book’

Determiner spreading has often been argued to instantiate a case of predicative modification (see especially \citealt[A3]{cinque:2010:adj} and references therein), but matters are not so clear (see \citealt{lekakou:2012:poly} for examples of adjectives that are \is{non-intersective modifier}intersective but resist predicative use). Hence, I will limit myself to the relevant facts for the discussion. While free order of modifiers has often been pointed out in relevant structures, a welcoming fact for an analysis as indirect modifiers \citep{alex:1998:adj, alex:2007:nounphrase, cinque:2010:adj}, it seems to be the case that the order among standard categories of \is{adjective ordering restrictions}qualitative modifiers is also significantly less strict in structures not involving determiner spreading, compared to \is{English} English or \is{German} German. This is pointed out in \citet[174]{laenzlinger:2011:elements} and was confirmed by Jim Kotopoulis (p.c.). In the following example, as in \is{Persian} Persian and Japanese, the modifiers of \textsc{dimension} and \textsc{color}, leaving factors such as focus, topicalization and emphasis aside, can appear in any order without deteriorating acceptability. Note that we only need to look at the prenominal position since in postnominal position, determiner spreading would be obligatory.  \is{Modern Greek} \is{determiner spreading}

\ex. \textsc{dimension} vs. \textsc{color}
\ag. to meghalo kokino vivlio\\
\textsc{det} big red book\\
\bg. to kokino meghalo vivlio\\
\textsc{det} red big book\\
`the big red book’

As in Persian, when direct-only modifiers are involved, a strict \is{adjective ordering restrictions}ordering can be observed. The surface order between the qualitative adjective \textit{neari} `young' and the \is{relational modifier}relational adjective \textit{piriniki} `nuclear’, derived via the suffix -\textsc{i}, must be \textit{qualitative adjective} $\succ$ \textit{relational adjective} $\succ$ \textit{noun}, the reverse order \ref{piriniki neari} is ungrammatical. Crucially, \textit{piriniki} is impossible in constructions involving \is{determiner spreading}determiner spreading \ref{piriniki ds}, therefore can never appear postnominally \ref{piriniki post}. It is also not readily available in predicative position \ref{piriniki pred}. These facts are interconnected and underline its direct status.\footnote{According to Jim Kotopoulis, \textit{piriniki} can, in some instances, occur predicatively in which case it has a strong flavor of an elliptical construction. One reason for this might be that, contrary to English, Greek does not need an article when a noun appears predicatively, so only the noun needs to be deleted. A similar phenomenon in Japanese was shown at the beginning of this chapter.}  \is{Modern Greek}

\ex. relational vs. qualitative modifier
\ag. i neari piriniki fisikos\\
\textsc{det} young nuclear physicist\\
`the young nuclear physicist'
\bg. *i piriniki neari fisikos\\
\textsc{det} nuclear young physicist \label{piriniki neari}\\
\cg. *i piriniki i fisikos\\
\textsc{det} nuclear \textsc{det} physicist \label{piriniki ds}\\
(`the nuclear physicist’)
\dg. *i fisikos i piriniki\\
\textsc{det} physicist \textsc{det} nuclear \label{piriniki post}\\
(`the nuclear physicist’)
\eg. *I fisikos ine piriniki.\\
\textsc{det} physicist \textsc{cop} nuclear\\
(`The physicist is nuclear.’) \label{piriniki pred}

An intriguing case is made by the modifier \is{relational modifier}\textit{theatriki} `theatrical’, for example modifying the noun \textit{parastasi} `play’. \textit{Theatriki} can (at least for some speakers including Jim Kotopoulis) occur in determiner spreading-structures, but only in prenominal position, never postnominally \ref{parastasi post}. Furthermore it can appear in predicative position \ref{parastasi pred}.

\ex. \ag. i theatriki (i) parastasi\\
\textsc{det} theatrical \textsc{det} play\\
`the theatrical play’
\bg. *i parastasi i theatriki\\
\textsc{det} play \textsc{det} theatrical \label{parastasi post}\\
\cg. I parastasi ine theatriki.\\
\textsc{det} play \textsc{cop} theatrical\\
`The play is theatrical.’ \label{parastasi pred}

Nevertheless, even when determiner spreading is invoked, there are still \is{adjective ordering restrictions}ordering restrictions observed when it co-occurs with a qualitative adjective such as \textit{orea} `nice’. Then, \textit{parastasi} has to occur closer to the noun. I show this below for a DP without \ref{orea no ds}, and a DP with determiner spreading \ref{orea ds}. \is{determiner spreading}

\ex. \ag. i orea (i) theatriki (i) parastasi\\
\textsc{det} nice \textsc{det} theatrical \textsc{det} play\\
`the nice theatrical play’
\bg. *i theatriki orea parastasi\\
\textsc{det} theatrical nice play \label{orea no ds}\\
\cg. *i theatriki i orea (i) parastasi\\
\textsc{det} theatrical \textsc{det} nice \textsc{det} play \label{orea ds}\\

These facts are highly challenging for existing theories. They suggest that the same, or at least a similar, inventory of \is{Units of Language}Units of Language is at stake for direct modification in Greek as in \is{Persian} Persian. Both languages, then, do not have in their inventory functional projections dedicated to categories of qualitative modifiers such as \textsc{dimension}. However, they do realize functional projections dedicated to specific direct modifiers, namely \is{relational modifier}relational modifiers. This explains the observed ordering principles. The DP structure for these two languages should thus look quite similar. \is{Modern Greek} \is{determiner spreading}

A remaining question is how cases in Greek can be explained where one, presumably direct, modifier occurs prenominally and another one occurs postnominally with \is{determiner spreading}determiner spreading.

\exg. i piriniki fisikos i neari\\
\textsc{det} nuclear physicist \textsc{det} young\\
`the young nuclear physicist’ \label{greek example for tree}

This receives a nice explanation under the structure presented here and the assumption of two DPs (as also argued in \citealt{laenzlinger:2011:elements}), a lower position of the \is{relational modifier}relational modifier and a combination of two movement operations. The underlying structure is presumably the one in \figref{DP structure Greek}. Assuming \textit{piriniki} is a direct modifier, it is merged in the direct domain, thus lower in the DP. The exact status of \textit{neari} cannot be determined, but it seems to be higher in any case.  \is{determiner spreading}

%\begin{figure}
%\centering
%\begin{tikzpicture}[baseline=0pt, remember picture]
%\tikzset{level distance=25pt, sibling distance=5pt}
% \Tree [.DP\textsubscript{external} {} [.D′\textsubscript{external} {} [.FP {} [.F′ F [.AgrP {} [.Agr′ Agr [.DP\textsubscript{internal} {} [.D′\textsubscript{internal} i [.AgrP {} [.Agr′ Agr [.FP [.AP \edge[roof]; neari ] [.F′ F [.AgrP {} [.Agr′ Agr [.FP\textsubscript{direct} [.AP \edge[roof]; piriniki ] [.F′ F [.NP \edge[roof]; fisikos; ] ] ] ] ] ] ] ] ] ] ] ] ] ] ] ] ]
% \end{tikzpicture}
%\caption{DP-Structure of Greek}
%\label{DP structure Greek}
%\end{figure}

%\todo{tree throws an error, was loaded as picture}

% \ex. \begin{tikzpicture}[baseline=0pt, remember picture]
% \tikzset{level distance=25pt, sibling distance=5pt}
%  \Tree [.DP\textsubscript{external} {}
%             [.D′\textsubscript{external} {}
%                 [.FP {}
%                     [.F′ F
%                         [.AgrP {}
%                             [.Agr′ Agr
%                                 [.DP\textsubscript{internal} {}
%                                     [.D′\textsubscript{internal} i [.AgrP \node (voitures2) {};
%                                         [.Agr′ Agr
%                                             [.FP
%                                                 [.AP \edge[roof]; neari ]
%                                                     [.F′ F
%                                                         [.AgrP \node (voitures3) {};
%                                                             [.Agr′ Agr [.FP\textsubscript{direct}
%                                                                 [.AP \edge[roof]; piriniki ]
%                                                                     [.F′ F
%                                                                         [.NP \edge[roof]; \node (voitures4) {fisikos}; ]
%                                                                     ]
%                                                                 ]
%                                                             ]
%                                                         ]
%                                                     ]
%                                                 ]
%                                             ]
%                                         ]
%                                     ]
%                                 ]
%                             ]
%                         ]
%                     ]
%                 ]
%             ]
%         ]
%  \end{tikzpicture}

% \begin{figure}
% \small
% \begin{forest}
% [DP\textsubscript{external}
%     [~]
%     [D′\textsubscript{external}
%         [~]
%         [FP
%             [~]
%             [F′
%                 [F]
%                 [AgrP
%                     [~]
%                     [Agr′
%                         [Agr]
%                         [DP\textsubscript{internal}
%                             [~]
%                             [D′\textsubscript{internal}
%                                 [i]
%                                 [AgrP
%                                     [~]
%                                     [Agr′
%                                         [Agr]
%                                         [FP
%                                             [AP
%                                                 [neari, roof]
%                                             ]
%                                             [F′
%                                                 [F]
%                                                 [AgrP
%                                                     [~]
%                                                     [Agr′
%                                                         [Agr]
%                                                         [FP\textsubscript{direct}
%                                                             [AP
%                                                                 [piriniki, roof]
%                                                             ]
%                                                             [F′
%                                                                 [F]
%                                                                 [NP
%                                                                     [fisikos, roof]
%                                                                 ]
%                                                             ]
%                                                         ]
%                                                     ]
%                                                 ]
%                                             ]
%                                         ]
%                                     ]
%                                 ]
%                             ]
%                         ]
%                     ]
%                 ]
%             ]
%         ]
%     ]
% ]
% \end{forest}
% \caption{DP-Structure of Greek}
% \label{DP structure Greek}
% \end{figure}

\begin{figure}[H]
\centering
%\vspace{-3cm}
\includegraphics[scale=1.2, width=12cm]{DP structure Greek.png} 
 \caption{DP-structure of Greek (exemplified with \ref{greek example for tree})}
 \label{DP structure Greek}
\end{figure}

The derivation likely proceeds as follows: First the NP moves in a \is{pied-piping}pictures-of-whom type movement operation to the specifier of a higher functional head (AgrP in \citealt{cinque:2017:micro, cinque:2023:linearization}), dragging along the relational modifier \textit{piriniki}, but crucially not affecting the order. Then this chunk moves across the qualitative modifier \textit{neari} to the \is{left periphery}left periphery across the determiner. Crucially, this modifier remains in situ and we get the order \textit{piriniki} $\succ$ \textit{fisikos} $\succ$ \textit{i} $\succ$ \textit{neari}. Next, the determiner moves to the DP-external domain and leaves a copy in the internal DP, as, for instance, argued in \citet{laenzlinger:2011:elements} and see \citet{grohmann:2003:prolific, grohmann:2015:demonstrative}.

\subsubsection{Summary}

The preceding discussion has shown that the DP structure presented in this study is applicable to other languages, emphasizing its universal status, but that parametric variation is permitted. Importantly, this variation is not random. There are languages which do not show ordering restrictions of categories such as \textsc{quality}, but do show ordering restrictions when it comes to direct-only modifiers. On the other hand, the reverse case is (at least so far) missing from the paradigm. Languages such as \is{English} English or \is{German} German show ordering restrictions in all scenarios.

On the basis of these facts, I propose the following \is{relational modifier}generalization. \is{Generalization on Cross-Linguistic Ordering Restrictions}

\ex. \emph{Generalization on Cross-Linguistic Ordering Restrictions}\\
If a language shows ordering restrictions among qualitative modifiers such as \textsc{quality} and \textsc{dimension}, it will also show ordering restrictions among direct-only modifiers, for example relational and non-intersective modifiers, among each other and with other categories. In other words, if a language has the language-specific inventory to realize functional projections for categories of \textsc{quality}, \textsc{dimension}, etc., it will also realize functional projections for direct-only modifiers.

This means that there are indeed languages between the English type, in which all categories adhere to specific, preferred orders in context-neutral utterances, and the Japanese type, in which there is the same universal spine observable, but not the language-specific entities for functional projections of fine-grained semantic categories. It is for future research to prove or refute this claim and to see whether more points can be added to the continuum and whether an even more fine-grained generalization can be put forth.

\subsubsection{Two more languages} \label{more languages}

Before closing this subsection, I would like to draw attention to two languages that prove the above made claim in different ways. First, \citet{brown:2024:adj} shows that in Naasioi, a Papuan language, direct adjectives do not adhere to \is{absence of ordering restrictions}ordering restrictions. Adjectives in this language can appear pre- or postnominally, seemingly without restrictions. Furthermore, direct can be differentiated from indirect adjectives via presence of an additional article in the latter case, making this construction appear very similar to determiner spreading in \is{Modern Greek} Modern Greek. First, direct modifiers must appear closer to the noun than indirect modifiers, illustrated in N-A-A sequences below. \is{Naasioi}

\ex. \ag. Te mosi pangkaing te mutaanu aatu-maang. (Naasioi)\\
\textsc{det} dog big \textsc{det} black sleep-\textsc{prs.prog}\\
`The big dog which is black is sleeping.’
\bg. *Te mosi te mutaanu pangkaing aatu-maang.\\
\textsc{det} dog \textsc{det} black big sleep-\textsc{prs.prog}\\ \null \hfill \citep[432]{brown:2024:adj}

Assuming all modifiers that appear without the additional article \textit{te} are direct, they should adhere to \is{absence of ordering restrictions}ordering restrictions in contexts with two or more direct modifiers. Alas, they do not, as \citet{brown:2024:adj} argues based on fieldwork data. This generalization seems to hold, at least, for modifiers of \textsc{dimension}, \textsc{color} and \textsc{quality} and is exemplified below with a string of three adjectives, which, according to \citet{brown:2024:adj} are fine in either order.

\exg. Te orara pangkaing arasi kiring dua-’ar-ing.\\
\textsc{det} bad big brown breadfruit fall-\textsc{pl-imm.pst}\\
`The bad big brown breadfruits fell.’ \\ \null \hfill \citep[429]{brown:2024:adj}

In other words, this language poses a problem for standard Cartography as well. In fact, since it seems to pattern with the languages depicted above (although no immediate conclusion can be drawn about relational and non-intersective modifiers), the same facts seem to be at stake, namely a direct domain with unspecific projections, in which various movement operations can take place.

Containment relations akin to the one presented for \is{Persian} Persian and \is{Modern Greek} Greek are given for \is{Ojibwe} Ojibwe, an Algonquian language, in \citet{rosen:2016:on}. This language seems to have a small class of direct modifiers, which are prenominally bound to the noun. The interesting part is that pairs of direct adjectives of \textsc{color} and \textsc{nationality} as well as \textsc{color} and \textsc{dimension} do not adhere to \is{adjective ordering restrictions}ordering restrictions, but pairs of \textsc{dimension} and \textsc{nationality} must do so (\textsc{ni} = inanimate noun).

\newpage
\ex. \textsc{color} vs. \textsc{nationality}: No ordering restrictions
\ag. misko-Anishinaabe-mashkimod\\
red-Ojibwe-bag.\textsc{ni}\\
\bg. Anishinaabe-misko-mashkimod\\
Ojibwe-red-bag.\textsc{ni}\\
`a red Ojibwe bag’ \null \hfill \citep[161]{rosen:2016:on}

\ex. \textsc{dimension} vs. \textsc{color}: No ordering restrictions
\ag. gichi-misko-waakaa’igan\\
big-red-house.\textsc{ni}\\
\bg. misko-gichi-waakaa’igan\\
red-big-house.\textsc{ni}\\
`a big red house’ \null \hfill \citep[161]{rosen:2016:on}

\ex. \textsc{dimension} vs. \textsc{nationality}: Ordering restrictions
\ag. gichi-Anishinaabe-mashkimod\\
big-Ojibwe-bag.\textsc{ni}\\
`a big Ojibwe bag’
\bg. *Anishinaabe-gichi-mashkimod\\
Ojibwe-big-bag.\textsc{ni} \\ \null \hfill \citep[162]{rosen:2016:on}

Furthermore, for some reason, \textsc{nationality} adjectives can \is{absence of ordering restrictions}appear to the left of \textsc{dimension} adjectives if an adjective of \textsc{color} intervenes \citep[162]{rosen:2016:on}.

\exg. Anishinaabe-misko-gichi-adoopowin\\
Ojibwe-red-big-table.\textsc{ni}\\
`a big red Ojibwe table’ \null \hfill \citep[162]{rosen:2016:on}

This order string \textsc{nationality} $\succ$ \textsc{color} $\succ$ \textsc{dimension} is the mirror-image order of modifiers compared to their assumed merge position. Naturally, it is unforeseen by any \is{adjective ordering restrictions}hierarchy. It should be impossible with modifiers in prenominal position according to Greenberg’s U20 and cannot be explained if we adopt the assumption that only the lexical head, the noun, can move. In other words, this surface order can only be explained if either, individual movement of modifiers, or, unspecific functional projections and free Merge positions are assumed. However, the latter assumption would have the big disadvantage that the strict order \textsc{dimension} $\succ$ \textsc{nationality} is unaccounted for. \is{U20}

Needless to say more work is necessary here, but one, very tentative, conclusion could be that, all things being equal, this language also does not realize functional heads based on specific semantic categories explaining the, in most cases, free order. However, there is a ban on the merge order of a modifier of \textsc{dimension} followed by a modifier of \textsc{nationality}  explaining why a string of two adjacent modifiers of those categories in that order is never available. I thus interpret these facts similarly to the Japanese case, but with the twist that modifiers of \textsc{dimension} might never be merged directly after modifiers of \textsc{nationality}.

\subsection{Remaining questions}

This section discusses two remaining questions, namely whether there is complete freedom of ordering in Japanese, and whether complex direct modifiers actually exist.

\subsubsection{\#1: Is there complete optionality in ordering?} \label{optionality}

If Japanese is a language, potentially one of many, with a broad direct domain lacking specific functional projections to dictate ordering, does that mean that no ordering principles exist at all? In other words: Is \is{absence of ordering restrictions}ordering of modifiers completely optional in this language? It was established that the semantic category of a modifier cannot be a factor, but in the following, I discuss potential other factors.


\subsubsubsection{Word class, morphology}



I took care in balancing out the word classes that I tested for ordering preferences. Not only do the two Japanese adjective groups not adhere to ordering restrictions within and across one another \citep{kohlich:2018:die}, the present study showed that direct \textit{no}-modifiers, not only among one another, but also together with \textit{i}- and \textit{na}-adjectives, also do not lead to any ordering preferences in speakers. This clearly suggests the absence of the factors \textit{word class} or \textit{morphology}, which readers can verify in the examples above. \is{word classes}


\subsubsubsection{Phonology}
\largerpage[-1]
It has been argued in \citet{nagano:2016:are} that phonology might play a role for \is{absence of ordering restrictions}ordering of modifiers, which is exactly why I have chosen to test modifiers with different numbers of morae. As it turns out, phonology does seem to be a factor, but only to a limited extent. For example, in none of the pairs \textit{Denmāku-no} `Denmark, Danish’ vs. \textit{chīku-no} `teak’ (5 vs. 3 morae) or \textit{ao-i} `blue’ vs. \textit{Chūgoku-no} `China, Chinese’ (2 vs. 4 morae) did ordering restrictions emerge. Note that I counted \is{mora length}mora without the monomoraic elements following.

The only strict ordering seems to be dictated by the single monomoraic modifier of \textsc{material} \textit{ki-no} `wood(en)’. Recall the example given in \citet{watanabe:2012:direct}, who argues that the combination with a trimoraic modifier leads to ungrammaticality if it is not placed closest to the noun. The relevant example is repeated from \ref{hokuo} below.

\ex. \ag. hokuō-no ki-no isu \\
North.Europe-\textsc{mod} wood-\textsc{mod} chair\\
`North European wood(en) chair'
\bg. *ki-no hokuō-no isu\\
wood-\textsc{mod} North.Europe-\textsc{mod} chair \\ \null \hfill \citep[508]{watanabe:2012:direct}

Now, this extends to the combination of \textit{ki} `wood(en)’ with other modifiers. The speakers I consulted prefer \textit{ki-no} closer to the head noun, but do not find the opposite example to be completely ungrammatical, just very marked. This is shown below. \is{mora length}

\ex. \ag. Denmāku-no ki-no isu\\
Denmark-\textsc{mod} wood-\textsc{mod} chair\\
\bg. \#ki-no Denmāku-no isu\\
wood-\textsc{mod} Denmark-\textsc{mod} chair\\
`Danish wood chair'

Note that it really is phonology that is the factor at stake here, not semantic category, as the equivalent \textit{mokusei-no} `wood(en)’ is possible in either position and not marked.

\ex. \ag. Denmāku-no mokusei-no isu\\
Denmark-\textsc{mod} wood-\textsc{mod} chair\\
\bg. mokusei-no Denmāku-no isu\\
wood-\textsc{mod} Denmark-\textsc{mod} chair\\
`Danish wood chair'

\largerpage[-1]
While this means that phonology does have an effect on ordering to a certain extent, it is questionable how big its effect actually is. As noted, \textit{ki-no} is the only monomoraic modifier in the data set. This means that we cannot test similar modifiers to establish strict paradigms. On the other hand, with bimoraic modifiers such as \textit{ao-i} `blue’, I did not detect such an effect.\footnote{A comment made by an anonymous reviewer, alongside a provided example, fits well with this observation. According to their judgment, all things being equal, the following minimal pair with a bimoraic modifier of \textsc{nationality} also does not offer a contrast.

\par

\ex. \ag. Tai-no mokusei-no isu\\
Thailand-\textsc{mod} wood-\textsc{mod} chair\\
\bg. mokusei-no Tai-no isu\\
wood-\textsc{mod} Thailand-\textsc{mod} chair\\
`Thai wood chair'

\par}


\subsubsubsection{Pragmatic factors, scope}


As already mentioned in Section \ref{japanesedp}, the only cases in which fixed order emerge involve different scope readings. In \ref{example scope}, the speaker is restricting the \is{scope}domain of German vases, not the domain of vases \textit{per se} to a concrete \textit{German vase} of which only one is available in the discourse. The reverse order, given the same context, is significantly degraded.

\ex. \label{example scope} \ag. yui'itsu-no Doitsu-no kabin\\
only-\textsc{mod} German-\textsc{mod} vase\\
`the only German vase'
\bg. ??Doitsu-no yui'itsu-no kabin\\
German-\textsc{mod} only-\textsc{mod} vase\\

This is thus a case in which the two modifiers occupy distinct merge positions (at least in relation to one another, the precise projection being unimportant), although both are inside the direct domain, as illustrated in \figref{structure yuiitsu-no doitsu-no}.

\begin{figure}
\centering
\begin{tikzpicture}[baseline=0pt, remember picture]
\tikzset{level distance=25pt, sibling distance=5pt}
\Tree [.DP\textsubscript{internal} {} [.D′\textsubscript{internal} [.FP\textsubscript{indirect} {} [.F′\textsubscript{indirect} [.dP [.FP\textsubscript{direct} [.ModP {} [.Mod′ [.XP \edge[roof]; yui'itsu ] [.Mod no ] ]] [.F′\textsubscript{direct} [.FP\textsubscript{direct} [.ModP {} [.Mod′ [.XP \edge[roof]; Doitsu ] [.Mod no ] ]] [.F′\textsubscript{direct} [.NP \edge[roof]; kabin ] F ] ] F ] ] d ] F ] ] D ] ]
\end{tikzpicture}
\caption{Structure of \textit{yui’its-no Doitsu-no kabin} `only German vase’}
\label{structure yuiitsu-no doitsu-no}
\end{figure}

Connected to this is the strong nominal character of modifiers of \textsc{nationality}. Some, but not all, speakers reported that the second order to them suggests a reading like `the only vase IN Germany’. This also becomes obvious when combining such modifiers with possessive modifiers as in the following. Here, speakers again find a specific order more natural when they want to identify the particular referent: That in which the \textsc{nationality} modifier is closer to the noun.

\ex. \ag. watashi-no Doitsu-no kabin\\
I-\textsc{mod} German-\textsc{mod} vase\\
`my German vase(-s)’
\bg. ??Doitsu-no watashi-no kabin\\
German-\textsc{mod} I-\textsc{mod} vase\\

Seemingly, scope is at play again. The possessive modifier \textit{watashi-no} `my’ takes scope over \textit{Doitsu-no kabin}, not just \textit{kabin} alone, restricting the domain to a specific vase, namely one belonging to the speaker. Again, the reverse order does not involve scope and means something like `my vase(-s), which happen(s) to be in German / form Germany’.

These facts highlight that ordering in Japanese is not completely unrestricted, but that there are influencing factors, among them, to a minor extent, phonology, and, to a larger extent, scope restrictions. Importantly, these factors mostly seem to be \is{scope}\textit{syntax-external}, as also reported for similar cases in the clause in \citet{ramchand:2014:deriving}. Another such factor, presented in Section \ref{complexity} for indirect modifiers, is \is{complexity}\textit{complexity}. This factor might not be really testable for direct modifiers, one of whose defining characteristics is to be, generally, less complex than their indirect counterparts. However, a related phenomenon will be discussed in the next section.

\subsubsection{\#2: Complex direct modifiers and the situation in Classical Japanese} \label{complexdirect}
\largerpage

Another remaining issue to address concerns the possible existence of \is{complex direct modifier}complex, that is clausal, direct modifiers. It is one of the characteristics of direct modifiers that they exhibit a less complex structure than indirect modifiers, which are arguably always embedded in a CP (\citealt{cinque:2010:adj, cinque:2020:relcl}, see Section \ref{indirect mod}). However, there is reason to believe that direct modifiers can also be complex. First, direct modifiers are not only functional, but can also be \textit{phrasal}. Evidence for this comes from the fact that direct adjectives can, in some languages, appear with overt adjectival arguments or adjuncts. This is shown below with an example from Bulgarian involving the \is{non-intersective modifier}direct (non-intersective) adjective `main’ and a prepositional complement (cf. \citealt[§ 4.1.2]{cinque:2010:adj} for more data). \is{Bulgarian}

\exg. glavnata po zna\v{c}enie pri\v{c}ina (Bulgarian)\\
main.the in significance reason\\
`the main reason for importance’ \null \hfill \citep[47]{cinque:2010:adj}

In Japanese, some direct modifiers are clearly \textit{clausal}. This concerns especially structures in which such modifiers co-occur with an overt subject argument. Some relevant data points are given in \citet{hoshi:2002:kaynean}, who shows that some of the alleged examples of non-intersective modifiers given in \citet{yamakido:2000:japanese} not only exhibit a predicative character, but also become intersective, highlighting this predicative character, when a subject appears (see \citealt{yamakido:2005:nature} for discussion). Observe \ref{hoshi example}. \is{complex direct modifier}

\exg. Pītā-wa [tsukiai-ga furu-i] tomodachi-da.\\
Peter-\textsc{top} association-\textsc{nom} long.time-\textsc{mod} friend-\textsc{cop}\\
`Peter has been a friend for a long time.’ (Lit.: `Peter is a friend whose association is old.’)  \hfill \citep[11]{hoshi:2002:kaynean} \label{hoshi example}

Here, the alleged direct modifier \textit{furu-i} `long.time’, which in combination with human entities only yields a non-intersective reading (Section \ref{non-intersective Japanese}), appears in a clausal structure with an overt subject. Furthermore, it now has acquired an intersective interpretation \textit{and} is possible in predicative position, evidenced in \ref{hoshi intersective}.

\exg. Pītā-wa [tsukiai-ga furu-i], soshite Pītā-wa tomodachi-da.\\
Peter-\textsc{top} association-\textsc{nom} long.time-\textsc{cop} and Peter-\textsc{top} friend-\textsc{cop}\\
`The association with Peter has been long and Peter is a friend.’ \\ \null \hfill \citep[11]{hoshi:2002:kaynean} \label{hoshi intersective}

Now admittedly, as has been argued, \textit{furu-i} can appear in predicative position as long as it is predicated on inanimate entities, which means it is not a direct-only modifier in the first place. However, a clausal construction is also possible for \is{relational modifier}relational modifiers, as shown below for \textit{Chūgoku-no} `China, Chinese’. \is{complex direct modifier}

\exg. seizōkoku-ga Chūgoku-no kuruma\\
country.of.production-\textsc{nom} China-\textsc{mod} car\\
`a car whose country of production is China’

In this case, the modifier also allows paraphrasing \textsc{no} with \textsc{de-ar-u} and inflection for past.

\ex. \ag. seizōkoku-ga Chūgoku-de-ar-u kuruma\\
country.of.production-\textsc{nom} China-\textsc{pred-cop-prs} car\\
`a car whose country of production is China’
\bg. seizōkoku-ga Chūgoku-d-at-ta kuruma\\
country.of.production-\textsc{nom} China-\textsc{pred-cop-pst} car\\
`a car whose country of production was China’

Two things are of importance here: Again, note that the modifiers in the preceding examples, unlike the examples in \citet{cinque:2010:adj}, are not only phrasal but clausal. Second, what \textit{Chūgoku} here actually refers to is the country \textit{China}, but not the notion \textit{from China} as in examples such as \textit{Chūgoku-no kabin} `Chinese vase’ presented in Section \ref{RAJapanese}. \textit{Chūgoku-no} is \textit{per se} not a direct-only modifier either. \is{complex direct modifier}

More tricky is the following example with the \is{rentaishi@\textit{rentaishi}}\textit{rentaishi} \textit{ōki-na} `big’, which, different to \textit{ōki-i}, cannot appear in predicative position. However, it can take a subject argument in attributive position.

\exg. yane-ga ōki-na ie\\
roof-\textsc{nom} big-\textsc{mod} house\\
`house whose roof is big’ \label{yane}

This means, the structural make-up, at least on the surface, for such constructions is similar to the following example, repeated from Section \ref{comparisonpreviousno}.

\exg. kami-ga kirei-na hito\\
hair-\textsc{nom} pretty-\textsc{mod} person\\
`a person whose hair is pretty' \null \hfill \citep[34]{kaplan:1995:the}

\newpage
The question is how to analyze such cases. An interesting option is the notion of direct clausal modifiers, i.e. relative clauses, which, to my knowledge, was only considered a possibility by \citet{kim:2019:syntax}. She offers the following example:

\ex. Politicians [who make extravagant promises] aren’t trusted. \\ \null \hfill \citep[139]{kim:2019:syntax}

Since this relative clause, from a semantic point of view, fulfills the function of ``characterizing people who are known as politicians'' \citep[139]{kim:2019:syntax}, it has an individual-level, not a \is{stage/individual-level}stage-level, reading. This reading, in turn, is associated with direct modification \citep{cinque:2010:adj}, the same as ambiguous adjectives such as \textit{visible}. \citet{kim:2019:syntax} leaves open what the precise structure of such relative clauses and their integration into the DP is, and indeed it is unclear what kind of structure we should assign to DPs such as \ref{yane}.

I assume, parallel to what is assumed for indirect modifiers with overt subjects, that the modifier \textit{ōki-na} `big’ projects a CP, making available a subject position in Spec, IP and providing a landing site for the relative clause \is{operator}operator in Spec, CP ensuring that the relative clause head is also correctly interpreted clause-internally. The whole CP is licensed by \textsc{mod}, which is realized as no overt tense appears (Section \ref{indirect mod}). This is given below.

\exg. [\textsubscript{ModP} [\textsubscript{CP} \textit{Op}\textsubscript{1} [\textsubscript{IP} yane-ga [\textsubscript{PredP} [\textsubscript{AP} ōki-]]]]na] ie\textsubscript{1}\\
{} {} Op {} roof-\textsc{nom} {} {} big-\textsc{mod} house\\
`house whose roof is big’

Taking into account the existence of complex, that is relative clause-like, direct modifiers, is intriguing, but would violate the definition of \textit{direct-only} adopted in this study forcing certain concessions. On the other hand, I do not see an immediately tenable alternative which would still maintain the proposal. For now, I leave this for future research.

Importantly, clausal modification is impossible for many of the prime examples of direct-only modifiers, for instance \textit{shōrai-no} `future’, which do not allow an overt subject argument DP-internally.

\exg. *Kare-ga seikō-ga shōrai-no purosakkāsenshu desu.\\
he-\textsc{nom} success-\textsc{nom} future-\textsc{mod} pro.soccer.player \textsc{cop.pol}\\
(`He is a professional soccer player whose success is (expected) in the future.’)

Furthermore, \ref{yane} could find another explanation which relates to the nature of \textit{ōki-na} as a \is{rentaishi@\textit{rentaishi}}\textit{rentaishi}. Recall that all verbally derived \textit{rentaishi} are fossilized in their attributive form (\textit{rentaikei}), and this also extends to adjectival \textit{rentaishi} co-occurring with the premodern copula \textit{nari}/\textit{tari}, inflected for \textit{naru}/\textit{taru} respectively, and potentially \textit{na}. As also argued in Section \ref{kaplan}, the attributive form of verbs in Classical Japanese clearly projected a C layer. Evidence for this comes from the fact that verbs obligatorily inflected for the attributive form in embedded wh-questions, where it can be analyzed as a complementizer (see \citealt{kinsui:1995:iwayuru, kaplan:1995:the, hiraiwa:2005:dimensions}). This suggests that such modifiers might indeed still project a CP-layer explaining the clausal behavior.\footnote{In this respect, note also the following example with a verbally derived \textit{rentaishi} and a nominal subject, which prompted discussions already in the middle of the 20th century. Several authors have argued that the transition from verb to \textit{rentaishi} had not been fully achieved, accounting for the possibility of overt subjects. See \citet{sekine:1937:fukutaishi}, \citet{matsushita:1976:kaisen} and \citet{kieda:1938:koto}.

\par

\exg. kimi-ga iwayuru shingi\\
you-\textsc{nom} so.called loyalty\\
`your so-called loyalty' (Lit. `what you call loyalty’) \null \hfill \citep[1:28]{sekine:1937:fukutaishi}

\par

} \is{complementizer}

This discussion also extends to other \is{rentaishi@\textit{rentaishi}}\textit{rentaishi} in Modern Japanese such as \textit{aru} `a certain’ in example \ref{aru}. There are (at least) two analyses possible for such lexemes: According to the first, they are analyzed as \textit{bona fide} (but fossilized) direct modifiers which do not project a higher functional layer, and which are somehow not licensed via \is{MOD@\textsc{mod}}\textsc{mod}. Under this analysis, the ending \textit{ru} is taken to be part of the stem. Alternatively, staying closer to the structure such modifiers presumably had in Classical Japanese, they are direct modifiers, but project higher IP and CP layers. Under this analysis, the ending \textit{ru} is a complementizer head (and I abstract away from the possible presence or absence of \textsc{mod}). Both possible analyses are depicted in \figref{analyses for aru gakusei}. \is{complex direct modifier}

\begin{figure}

\begin{subfigure}{\textwidth}
\centering
\begin{tikzpicture}[baseline=0pt, remember picture]
\tikzset{level distance=25pt, sibling distance=5pt}
 \Tree [.DP\textsubscript{internal} {} [.D′\textsubscript{internal} [.FP\textsubscript{direct} [.XP \edge[roof]; aru ] [.F′\textsubscript{direct} [.NP \edge[roof]; gakusei ] F ]] D ] ]
\end{tikzpicture}
\caption{\textit{aru gakusei}: Analysis 1}
\label{aru gakusei analysis 1}
\end{subfigure}

\begin{subfigure}{\textwidth}
\centering
\begin{tikzpicture}[baseline=0pt, remember picture]
\tikzset{level distance=25pt, sibling distance=5pt}
 \Tree [.DP\textsubscript{internal} {} [.D′\textsubscript{internal} [.FP\textsubscript{direct} [.CP {} [.C′ [.IP {} [.I′ [.VP \edge[roof]; a ] I ] ] [.C ru ] ] ] [.F′\textsubscript{direct} [.NP \edge[roof]; gakusei ] F ]] D ] ]
\end{tikzpicture}
\caption{\textit{aru gakusei}: Analysis 2}
\label{aru gakusei analysis 2}
\end{subfigure}

\caption{Two possible analyses for \textit{aru gakusei} `a mere student’ (example \ref{aru}}
\label{analyses for aru gakusei}

\end{figure}

\largerpage[2]
This discussion also concerns \is{Classical Japanese}Classical Japanese, where all verbally and adjectivally derived lexemes analyzed as \is{rentaishi@\textit{rentaishi}}\textit{rentaishi} in Modern Japanese were still inflecting parts of speech in Classical Japanese. This means, in attributive position, they arguably could perform a role as direct modifiers casting doubt on the analysis of the attributive form (\textit{rentaikei}) in Classical Japanese -- although probably projecting a C layer (\citealt{kinsui:1995:iwayuru, kaplan:1995:the, hiraiwa:2005:dimensions}; cf. Section \ref{kaplan}) -- as being necessarily tied to indirect modification. Under this account, all modifiers had access to both kinds of modification at this stage of the language as well. For example, \textit{i}-adjectives, whose attributive form was represented via the suffix -\textsc{ki}, were presumably structurally ambiguous between a direct and an indirect status, similar to the situation in Modern Japanese. \is{complex direct modifier}

\exg. taka-ki mono (Classical Japanese)\\
high-\textsc{rtk} thing\\
`a high thing' \null \hfill

To summarize, the existence of clausal direct modifiers provides a challenge to the theory presented here, but could be explained via assuming either the existence of complex direct modifiers, or giving a special status to certain lexemes that are only direct in certain usage domains. A bigger challenge is posed by direct modifiers historically derived from verbs, and this extends to the situation in Classical Japanese. This is a fruitful area for future research.

\subsection{Summary}

In this section, I argued that Japanese lacks the language-specific categories to realize specific functional projections. This explains the lack of strict ordering and shows that direct modification and ordering are orthogonal issues. I then discussed the cross-linguistic implications of this claim and put forth a generalization on ordering principles. Finally, I showed that ordering in Japanese is in fact not entirely ungoverned either and discussed the possibility of complex direct modifiers.

\section{Chapter summary and closing remarks}

This chapter started with a discussion of why direct modification has so far been neglected, or entirely dismissed, in Japanese. Then, criteria for direct modification were introduced and it was shown that the lack of predicativity is the leading criterion, which is often accompanied by other criteria such as non-gradability. The following observations emerged in the course of the discussion:

\ex. \a. Japanese exhibits a variety of different direct-only modifiers: The equivalents of cross-linguistically well-known categories of non-intersective and relational modifiers, as well as cases of direct qualitative modifiers, idiomatic modifiers, direct quantificational modifiers and the language-specific group of \textit{rentaishi}.
\b. Almost all cases of direct-only modifiers are found among \textit{no}-modifiers.
\c. There is a lack of ordering restrictions which is due to a lack of projections dedicated to specific categories inside the direct domain.
\d. The reason for this is that Japanese does not have the language-specific entities to realize such dedicated functional projections.

I believe that the challenges these findings pose for Cartography are also chances to expand this research field and adapt it in a way so that the very elaborate predictions it makes become even more powerful. Postulating a broader, more fluid set of functional projections inside the direct domain preempts the common criticism regarding superfluous projections as well as alleged violations of ordering restrictions. The concept of cross-linguistic variation enriches our understanding of the DP and generates new insights for other languages.

It is definitely a desideratum for future research to establish more cases of clear-cut ordering restrictions in Japanese, in addition to scopal relations or restrictions among direct and indirect modifiers. This would not only help expanding the hierarchy of direct modifiers proposed in this book in a language-specific and cross-linguistic perspective, but also could expand the structure of the direct domain I have proposed here and add a more fine-grained dimension to it.
